% Traduction : comparatifs et superlatifs ; redacción : el currículum vítae — Espagnol Tle A — Cours signé OnBuch+
% Programme officiel MINESEC. Compiler avec : tectonic E17-comparativos-superlativos-curriculum.tex
\newcommand{\DOCMATIERE}{Espagnol}
\newcommand{\DOCNIVEAU}{Tle A}
\newcommand{\DOCMODULE}{Módulo 4 : Activités économiques et monde du travail}
\newcommand{\DOCLECON}{E17}
\newcommand{\DOCTITRE}{Traduction : comparatifs et superlatifs ; redacción : el currículum vítae}
% OnBuch+ — préambule « cours signés » ENRICHI (compilé avec tectonic / XeLaTeX)
% Système visuel « Pop » d'OnBuch+ : fond crème (jamais blanc), bordures encre
% épaisses, ombres « dures » décalées, titres Archivo Black en capitales.
% Version MATHÉMATIQUES : graphes (pgfplots/tikz), boîtes pédagogiques
% (définition, propriété, méthode, attention, le savais-tu, exercice résolu,
% exercice, TP, bilan), page de garde et signature OnBuch+ ancrée en bas.
%
% Macros attendues dans main.tex AVANT \input{preamble} :
%   \DOCMATIERE  \DOCNIVEAU  \DOCMODULE  \DOCLECON (numéro)  \DOCTITRE
\documentclass[11pt]{article}

\usepackage{fontspec}
\usepackage[spanish,french]{babel} % français = langue principale ; l'espagnol via \esp{..} / espagnol
\usepackage[a4paper,margin=2cm,top=3.4cm,bottom=2.8cm,headheight=1.4cm,footskip=1.3cm]{geometry}
\usepackage{amsmath,amssymb,amsfonts}
\usepackage{mathtools} % \xmapsto, \coloneqq... souvent utilisés par les modèles
\usepackage{cancel} % simplifications barrées (bilans de forces)
% \abs{z} (module, valeur absolue) et \norm{u} (norme), utilisés par les modèles
\providecommand{\abs}{}\let\abs\relax\DeclarePairedDelimiter{\abs}{\lvert}{\rvert}
\providecommand{\norm}{}\let\norm\relax\DeclarePairedDelimiter{\norm}{\lVert}{\rVert}
\usepackage[table]{xcolor} % [table] : fournit aussi \rowcolors (alternance de lignes)
\usepackage{pagecolor}
\usepackage{tikz}
\usetikzlibrary{arrows.meta,positioning,calc,shapes.geometric,shapes.misc,shapes.arrows,decorations.pathmorphing,decorations.pathreplacing,decorations.markings,patterns,shadows,mindmap,trees,fit,backgrounds,matrix,intersections,angles,quotes}
\usepackage{pgfplots}
\usepackage[version=4]{mhchem} % équations nucléaires \ce{^{235}_{92}U}
\usepackage[european,siunitx]{circuitikz} % schémas électriques
\pgfplotsset{compat=1.18}
\usepgfplotslibrary{fillbetween,groupplots} % aires entre courbes (calcul intégral)
\usepackage{enumitem}
\usepackage{fancyhdr}
% [most] charge toutes les bibliothèques tcolorbox (listings, minted, raster,
% documentation...) et gonfle inutilement la mémoire TeX sur un document long
% (~300 boîtes) : on ne charge que ce qui est réellement utilisé.
\usepackage[skins,breakable]{tcolorbox}
\usepackage{graphicx}
\usepackage{colortbl}
\usepackage{booktabs}
\usepackage{tabularx}
\usepackage{diagbox} % case d'en-tête diagonale des tableaux à double entrée
% Colonnes extensibles centrée / gauche / droite (C, L, R), souvent utilisées par les modèles
\newcolumntype{C}{>{\centering\arraybackslash}X}
\newcolumntype{L}{>{\raggedright\arraybackslash}X}
\newcolumntype{R}{>{\raggedleft\arraybackslash}X}
\usepackage{array}
\usepackage{multicol}
\usepackage{siunitx}
% parse-numbers=false : accepte aussi \SI{2,5 \times 10^{-3}}{...}
\sisetup{locale=FR,output-decimal-marker={,},group-minimum-digits=4,parse-numbers=false}
\DeclareSIUnit\ohm{\ensuremath{\symup{\Omega}}} % U+2126 absent de la police
\DeclareSIUnit\division{div} % réglages d'oscilloscope (ms/div, V/div)
\DeclareSIUnit\tr{tr} % tours (vitesse de rotation, ex. \SI{1500}{\tr\per\minute})
\DeclareSIUnit\molar{M} % concentration molaire (ex. \SI{3}{\molar}, \SI{1}{\micro\molar})
\DeclareSIUnit\an{an} % année (le modèle confond parfois avec \year, un compteur TeX) — ex. \SI{5}{\kilo\meter\per\an}
\DeclareSIUnit\FCFA{FCFA} % franc CFA (ex. \SI{500}{\FCFA\per\kilo\gram})
\DeclareSIUnit\degreeBrix{\textdegree Brix} % teneur en sucre d'un sirop/jus (réfractométrie)
\DeclareSIUnit\month{mois} % le modèle utilise parfois \month, un compteur TeX, comme unité de durée
\usepackage{qrcode}
% \degree : commande fréquemment utilisée par les modèles pour un angle ou une
% température en dehors de \SI{}{} (non fournie par les paquets chargés).
\providecommand{\degree}{^{\circ}}
% \qed (fin de démonstration), utilisé par les modèles sans amsthm
\providecommand{\qed}{\hfill$\square$}
% \barre : double barre des valeurs interdites dans un tableau de variations (tkz-tab)
\providecommand{\barre}{\ensuremath{\|}}
% environnement proof (amsthm non chargé) utilisé par les modèles
\ifdefined\proof\else\newenvironment{proof}[1][Démonstration]{\par\noindent\textit{#1.}\ }{\qed\par}\fi
\usepackage{lastpage}
\usepackage{hyperref}
\hypersetup{hidelinks}

% ============================================================
% Palette « Pop » OnBuch+ — mêmes hex que la classe P (plus_ui.dart)
% ============================================================
\definecolor{poporange}{HTML}{F59321}
\definecolor{popdark}{HTML}{E07A0C}
\definecolor{popcream}{HTML}{FFF6E8}
\definecolor{popink}{HTML}{1C1714}
\definecolor{poppurple}{HTML}{7A5AE0}
\definecolor{popgreen}{HTML}{2E9E6A}
\definecolor{poppink}{HTML}{E0457B}
\definecolor{popblue}{HTML}{2F7DE1}
\definecolor{popgold}{HTML}{C98A12}
\definecolor{popmuted}{HTML}{8A7F76}
\definecolor{popline}{HTML}{EADFD2}
\definecolor{popgray}{HTML}{8A7F76}
\colorlet{popgrey}{popgray}
% Alias tolérants (le modèle nomme parfois les couleurs différemment)
\colorlet{popwhite}{white}
\colorlet{popblack}{popink}
\colorlet{popred}{poppink}
\colorlet{popyellow}{popgold}
% Teintes claires (fonds de tableaux/encadrés) — le modèle nomme parfois
% les couleurs avec un suffixe L (ex. popblueL) sans les avoir définies.
\colorlet{poporangeL}{poporange!20}
\colorlet{popdarkL}{popdark!20}
\colorlet{poppurpleL}{poppurple!20}
\colorlet{popgreenL}{popgreen!20}
\colorlet{poppinkL}{poppink!20}
\colorlet{popblueL}{popblue!20}
\colorlet{popgoldL}{popgold!20}
\colorlet{popredL}{popred!20}
\colorlet{popgrayL}{popgray!20}
% Teintes claires pour fonds de tableaux / zones
\colorlet{popblueL}{popblue!10!white}
\colorlet{poporangeL}{poporange!14!white}
\colorlet{popgreenL}{popgreen!12!white}
\colorlet{poppurpleL}{poppurple!12!white}
\colorlet{poppinkL}{poppink!12!white}
\colorlet{popgoldL}{popgold!14!white}

\pagecolor{popcream}

% ============================================================
% Typographie
% ============================================================
\defaultfontfeatures{Ligatures=TeX}
\setmainfont{PlusJakartaSans-Medium.ttf}[Path=../../../fonts/,BoldFont=PlusJakartaSans-Bold.ttf]
% Maths en sans-serif (Fira Math) avec chiffres et lettres droites de Plus
% Jakarta, pour que les formules se fondent dans le texte.
\usepackage[math-style=ISO,bold-style=ISO]{unicode-math}
\setmathfont{FiraMath-Regular.otf}[Path=../../../fonts/]
\setmathfont{PlusJakartaSans-Medium.ttf}[Path=../../../fonts/,range={up/num,up/{Latin,latin},\mathup}]
\usepackage{icomma} % virgule décimale sans espace en mode math
% Caractères mathématiques Unicode absents des polices de texte (Plus Jakarta,
% Archivo Black) : rendus par leur équivalent mathématique, en texte comme en
% formule — sinon ils s'affichent en carré vide (ex. « Arithmétique dans ℤ »).
\usepackage{newunicodechar}
\newunicodechar{ℝ}{\ensuremath{\mathbb{R}}}
\newunicodechar{ℤ}{\ensuremath{\mathbb{Z}}}
\newunicodechar{ℂ}{\ensuremath{\mathbb{C}}}
\newunicodechar{ℕ}{\ensuremath{\mathbb{N}}}
\newunicodechar{ℚ}{\ensuremath{\mathbb{Q}}}
\newunicodechar{𝔻}{\ensuremath{\mathbb{D}}}
\newunicodechar{α}{\ensuremath{\alpha}}
\newunicodechar{β}{\ensuremath{\beta}}
\newunicodechar{γ}{\ensuremath{\gamma}}
\newunicodechar{δ}{\ensuremath{\delta}}
\newunicodechar{ε}{\ensuremath{\varepsilon}}
\newunicodechar{θ}{\ensuremath{\theta}}
\newunicodechar{λ}{\ensuremath{\lambda}}
\newunicodechar{μ}{\ensuremath{\mu}}
\newunicodechar{σ}{\ensuremath{\sigma}}
\newunicodechar{τ}{\ensuremath{\tau}}
\newunicodechar{φ}{\ensuremath{\varphi}}
\newunicodechar{ψ}{\ensuremath{\psi}}
\newunicodechar{ω}{\ensuremath{\omega}}
\newunicodechar{Σ}{\ensuremath{\Sigma}}
% majuscules produites par \MakeUppercase dans les titres (α-aminés -> Α)
\newunicodechar{Α}{\ensuremath{\alpha}}
\newunicodechar{Β}{\ensuremath{\beta}}
\newunicodechar{Γ}{\ensuremath{\Gamma}}
\newunicodechar{Δ}{\ensuremath{\Delta}}
\newunicodechar{Ε}{\ensuremath{\varepsilon}}
\newunicodechar{Θ}{\ensuremath{\Theta}}
\newunicodechar{Λ}{\ensuremath{\Lambda}}
\newunicodechar{Μ}{\ensuremath{\mu}}
\newunicodechar{Τ}{\ensuremath{\tau}}
\newunicodechar{Φ}{\ensuremath{\Phi}}
\newunicodechar{Ψ}{\ensuremath{\Psi}}
\newunicodechar{Ω}{\ensuremath{\Omega}}
\newunicodechar{↦}{\ensuremath{\mapsto}}
\newunicodechar{∈}{\ensuremath{\in}}
\newunicodechar{∉}{\ensuremath{\notin}}
\newunicodechar{∧}{\ensuremath{\wedge}}
\newunicodechar{⇔}{\ensuremath{\Leftrightarrow}}
\newunicodechar{⟺}{\ensuremath{\Longleftrightarrow}}
\newunicodechar{⇒}{\ensuremath{\Rightarrow}}
\newunicodechar{⟹}{\ensuremath{\Longrightarrow}}
\newunicodechar{∘}{\ensuremath{\circ}}
\newunicodechar{∪}{\ensuremath{\cup}}
\newunicodechar{∩}{\ensuremath{\cap}}
\newunicodechar{≡}{\ensuremath{\equiv}}
\newunicodechar{⊂}{\ensuremath{\subset}}
\newunicodechar{⊥}{\ensuremath{\perp}}
\newunicodechar{∃}{\ensuremath{\exists}}
\newunicodechar{∀}{\ensuremath{\forall}}
\newunicodechar{∛}{\ensuremath{\sqrt[3]{\,}}}
\newunicodechar{∜}{\ensuremath{\sqrt[4]{\,}}}
\newunicodechar{⃗}{\ensuremath{{}^{\to}}}
\newunicodechar{ⁿ}{\ifmmode^{n}\else\textsuperscript{\ensuremath{n}}\fi}
\newunicodechar{ˣ}{\ifmmode^{x}\else\textsuperscript{\ensuremath{x}}\fi}
\newunicodechar{ᵇ}{\ifmmode^{b}\else\textsuperscript{\ensuremath{b}}\fi}
\newunicodechar{ʳ}{\ifmmode^{r}\else\textsuperscript{\ensuremath{r}}\fi}
\newunicodechar{ᵘ}{\ifmmode^{u}\else\textsuperscript{\ensuremath{u}}\fi}
\newunicodechar{ʸ}{\ifmmode^{y}\else\textsuperscript{\ensuremath{y}}\fi}
\newunicodechar{ᵃ}{\ifmmode^{a}\else\textsuperscript{\ensuremath{a}}\fi}
\newunicodechar{ᵉ}{\ifmmode^{e}\else\textsuperscript{\ensuremath{e}}\fi}
\newunicodechar{ᵏ}{\ifmmode^{k}\else\textsuperscript{\ensuremath{k}}\fi}
\newunicodechar{ᵖ}{\ifmmode^{p}\else\textsuperscript{\ensuremath{p}}\fi}
\newunicodechar{ᴺ}{\ifmmode^{N}\else\textsuperscript{\ensuremath{N}}\fi}
\newunicodechar{ᶜ}{\ifmmode^{c}\else\textsuperscript{\ensuremath{c}}\fi}
\newunicodechar{ʰ}{\ifmmode^{h}\else\textsuperscript{\ensuremath{h}}\fi}
\newunicodechar{ᵠ}{\ifmmode^{q}\else\textsuperscript{\ensuremath{q}}\fi}
\newunicodechar{ₙ}{\ifmmode_{n}\else\textsubscript{\ensuremath{n}}\fi}
\newunicodechar{ₐ}{\ifmmode_{a}\else\textsubscript{\ensuremath{a}}\fi}
\newunicodechar{ᵢ}{\ifmmode_{i}\else\textsubscript{\ensuremath{i}}\fi}
\newunicodechar{ᵣ}{\ifmmode_{r}\else\textsubscript{\ensuremath{r}}\fi}
\newunicodechar{ₖ}{\ifmmode_{k}\else\textsubscript{\ensuremath{k}}\fi}
\newunicodechar{ᵦ}{\ifmmode_{\beta}\else\textsubscript{\ensuremath{\beta}}\fi}
\newunicodechar{ₓ}{\ifmmode_{x}\else\textsubscript{\ensuremath{x}}\fi}
\newfontfamily\popdisplay{ArchivoBlack-Regular.ttf}[Path=../../../fonts/]
\newfontfamily\poplogo{SpaceGrotesk-Bold.ttf}[Path=../../../fonts/]
\newfontfamily\popsemi{PlusJakartaSans-SemiBold.ttf}[Path=../../../fonts/]
\newfontfamily\popxbold{PlusJakartaSans-ExtraBold.ttf}[Path=../../../fonts/]
\newfontfamily\popxboldls{PlusJakartaSans-ExtraBold.ttf}[Path=../../../fonts/,LetterSpace=3.0]

\setlist[enumerate]{leftmargin=1.7em,itemsep=5pt,topsep=4pt}
\setlist[itemize]{leftmargin=1.7em,itemsep=4pt,topsep=4pt,label={\color{poporange}\textbullet}}
\setlength{\parindent}{0pt}
\setlength{\parskip}{5pt}
\renewcommand{\arraystretch}{1.25}

% pgfplots : style Pop par défaut
\pgfplotsset{
  every axis/.append style={axis line style={popink,line width=1pt},tick style={popink},
    grid style={popline,line width=0.5pt},label style={font=\small},tick label style={font=\footnotesize}},
  popcurve/.style={poporange,line width=1.8pt,no markers,smooth},
}
% Même style disponible dans un \draw TikZ simple (hors pgfplots)
\tikzset{popcurve/.style={poporange,line width=1.8pt}}
\tikzset{popgrid/.style={popline,very thin}}
\colorlet{popcurve}{poporange} % le nom du style est parfois utilisé comme couleur

% ============================================================
% Pill / squiggle
% ============================================================
\newcommand{\poppill}[3]{%
  \tikz[baseline=-0.55ex]\node[draw=popink,line width=1.2pt,rounded corners=2.6pt,%
    fill=#2,inner xsep=6.5pt,inner ysep=3pt]{%
    \popxboldls\fontsize{7}{7}\selectfont\color{#3}\MakeUppercase{#1}};%
}
\newcommand{\popsquiggle}[1]{%
  \tikz[baseline=-0.4ex]{
    \draw[#1,line width=2.6pt,line cap=round]
      plot [smooth] coordinates {(0,0.09) (0.34,0.01) (0.68,0.21) (1.02,0.01) (1.36,0.10)};
  }%
}
% Mot-clé surligné (marqueur orange sous le texte)
\newcommand{\cle}[1]{{\popxbold\color{popdark}#1}}

% ============================================================
% En-tête / pied
% ============================================================
\pagestyle{fancy}
\fancyhf{}
\renewcommand{\headrulewidth}{0pt}
\renewcommand{\footrulewidth}{0pt}
\lhead{\raisebox{-0.15ex}{\poplogo\fontsize{13}{13}\selectfont\color{popink}OnBuch\color{poporange}+}}
\rhead{\poppill{\DOCMATIERE\ \textbullet\ \DOCNIVEAU\ \textbullet\ Leçon \DOCLECON}{white}{popink}}
\lfoot{\popxbold\fontsize{7.6}{7.6}\selectfont\color{popmuted}\MakeUppercase{Cours signé OnBuch+}\ \textbullet\ {\popsemi\DOCTITRE}}
\rfoot{\tikz[baseline=-0.6ex]\node[rounded corners=8.5pt,fill=popink,minimum height=17pt,minimum width=17pt,inner xsep=5pt,inner sep=0]{\popxbold\fontsize{8}{8}\selectfont\color{popcream}\thepage\,/\,\pageref*{LastPage}};}

% ============================================================
% Titres
% ============================================================
\newcommand{\chaptitre}[2]{%
  \begin{tcolorbox}[enhanced,colback=poporange,colframe=popink,boxrule=2.6pt,arc=11pt,
      drop shadow=popink,left=15pt,right=15pt,top=12pt,bottom=11pt,before skip=3pt,after skip=18pt]
    \poppill{#2}{popink}{popcream}\par\vspace{9pt}
    {\raggedright\hyphenpenalty=10000\exhyphenpenalty=10000\popdisplay\fontsize{22}{24}\selectfont\color{popcream}\MakeUppercase{#1}\par}
  \end{tcolorbox}
}
% Section numérotée : pastille orange avec le numéro + titre Archivo
\newcounter{popsec}
\newcommand{\coursec}[1]{%
  \stepcounter{popsec}\setcounter{popsub}{0}%
  \par\vspace{16pt}\noindent
  \begin{minipage}{\linewidth}
  \tikz[baseline=-0.7ex]\node[draw=popink,line width=1.6pt,fill=poporange,rounded corners=4pt,minimum size=22pt,inner sep=0]{\popdisplay\fontsize{12}{12}\selectfont\color{popcream}\thepopsec};%
  \hspace{8pt}{\popdisplay\fontsize{15}{17}\selectfont\color{popink}\MakeUppercase{#1}}\par
  \vspace{4pt}\noindent{\color{poporange}\rule{36pt}{3.6pt}}
  \end{minipage}\par\nopagebreak\vspace{8pt}
}
\newcounter{popsub}
\newcommand{\courssub}[1]{%
  \stepcounter{popsub}%
  \par\vspace{9pt}\noindent{\popxbold\fontsize{11.5}{13}\selectfont\color{popink}{\color{poporange}\thepopsec.\thepopsub}\hspace{6pt}#1}\par\nopagebreak\vspace{4pt}
}

% ============================================================
% Boîtes pédagogiques — même recette : fond blanc, bordure épaisse colorée,
% ombre dure encre, badge plein. Titre optionnel [#1].
% ============================================================
\tcbset{popbase/.style={breakable,enhanced,colback=white,boxrule=2.3pt,arc=9pt,
  shadow={3pt}{-3pt}{0pt}{popink},left=14pt,right=14pt,top=11pt,bottom=11pt,before skip=11pt,after skip=15pt,
  fontupper=\popsemi\fontsize{10.6}{14.5}\selectfont\color{popink},
  fontlower=\popsemi\fontsize{10.6}{14.5}\selectfont\color{popink}}}
% Coche dessinée (le glyphe ✓ n'existe pas dans les polices du document)
\newcommand{\popcheck}[1]{\tikz[baseline=-0.5ex]\draw[#1,line width=1.6pt,line cap=round,line join=round](-0.13,0.0)--(-0.04,-0.09)--(0.13,0.1);}
\providecommand{\coche}{\popcheck{popgreen}}
\providecommand{\resultat}[1]{\par\smallskip\noindent\fcolorbox{popgreen}{popgreenL}{\parbox{\dimexpr\linewidth-2\fboxsep-2\fboxrule\relax}{\textbf{Résultat.} #1}}\par\smallskip}
\newcommand{\popboxhead}[3]{\poppill{#1}{#2}{white}\ifx\relax#3\relax\else\hspace{6pt}{\popxbold\fontsize{10.6}{12}\selectfont\color{popink}#3}\fi\par\vspace{7pt}}

\newtcolorbox{aretenir}[1][]{popbase,colframe=popgreen,before upper={\popboxhead{À retenir}{popgreen}{#1}}}
% Texte en espagnol (document de lecture, dialogue, poème, extrait) : cadre
% d'étude. Un titre optionnel (source, auteur) s'affiche dans l'en-tête.
\newtcolorbox{textebox}[1][]{popbase,colframe=popblue,before upper={\popboxhead{Texto}{popblue}{#1}\begin{otherlanguage}{spanish}},after upper={\end{otherlanguage}}}
% Marques ✓ / ✗ (forme correcte / fautive) souvent utilisées par les modèles
\providecommand{\cross}{\textcolor{poppink}{\ensuremath{\times}}}
\providecommand{\xmark}{\cross}\providecommand{\crossmark}{\cross}\providecommand{\cmark}{\ensuremath{\checkmark}}
% Ligne de réponse / justification à compléter (inventée par les modèles)
\providecommand{\justification}[1]{\par\noindent\textit{Justification~:} \dotfill\par}
% \textipa (package tipa non chargé) : la police ne couvre pas l'API -> simple passage
\providecommand{\textipa}[1]{#1}
% Soulignement (ulem non chargé) : \uline, \ul
\providecommand{\uline}[1]{\underline{#1}}\providecommand{\ul}[1]{\underline{#1}}
% Mot ou phrase en espagnol à l'intérieur d'un paragraphe français
\newcommand{\esp}[1]{\ifmmode\text{\textit{#1}}\else\begingroup\selectlanguage{spanish}\itshape #1\endgroup\fi}
\newtcolorbox{exemplebox}[1][]{popbase,colframe=poporange,before upper={\popboxhead{Exemple}{poporange}{#1}}}
\newtcolorbox{definition}[1][]{popbase,colframe=popblue,before upper={\popboxhead{Définition}{popblue}{#1}}}
\newtcolorbox{propriete}[1][]{popbase,colframe=poppurple,before upper={\popboxhead{Propriété}{poppurple}{#1}}}
\newtcolorbox{methode}[1][]{popbase,colframe=popgold,before upper={\popboxhead{Méthode}{popgold}{#1}}}
\newtcolorbox{attention}[1][]{popbase,colframe=poppink,before upper={\popboxhead{Attention}{poppink}{#1}}}
\newtcolorbox{experience}[1][]{popbase,colframe=popblue,colback=popblueL,before upper={\popboxhead{Expérience}{popblue}{#1}}}
% « Le savais-tu ? » : carte violette pleine (contexte, culture, Cameroun)
\newtcolorbox{savaistu}[1][]{popbase,colback=poppurple,colframe=popink,
  fontupper=\popsemi\fontsize{10.4}{14.2}\selectfont\color{white},
  before upper={\poppill{Le savais-tu\,?}{white}{poppurple}\ifx\relax#1\relax\else\hspace{6pt}{\popxbold\color{white}#1}\fi\par\vspace{7pt}}}
% Exercice résolu : énoncé en haut, \tcblower puis la solution sur fond crème
\newcounter{popexo}\newcounter{popexr}
\newtcolorbox{exoresolu}[1][]{popbase,colframe=poporange,colbacklower=popcream,
  segmentation style={popink,line width=1.2pt,dashed},
  before upper={\stepcounter{popexr}\popboxhead{Exercice résolu \thepopexr}{poporange}{#1}},
  before lower={\poppill{Solution}{popink}{popcream}\par\vspace{6pt}}}
% Exercice (non résolu en place) — la correction est dans la partie Corrigés
\newtcolorbox{exercice}[1][]{popbase,colframe=popink,
  before upper={\stepcounter{popexo}\popboxhead{Exercice \thepopexo}{popink}{#1}}}
% Corrigé
\newtcolorbox{corrige}[1][]{popbase,colframe=popgreen,colback=popgreenL,
  before upper={\popboxhead{Corrigé}{popgreen}{#1}}}
% Objectifs / prérequis / bilan
\newtcolorbox{objectifs}{popbase,colframe=popink,colback=poporangeL,
  before upper={\popboxhead{Objectifs de la leçon}{popink}{}}}
\newtcolorbox{prerequis}{popbase,colframe=popink,colback=popblueL,
  before upper={\popboxhead{Prérequis}{popblue}{}}}
\newtcolorbox{bilan}{popbase,colframe=popink,colback=white,boxrule=2.8pt,
  before upper={\popboxhead{Fiche bilan}{poporange}{L'essentiel à connaître pour l'examen}}}
% Figure encadrée avec légende
\newtcolorbox{popfigure}[1][]{enhanced,breakable=false,colback=white,colframe=popline,boxrule=1.4pt,arc=8pt,
  left=10pt,right=10pt,top=10pt,bottom=8pt,before skip=10pt,after skip=12pt,halign=center,
  lower separated=false,halign lower=center,
  fontlower=\popsemi\fontsize{9}{11.5}\selectfont\color{popmuted},
  #1}
\newcommand{\legende}[1]{\par\vspace{6pt}{\popsemi\fontsize{9}{11.5}\selectfont\color{popmuted}#1}}

% ============================================================
% Signature OnBuch+
% ============================================================
\newcommand{\poplogoinline}[1]{{\poplogo\fontsize{#1}{#1}\selectfont\color{popink}OnBuch\color{poporange}+}}
\newcommand{\signatureonbuch}{%
  \begin{tcolorbox}[enhanced,colback=popink,colframe=popink,boxrule=2.2pt,arc=10pt,
      drop shadow=poporange,left=14pt,right=14pt,top=10pt,bottom=10pt,before skip=0pt,after skip=0pt]
    \tikz[baseline=-0.7ex]\node[circle,fill=popgreen,draw=popcream,line width=1.3pt,minimum size=22pt,inner sep=0]{\popcheck{white}};%
    \hspace{9pt}\begin{minipage}[c]{0.62\linewidth}
      {\popxbold\fontsize{12}{13}\selectfont\color{popcream}Signé\ \textbullet\ {\poplogo OnBuch\color{poporange}+}}\par\vspace{2pt}
      {\raggedright\popsemi\fontsize{8}{10.5}\selectfont\color{popline}\MakeUppercase{Équipe pédagogique OnBuch+}\\\MakeUppercase{Conforme au programme officiel MINESEC}\par}
    \end{minipage}\hfill
    {\poplogo\fontsize{22}{22}\selectfont\color{popcream}OnBuch\color{poporange}+}
  \end{tcolorbox}%
}

% ============================================================
% Page de garde
% #1 titre · #2 matière · #3 niveau · #4 module · #5 n° leçon · #6 durée ·
% #7 description / objectifs (texte ou itemize)
% ============================================================
\newcommand{\pagegarde}[7]{%
  \thispagestyle{empty}
  \newgeometry{a4paper,margin=1.7cm,top=1.6cm,bottom=1.4cm}%
  \begin{tikzpicture}[remember picture,overlay]
    \fill[poporange!18] ([xshift=-3.2cm,yshift=-3.6cm]current page.north east) circle (4.6cm);
    \fill[poppurple!14] ([xshift=2.4cm,yshift=7.6cm]current page.south west) circle (3.8cm);
  \end{tikzpicture}%
  {\poplogo\fontsize{30}{30}\selectfont\color{popink}OnBuch\color{poporange}+}\hfill
  \raisebox{0.6ex}{\poppill{Cours signé \textbullet\ Premium}{popink}{popcream}}\par
  \vspace{1pt}\popsquiggle{poppurple}\par
  \vspace{8pt}
  \begin{tcolorbox}[enhanced,colback=poporange,colframe=popink,boxrule=2.8pt,arc=13pt,
      drop shadow=popink,left=18pt,right=18pt,top=12pt,bottom=13pt,after skip=10pt]
    \poppill{#2 \textbullet\ #3}{popink}{popcream}\hspace{5pt}\poppill{#4}{white}{popink}\par\vspace{8pt}
    {\popdisplay\fontsize{14}{14}\selectfont\color{popink}LEÇON #5}\par\vspace{4pt}
    {\raggedright\hyphenpenalty=10000\exhyphenpenalty=10000\popdisplay\fontsize{25}{27}\selectfont\color{popcream}\MakeUppercase{#1}\par}
    \vspace{8pt}
    {\popxbold\fontsize{9.5}{9.5}\selectfont\color{popink}\MakeUppercase{Durée conseillée : #6}}
  \end{tcolorbox}
  \begin{tcolorbox}[enhanced,colback=white,colframe=popink,boxrule=2.2pt,arc=10pt,
      drop shadow=popink,left=16pt,right=16pt,top=10pt,bottom=10pt,after skip=8pt]
    \poppill{Dans cette leçon}{poppurple}{white}\par\vspace{5pt}
    \popsemi\fontsize{9.3}{12.6}\selectfont\color{popink}#7
  \end{tcolorbox}
  \noindent\begin{minipage}[t]{0.487\linewidth}
    \begin{tcolorbox}[enhanced,colback=popgreen,colframe=popink,boxrule=2.2pt,arc=10pt,
        drop shadow=popink,left=11pt,right=11pt,top=10pt,bottom=10pt,height=3.1cm,valign=center]
      \tcbox[colback=white,colframe=popink,boxrule=1.3pt,arc=5pt,boxsep=0pt,nobeforeafter,
        left=3pt,right=3pt,top=3pt,bottom=3pt,valign=center]{\qrcode[height=1.9cm]{https://play.google.com/store/apps/details?id=com.luvvix.onbuch&hl=fr}}%
      \hspace{8pt}\begin{minipage}[c]{0.5\linewidth}
        {\popdisplay\fontsize{11}{12.5}\selectfont\color{white}\MakeUppercase{Télécharge OnBuch}}\par\vspace{4pt}
        {\raggedright\popsemi\fontsize{8.4}{11}\selectfont\color{white}Gratuit \textbullet\ Google Play\\Tuteur IA, annales, concours, résultats\par}
      \end{minipage}
    \end{tcolorbox}
  \end{minipage}\hfill
  \begin{minipage}[t]{0.487\linewidth}
    \begin{tcolorbox}[enhanced,colback=poppurple,colframe=popink,boxrule=2.2pt,arc=10pt,
        drop shadow=popink,left=11pt,right=11pt,top=10pt,bottom=10pt,height=3.1cm,valign=center]
      \tikz[baseline=-0.7ex]\node[circle,fill=white,draw=popink,line width=1.3pt,minimum size=1.6cm,inner sep=0]{\popdisplay\fontsize{24}{24}\selectfont\color{poppurple}+};%
      \hspace{8pt}\begin{minipage}[c]{0.6\linewidth}
        {\popdisplay\fontsize{11}{12.5}\selectfont\color{white}\MakeUppercase{Passe à OnBuch+}}\par\vspace{4pt}
        {\raggedright\popsemi\fontsize{8.4}{11}\selectfont\color{white}Tous les cours signés, Tuteur IA illimité, fiches PDF — onglet OnBuch+ de l'app\par}
      \end{minipage}
    \end{tcolorbox}
  \end{minipage}
  \vspace{8pt}
  \popcontenu
  \vfill
  \signatureonbuch
  \restoregeometry
  \newpage
}

% Rangée « Ce cours contient » (page de garde)
\newcommand{\poptile}[3]{% #1 couleur #2 gros texte #3 légende
  \begin{tcolorbox}[enhanced,colback=#1,colframe=popink,boxrule=2pt,arc=9pt,shadow={2.5pt}{-2.5pt}{0pt}{popink},
      left=6pt,right=6pt,top=9pt,bottom=9pt,halign=center,valign=center,height=1.75cm,width=\linewidth,nobeforeafter]
    {\popdisplay\fontsize{12.5}{13}\selectfont\color{white}\MakeUppercase{#2}}\par\vspace{3pt}
    {\centering\popsemi\fontsize{7.8}{9.5}\selectfont\color{white}#3\par}
  \end{tcolorbox}}
\newcommand{\popcontenu}{%
  \par\noindent{\popxboldls\fontsize{7.5}{7.5}\selectfont\color{popmuted}CE COURS SIGNÉ CONTIENT}\par\vspace{7pt}
  \noindent\begin{minipage}[t]{0.235\linewidth}\poptile{popblue}{Cours}{complet, illustré, pas à pas}\end{minipage}\hfill
  \begin{minipage}[t]{0.235\linewidth}\poptile{poporange}{Méthodes}{et exercices résolus}\end{minipage}\hfill
  \begin{minipage}[t]{0.235\linewidth}\poptile{poppink}{Exercices}{type examen + corrigés}\end{minipage}\hfill
  \begin{minipage}[t]{0.235\linewidth}\poptile{popgreen}{Fiche bilan}{l'essentiel à retenir}\end{minipage}\par}

% Sommaire
\newtcolorbox{sommaire}{enhanced,colback=white,colframe=popink,boxrule=2.2pt,arc=10pt,
  drop shadow=popink,left=18pt,right=18pt,top=13pt,bottom=13pt,before skip=6pt,after skip=20pt,
  before upper={\poppill{Sommaire}{poppurple}{white}\par\vspace{9pt}\popsemi\fontsize{10.6}{14.5}\selectfont\color{popink}}}

% Signature de fin de cours — poussée tout en bas de la dernière page
\newcommand{\findecours}{%
  \par\vfill\nopagebreak
  \begin{center}\popsquiggle{poporange}\end{center}\vspace{4pt}
  \signatureonbuch
}
\AtBeginDocument{\def\llbracket{\mathopen{[\mkern-3.5mu[}}\def\rrbracket{\mathclose{]\mkern-3.5mu]}}} % intervalles d'entiers (glyphe ⟦ absent de la police)
% Glyphes absents de Fira Math : redéfinis à partir de glyphes disponibles
\AtBeginDocument{%
  \def\setminus{\mathbin{\backslash}}\let\smallsetminus\setminus
  \def\vdots{\mathord{\vbox{\baselineskip3.2pt\lineskiplimit0pt\kern4pt\hbox{.}\hbox{.}\hbox{.}}}}%
  \def\ddots{\mathinner{\mkern1mu\raise6.5pt\hbox{.}\mkern2mu\raise3.6pt\hbox{.}\mkern2mu\raise0.7pt\hbox{.}\mkern1mu}}%
  \def\triangle{\mathord{\scalebox{0.9}{$\symup{\Delta}$}}}%
  \def\nabla{\mathord{\raisebox{\depth}{\scalebox{1}[-1]{$\symup{\Delta}$}}}}%
  \def\checkmark{\popcheck{popdark}}%
  \def\star{\mathbin{\ast}}\def\bigstar{\mathord{\ast}}%
  \def\diamond{\mathbin{\scalebox{0.6}{$\Diamond$}}}%
  \def\bigcup{\mathop{\mathchoice{\raisebox{-0.3em}{\scalebox{1.7}{$\cup$}}}{\scalebox{1.25}{$\cup$}}{\cup}{\cup}}\displaylimits}%
  \def\bigcap{\mathop{\mathchoice{\raisebox{-0.3em}{\scalebox{1.7}{$\cap$}}}{\scalebox{1.25}{$\cap$}}{\cap}{\cap}}\displaylimits}%
  \def\lesseqqgtr{\mathrel{\lessgtr}}\def\bigoplus{\mathop{\oplus}\displaylimits}%
}
\providecommand{\ssi}{si et seulement si} % abréviation fréquente
\AtBeginDocument{\providecommand{\wideparen}{\overparen}} % arc de cercle

%% FIGFIX-BEGIN v3 — correctifs globaux des figures (docs/onbuch-generation/tools/figfix)
\usepackage{adjustbox}\usepackage{etoolbox}
% toute figure TikZ/pgfplots plus large que la ligne est réduite à la largeur disponible
\BeforeBeginEnvironment{tikzpicture}{\begin{adjustbox}{max width=\linewidth}}
\AfterEndEnvironment{tikzpicture}{\end{adjustbox}}
% légendes pgfplots lisibles par-dessus les courbes
\pgfplotsset{every axis/.append style={legend style={fill=white,fill opacity=0.92,text opacity=1,draw=black!15,inner sep=3pt}}}
% étiquettes d'arêtes (syntaxe edge["..."]) avec fond blanc
\tikzset{every edge quotes/.append style={fill=white,inner sep=1.2pt}}
% bulles de cartes mentales : texte à la ligne dans la bulle, bulle un peu plus grande
\tikzset{every concept/.append style={text width=2.0cm,align=center,minimum size=2.3cm,inner sep=2pt}}
%% FIGFIX-END
\begin{document}

\pagegarde{Traduction : comparatifs et superlatifs ; redacción : el currículum vítae}{Espagnol}{Tle A}{Módulo 4 : Activités économiques et monde du travail}{E17}{2 h}{Leçon mixte du Módulo 4 : d'une part, la grammaire de la comparaison (comparatifs réguliers et irréguliers, superlatifs relatif et absolu) avec ses règles, ses pièges et ses exercices de traduction ; d'autre part, la rédaction du currículum vítae en espagnol, outil indispensable du monde du travail.
\vspace{4pt}
\begin{multicols}{2}\small
\begin{itemize}
\item À la fin de cette leçon, je sais former et traduire les comparatifs de supériorité, d'infériorité et d'égalité (más… que, menos… que, tan… como, tanto… como)
\item À la fin de cette leçon, je sais employer les comparatifs irréguliers mejor, peor, mayor, menor et distinguer leurs emplois
\item À la fin de cette leçon, je sais construire le superlatif relatif (el/la/los/las más… de) et le superlatif absolu (-ísimo, muy, sumamente)
\item À la fin de cette leçon, je sais éviter les pièges francophones : « plus de » devant un nom, « que » vs « de » devant un chiffre, l'accord de tanto
\item À la fin de cette leçon, je sais traduire un texte du monde du travail contenant des comparaisons (version) et traduire des phrases du français vers l'espagnol (thème)
\item À la fin de cette leçon, je sais identifier les rubriques d'un CV espagnol : datos personales, formación académica, experiencia profesional, idiomas, competencias, otros datos
\end{itemize}\end{multicols}}

\begin{sommaire}
\begin{multicols}{2}
\begin{enumerate}
\item Observation : la comparaison en contexte
\item Les comparatifs réguliers et irréguliers
\item Les superlatifs : relatif et absolu
\item Traduction guidée : thème et version
\item Le currículum vítae en espagnol
\item Production guidée : mon CV et mon autoportrait
\item Activité d'intégration
\item Méthodes et exercices résolus
\item Exercices
\item Corrigés des exercices
\item Fiche bilan
\end{enumerate}
\end{multicols}
\end{sommaire}

\begin{objectifs}
\begin{itemize}
\item À la fin de cette leçon, je sais former et traduire les comparatifs de supériorité, d'infériorité et d'égalité (más… que, menos… que, tan… como, tanto… como)
\item À la fin de cette leçon, je sais employer les comparatifs irréguliers mejor, peor, mayor, menor et distinguer leurs emplois
\item À la fin de cette leçon, je sais construire le superlatif relatif (el/la/los/las más… de) et le superlatif absolu (-ísimo, muy, sumamente)
\item À la fin de cette leçon, je sais éviter les pièges francophones : « plus de » devant un nom, « que » vs « de » devant un chiffre, l'accord de tanto
\item À la fin de cette leçon, je sais traduire un texte du monde du travail contenant des comparaisons (version) et traduire des phrases du français vers l'espagnol (thème)
\item À la fin de cette leçon, je sais identifier les rubriques d'un CV espagnol : datos personales, formación académica, experiencia profesional, idiomas, competencias, otros datos
\item À la fin de cette leçon, je sais rédiger mon propre CV en espagnol avec un vocabulaire précis et des formulations idiomatiques
\item À la fin de cette leçon, je sais décrire oralement mon parcours scolaire et mes compétences en vue d'un entretien
\end{itemize}
\end{objectifs}

\begin{prerequis}
\begin{itemize}
\item L'accord de l'adjectif en genre et en nombre (classe de Première)
\item Les adjectifs qualificatifs courants et leur place dans la phrase
\item Les articles définis et indéfinis, les possessifs
\item Le présent de l'indicatif et le passé composé (pretérito perfecto) pour décrire son parcours
\item Le vocabulaire de base de l'école et des professions (Módulo 4, leçons précédentes)
\end{itemize}
\end{prerequis}

\newpage

\coursec{Observation : la comparaison en contexte}

\courssub{Situation d'accroche : deux candidats, une seule place}

C'est la rentrée à Yaoundé. Ton grand frère, fraîchement diplômé, postule à un stage dans une entreprise espagnole installée au port de Douala. Deux surprises l'attendent : on lui demande un \esp{currículum vítae} rédigé entièrement en espagnol, et il doit comparer deux offres d'emploi qu'on lui propose en même temps. Sur la table, deux annonces :

\esp{Este puesto es mejor pagado, pero el otro es menos exigente.} — « Ce poste est mieux payé, mais l'autre est moins exigeant. »

Comment traduit-il correctement ces comparaisons ? Comment dit-on « plus expérimenté que », « moins de temps que », « autant de langues que », « le meilleur candidat » ? Et comment présente-t-il son parcours, ses diplômes et ses compétences dans un CV qui respecte les codes du monde hispanophone ?

\textbf{Problématique de la leçon :} comment la langue espagnole construit-elle la comparaison et le superlatif, et comment ces structures permettent-elles de décrire un parcours professionnel et de rédiger un \esp{currículum vítae} efficace ?

Dans cette première partie, nous allons d'abord \emph{observer} la comparaison en contexte, dans un texte authentique du monde du travail, avant d'en \emph{déduire} les structures et les règles.

\courssub{Lecture : dos candidatos para un puesto}

Lis attentivement le texte suivant. Il s'agit d'un extrait du compte rendu d'un jury de recrutement dans une entreprise d'import-export de Douala, qui hésite entre deux candidats : Ana Mbarga et Luis Esono.

\begin{textebox}[Dos candidatos para un puesto]
El departamento de recursos humanos de la empresa CamerExport ha recibido más de cien candidaturas para el puesto de asistente comercial. Después de la primera selección, quedan dos finalistas: Ana Mbarga y Luis Esono.

Ana es más experimentada que Luis: ha trabajado cinco años en una empresa de logística en Douala, mientras que Luis tiene menos experiencia que ella. Sin embargo, Luis habla tantos idiomas como Ana: los dos dominan el español, el francés y el inglés. Además, Luis es menor que Ana — tiene veintitrés años — y los directivos piensan que es más dinámico que los otros candidatos.

En la entrevista, Ana ha sido la candidata más preparada de todas. Su carta de motivación es buenísima: está mejor redactada que la de Luis y es más convincente que cualquier otra. Luis, en cambio, ha respondido peor a las preguntas técnicas, pero ha mostrado más entusiasmo que nadie.

El salario propuesto es más alto que el del año pasado, pero el puesto exige trabajar más horas que antes. El director afirma: « Necesitamos al candidato más completo, no necesariamente al más joven ni al más barato ».

Mañana, el jurado tomará la decisión final. ¿Quién conseguirá el puesto? ¿La candidata con más experiencia o el candidato con más energía?
\end{textebox}

\textbf{Glossaire :}

\begin{center}
\begin{tabularx}{\linewidth}{|X|X|}
\hline
\rowcolor{popblueL} \textbf{Español} & \textbf{Français} \\
\hline
recursos humanos & ressources humaines \\
\hline
el puesto (de trabajo) & le poste (de travail) \\
\hline
la candidatura & la candidature \\
\hline
los finalistas & les finalistes \\
\hline
dominar (un idioma) & maîtriser (une langue) \\
\hline
la entrevista & l'entretien (d'embauche) \\
\hline
la carta de motivación & la lettre de motivation \\
\hline
redactar & rédiger \\
\hline
convincente & convaincant(e) \\
\hline
el salario / el sueldo & le salaire \\
\hline
exigir & exiger \\
\hline
tomar una decisión & prendre une décision \\
\hline
conseguir el puesto & obtenir le poste \\
\hline
\end{tabularx}
\end{center}

\textbf{Questions de compréhension :}

\begin{enumerate}
\item ¿Cuántas candidaturas ha recibido la empresa? (« Combien de candidatures l'entreprise a-t-elle reçues ? »)
\item ¿Quién tiene más experiencia, Ana o Luis? (« Qui a le plus d'expérience ? »)
\item ¿Qué idiomas hablan los dos candidatos? (« Quelles langues parlent les deux candidats ? »)
\item Según el director, ¿qué tipo de candidato necesita la empresa? (« Selon le directeur, quel type de candidat l'entreprise recherche-t-elle ? »)
\end{enumerate}

\begin{corrige}[Questions de compréhension]
\begin{enumerate}
\item Ha recibido más de cien candidaturas. — Elle a reçu plus de cent candidatures.
\item Ana tiene más experiencia que Luis. — Ana a plus d'expérience que Luis.
\item Hablan español, francés e inglés. — Ils parlent espagnol, français et anglais.
\item Necesita al candidato más completo. — Elle a besoin du candidat le plus complet (le plus polyvalent).
\end{enumerate}
\end{corrige}

\courssub{Repérage guidé : souligner et classer les comparaisons}

Relisons maintenant le texte avec un regard de grammairien. Souligne (ou recopie) toutes les structures qui expriment une \cle{comparaison} ou un \cle{superlatif}. En voici la liste complète, dans l'ordre d'apparition :

\begin{enumerate}
\item \esp{más experimentada que Luis} — plus expérimentée que Luis
\item \esp{menos experiencia que ella} — moins d'expérience qu'elle
\item \esp{tantos idiomas como Ana} — autant de langues qu'Ana
\item \esp{menor que Ana} — plus jeune (litt. « moindre ») qu'Ana
\item \esp{más dinámico que los otros} — plus dynamique que les autres
\item \esp{la candidata más preparada de todas} — la candidate la mieux préparée de toutes
\item \esp{buenísima} — excellente (extrêmement bonne)
\item \esp{mejor redactada que la de Luis} — mieux rédigée que celle de Luis
\item \esp{más convincente que cualquier otra} — plus convaincante que toute autre
\item \esp{peor} — plus mal (pire)
\item \esp{más entusiasmo que nadie} — plus d'enthousiasme que personne
\item \esp{más alto que el del año pasado} — plus élevé que celui de l'année dernière
\item \esp{más horas que antes} — plus d'heures qu'avant
\item \esp{el candidato más completo} — le candidat le plus complet
\item \esp{al más joven ni al más barato} — le plus jeune ni le moins cher
\end{enumerate}

\begin{exemplebox}[Tableau de classement des structures repérées]
\begin{center}
\begin{tabularx}{\linewidth}{|X|X|X|}
\hline
\rowcolor{popblueL} \textbf{Supériorité (más… que)} & \textbf{Infériorité (menos… que)} & \textbf{Égalité (tan/tanto… como)} \\
\hline
más experimentada que Luis & menos experiencia que ella & tantos idiomas como Ana \\
\hline
más dinámico que los otros & peor (a las preguntas) & \\
\hline
más convincente que cualquier otra & & \\
\hline
más alto que el del año pasado & & \\
\hline
más horas que antes & & \\
\hline
más entusiasmo que nadie & & \\
\hline
\rowcolor{popgoldL} \multicolumn{3}{|l|}{\textbf{Superlatifs :} la candidata más preparada de todas ; buenísima ; el candidato más completo ; el más joven ; el más barato} \\
\hline
\end{tabularx}
\end{center}
\end{exemplebox}

\begin{attention}[Ne confonds pas les trois rapports]
Le français et l'espagnol construisent les mêmes trois rapports de comparaison, mais avec des mots-outils différents : \esp{más… que} = plus… que ; \esp{menos… que} = moins… que ; \esp{tan… como} = aussi… que. Le second terme est introduit par \esp{que} (supériorité/infériorité) ou par \esp{como} (égalité). Ne dis jamais \esp{más… como} : c'est un calque impossible.
\end{attention}

\courssub{Déduction : quelles structures se dégagent ?}

Observons les exemples avant d'énoncer la règle. Regroupe mentalement les phrases repérées : que remarques-tu ?

\begin{itemize}
\item Devant l'adjectif ou l'adverbe comparé, on trouve toujours \esp{más}, \esp{menos} ou \esp{tan} : \esp{más alto}, \esp{menos exigente}, \esp{tan preparada}.
\item Le second terme de la comparaison est introduit par \esp{que} ou \esp{como} : \esp{más alto que…}, \esp{tantos idiomas como…}.
\item Pour l'égalité avec un \emph{nom}, le mot \esp{tanto} s'accorde : \esp{tantos idiomas como} (masculin pluriel).
\item Certaines formes ne suivent pas le schéma \esp{más + adjectif} : \esp{mejor} (meilleur/mieux), \esp{peor} (pire/plus mal), \esp{menor} (plus jeune/moindre). Ce sont des \cle{comparatifs irréguliers}.
\item Pour exprimer « le plus… de », l'espagnol utilise l'article + \esp{más} + adjectif + \esp{de} : \esp{la candidata más preparada de todas}. C'est le \cle{superlatif relatif}.
\item La forme \esp{buenísima} (bon + ísima) exprime un degré très élevé sans comparaison : c'est le \cle{superlatif absolu}.
\end{itemize}

\begin{propriete}[Première synthèse : les structures de la comparaison]
\begin{itemize}
\item \textbf{Supériorité :} \esp{más + adjectif/adverbe/nom + que} — \esp{Este trabajo es más difícil que el otro.} « Ce travail est plus difficile que l'autre. »
\item \textbf{Infériorité :} \esp{menos + adjectif/adverbe/nom + que} — \esp{Gana menos dinero que su hermano.} « Il gagne moins d'argent que son frère. »
\item \textbf{Égalité :} \esp{tan + adjectif/adverbe + como} ; \esp{tanto/a/os/as + nom + como} — \esp{Luis habla tantos idiomas como ella.} « Luis parle autant de langues qu'elle. »
\item \textbf{Superlatif relatif :} article + \esp{más/menos} + adjectif + \esp{de} — \esp{Ana es la candidata más preparada de todas.} « Ana est la candidate la mieux préparée de toutes. »
\item \textbf{Superlatif absolu :} adjectif + \esp{-ísimo/a/os/as} ou \esp{muy} + adjectif — \esp{Su carta es buenísima.} « Sa lettre est excellente. »
\end{itemize}
\end{propriete}

\begin{attention}[Sur quoi porte la comparaison ?]
La comparaison ne porte pas seulement sur des adjectifs. Elle peut porter sur :
\begin{itemize}
\item un \textbf{adjectif} : \esp{más difícil que} (plus difficile que) ;
\item un \textbf{adverbe} : \esp{habla más despacio que yo} (il parle plus lentement que moi) ;
\item un \textbf{nom} : \esp{menos experiencia que ella} (moins d'expérience qu'elle) ;
\item un \textbf{verbe} : \esp{trabaja más que antes} (il travaille plus qu'avant).
\end{itemize}
En français aussi : « il travaille \emph{plus} qu'avant », « il a \emph{autant de} courage que toi ». Repère toujours la nature du mot comparé avant de traduire.
\end{attention}

\begin{methode}[Analyser une comparaison en trois étapes]
\begin{enumerate}
\item \textbf{Identifier l'élément comparé} : est-ce un adjectif (\esp{alto}), un adverbe (\esp{despacio}), un nom (\esp{experiencia}) ou un verbe (\esp{trabaja}) ?
\item \textbf{Repérer le second terme} : il est introduit par \esp{que} (rapport inégal) ou \esp{como} (rapport d'égalité) ; au superlatif relatif, le groupe est introduit par \esp{de} (\esp{de todas}).
\item \textbf{Déterminer le type de rapport} : supériorité (\esp{más}), infériorité (\esp{menos}) ou égalité (\esp{tan/tanto… como}) — puis choisir l'équivalent français : plus… que, moins… que, aussi… que / autant de… que.
\end{enumerate}
\end{methode}

\begin{exoresolu}[Analyse guidée d'une phrase du texte]
Énoncé : analyse la phrase \esp{Su carta está mejor redactada que la de Luis y es más convincente que cualquier otra.} selon les trois étapes de la méthode, puis traduis-la.
\tcblower
Solution détaillée :
\begin{enumerate}
\item Éléments comparés : \esp{redactada} (participe employé comme adjectif) et \esp{convincente} (adjectif).
\item Seconds termes : \esp{que la de Luis} (« que celle de Luis » — \esp{la} reprend \esp{carta}) et \esp{que cualquier otra} (« que toute autre »).
\item Rapports : \esp{mejor… que} = supériorité avec comparatif irrégulier (\esp{mejor} = mieux) ; \esp{más… que} = supériorité régulière.
\end{enumerate}
Traduction : « Sa lettre est mieux rédigée que celle de Luis et elle est plus convaincante que toute autre. »
Remarque : le français reprend le nom par « celle de », l'espagnol par l'article \esp{la de} : même procédé, même logique.
\end{exoresolu}

\begin{exercice}[À toi de jouer]
Classe les structures suivantes selon le rapport exprimé (supériorité, infériorité, égalité, superlatif), puis traduis-les en français :
\begin{enumerate}
\item \esp{Este puesto es mejor pagado que el otro.}
\item \esp{Mi hermano tiene tantos diplomas como yo.}
\item \esp{Es la empresa más importante de Douala.}
\item \esp{El tren es menos rápido que el avión.}
\item \esp{Su currículum es interesantísimo.}
\end{enumerate}
\end{exercice}

\begin{corrige}[Exercice : À toi de jouer]
\begin{enumerate}
\item Supériorité (irrégulier \esp{mejor}) : « Ce poste est mieux payé que l'autre. »
\item Égalité sur un nom (\esp{tantos… como}) : « Mon frère a autant de diplômes que moi. »
\item Superlatif relatif : « C'est l'entreprise la plus importante de Douala. »
\item Infériorité : « Le train est moins rapide que l'avion. »
\item Superlatif absolu : « Son CV est extrêmement intéressant. »
\end{enumerate}
\end{corrige}

\begin{savaistu}[Le monde du travail hispanophone et le Cameroun]
L'espagnol est une langue de travail stratégique pour les jeunes Camerounais : la Guinée équatoriale, voisine immédiate, a l'espagnol pour langue officielle, et de nombreuses entreprises espagnoles et latino-américaines recrutent à Douala et à Yaoundé dans le cacao, le bois, la logistique portuaire et le tourisme. Dans tout le monde hispanophone, le \esp{currículum vítae} (souvent abrégé \esp{CV}) suit des rubriques codifiées : \esp{datos personales}, \esp{formación académica}, \esp{experiencia laboral}, \esp{idiomas}, \esp{otros datos de interés}. Nous les étudierons dans la cinquième partie de cette leçon.
\end{savaistu}

\begin{aretenir}[Ce qu'il faut retenir de cette observation]
\begin{itemize}
\item Trois rapports de comparaison : supériorité \esp{más… que}, infériorité \esp{menos… que}, égalité \esp{tan… como} / \esp{tanto… como}.
\item Deux superlatifs : relatif (\esp{el/la más… de}) et absolu (\esp{-ísimo}, \esp{muy}).
\item Des formes irrégulières à mémoriser : \esp{mejor}, \esp{peor}, \esp{mayor}, \esp{menor}.
\item La comparaison peut porter sur un adjectif, un adverbe, un nom ou un verbe : repère toujours l'élément comparé avant de traduire.
\end{itemize}
\end{aretenir}

\coursec{Les comparatifs réguliers et irréguliers}

Après avoir observé dans la section précédente comment la comparaison structure le discours journalistique et professionnel (offres d'emploi, rapports économiques), nous allons maintenant fixer les formes grammaticales. La maîtrise de ces structures est indispensable pour \textbf{traduire} fidèlement du français vers l'espagnol (thème) et pour \textbf{réussir votre \emph{currículum vítae}} en valorisant vos atouts par rapport aux autres candidats.

\courssub{Comparatif de supériorité : \esp{más ... que}}

En espagnol, comme en français, la supériorité s'exprime à l'aide d'un mot-outil invariable placé devant l'élément comparé (adjectif, adverbe ou nom) et suivi de \esp{que}.

\begin{propriete}[Comparatif de supériorité]
\textbf{Forme :} \esp{más + \{adjectif | adverbe | nom\} + que} \\
\textbf{Emploi :} Indiquer qu'un terme possède une qualité, une quantité ou une intensité supérieure à l'autre terme de comparaison. \\
\textbf{Accord :} \esp{más} est invariable. L'adjectif s'accorde en genre et en nombre avec le nom qu'il qualifie. Le nom suit directement \esp{más}.
\end{propriete}

\begin{exemplebox}[Exemples variés de supériorité]
\esp{Este puesto está \textbf{más \emph{bien pagado}} que el anterior.} « Ce poste est mieux payé que le précédent. » (adverbe composé) \\
\esp{La empresa ofrece \textbf{más \emph{vacaciones}} que la competencia.} « L'entreprise offre plus de vacances que la concurrence. » (nom) \\
\esp{María habla inglés \textbf{más \emph{fluidamente}} que Juan.} « María parle anglais plus couramment que Juan. » (adverbe en \emph{-mente}) \\
\esp{Mi salario es \textbf{más \emph{alto}} que el tuyo.} « Mon salaire est plus élevé que le tien. » (adjectif)
\end{exemplebox}

\courssub{Comparatif d'infériorité : \esp{menos ... que}}

La structure est symétrique à la supériorité. \esp{Menos} est également invariable.

\begin{propriete}[Comparatif d'infériorité]
\textbf{Forme :} \esp{menos + \{adjectif | adverbe | nom\} + que} \\
\textbf{Emploi :} Indiquer une qualité, quantité ou intensité moindre. \\
\textbf{Remarque :} En français, « moins de » devant un nom se traduit par \esp{menos} + nom (sans \esp{de}). Ex. : \esp{Tengo menos tiempo que tú} (J'ai moins de temps que toi).
\end{propriete}

\begin{exemplebox}[Exemples d'infériorité]
\esp{Este contrato tiene \textbf{menos \emph{cláusulas}} que el estándar.} « Ce contrat a moins de clauses que le standard. » \\
\esp{El candidato B tiene \textbf{menos \emph{experiencia}} que la candidata A.} « Le candidat B a moins d'expérience que la candidate A. » \\
\esp{Trabaja \textbf{menos \emph{eficazmente}} que su compañero.} « Il travaille moins efficacement que son collègue. »
\end{exemplebox}

\courssub{Comparatif d'égalité : \esp{tan ... como} et \esp{tanto ... como}}

Ici, la langue espagnole distingue selon la nature grammaticale de l'élément comparé. C'est une source d'erreurs majeure pour les francophones.

\begin{propriete}[Comparatif d'égalité -- Trois structures]
\begin{enumerate}
\item \textbf{Adjectif ou adverbe :} \esp{tan + \{adjectif | adverbe\} + como} \\
      \esp{tan rápido como} (aussi vite que), \esp{tan competente como} (aussi compétent que).
\item \textbf{Nom :} \esp{tanto/a/os/as + nom + como} \\
      \emph{Tanto} s'accorde en genre et en nombre avec le nom qui suit. \\
      \esp{tanta \emph{paciencia} como}, \esp{tantos \emph{empleados} como}, \esp{tantas \emph{horas} como}.
\item \textbf{Verbe (quantité d'action) :} \esp{verbe + tanto como} \\
      \esp{Trabaja tanto como tú.} (Il travaille autant que toi.)
\end{enumerate}
\end{propriete}

\begin{center}
\begin{tabularx}{\linewidth}{|X|X|X|}
\hline
\rowcolor{popblueL} \textbf{Français} & \textbf{Español} & \textbf{Structure} \\
\hline
Aussi compétent que & \esp{tan competente como} & \esp{tan + adj. + como} \\
Aussi vite que & \esp{tan rápido como} & \esp{tan + adv. + como} \\
Autant de patience que & \esp{\textbf{tanta} paciencia como} & \esp{tanto (accordé) + nom + como} \\
Autant de candidats que & \esp{\textbf{tántos} candidatos como} & \esp{tanto (accordé) + nom + como} \\
Autant d'expérience que & \esp{\textbf{tanta} experiencia como} & \esp{tanto (accordé) + nom + como} \\
Il travaille autant que moi & \esp{Trabaja tanto como yo} & \esp{verbe + tanto como} \\
\hline
\end{tabularx}
\end{center}

\begin{attention}[Piège n°1 : L'accord obligatoire de \esp{tanto} devant un nom]
En français, « autant de » est invariable. En espagnol, \esp{tanto} est un adjectif indéfini qui \textbf{s'accorde en genre et en nombre} avec le nom qu'il détermine.
\begin{itemize}
\item \esp{Tengo \textbf{tanta} formación como tú.} (féminin singulier)
\item \esp{Recibí \textbf{tántos} currículos como ayer.} (masculin pluriel)
\item \esp{Ofrecen \textbf{tantas} ventajas como la otra empresa.} (féminin pluriel)
\end{itemize}
\textbf{Erreur fréquente :} \esp{*tanto experiencia como} $\rightarrow$ \textbf{Correct :} \esp{\textbf{tanta} experiencia como}.
\end{attention}

\begin{exemplebox}[Égalité en contexte professionnel]
\esp{El nuevo software es \textbf{tan intuitivo como} el anterior.} « Le nouveau logiciel est aussi intuitif que l'ancien. » \\
\esp{La becaria muestra \textbf{tanta} iniciativa \textbf{como} un empleado fijo.} « La stagiaire fait preuve d'autant d'initiative qu'un employé titulaire. » \\
\esp{En la entrevista, \textbf{respondió tanto como} pudo.} « Lors de l'entretien, il a répondu autant qu'il a pu. »
\end{exemplebox}

\courssub{Les comparatifs irréguliers : \esp{mejor, peor, mayor, menor}}

Quatre adjectifs très fréquents possèdent des formes comparatives synthétiques (un seul mot) qui remplacent \esp{más bueno}, \esp{más malo}, \esp{más grande}, \esp{más pequeño}. Leur utilisation est \textbf{obligatoire} dans la langue standard.

\begin{propriete}[Comparatifs irréguliers -- Formes et sens]
\begin{center}
\begin{tabular}{|l|l|l|l|}
\hline
\rowcolor{popblueL} \textbf{Positif} & \textbf{Comparatif irrégulier} & \textbf{Traduction} & \textbf{Exemple} \\
\hline
\esp{bueno} (bon) & \esp{mejor} & meilleur & \esp{Es \textbf{mejor} candidato.} \\
\esp{malo} (mauvais) & \esp{peor} & pire & \esp{Es el \textbf{peor} escenario.} \\
\esp{grande} (grand) & \esp{mayor} & plus grand / plus âgé & \esp{Mi hermano es \textbf{mayor} que yo.} \\
\esp{pequeño} (petit) & \esp{menor} & plus petit / plus jeune & \esp{Es el \textbf{menor} de la familia.} \\
\hline
\end{tabular}
\end{center}
\end{propriete}

\begin{attention}[Piège n°2 : Interdiction du double comparatif \esp{*más mejor} / \esp{*más malo}]
\begin{itemize}
\item \textbf{Jamais} \esp{*más bueno} $\rightarrow$ Toujours \esp{mejor}.
\item \textbf{Jamais} \esp{*más malo} $\rightarrow$ Toujours \esp{peor}.
\item \textbf{Jamais} \esp{*más grande} (pour l'âge/hiérarchie) $\rightarrow$ \esp{mayor}.
\item \textbf{Jamais} \esp{*más pequeño} (pour l'âge/hiérarchie) $\rightarrow$ \esp{menor}.
\item \textbf{Jamais} \esp{*más mejor}, \esp{*más peor} (double marquage, faute grave).
\end{itemize}
\end{attention}

\begin{savaistu}[Le double comparatif : une faute « populaire »]
En français familier, on entend parfois « \emph{plus meilleur} » ou « \emph{plus pire} ». En espagnol, \esp{más mejor} et \esp{más peor} sont considérés comme des \emph{barbarismes} (fautes graves de morphologie). Les formes \esp{mejor} et \esp{peor} contiennent déjà l'idée de « plus ». De même, \esp{mayor} et \esp{menor} intègrent la comparaison.
\end{savaistu}

\begin{propriete}[Emplois restreints de \esp{mayor} et \esp{menor}]
\esp{Mayor} et \esp{menor} s'emploient principalement pour :
\begin{enumerate}
\item \textbf{L'âge :} \esp{Mi hermana es \textbf{mayor} que yo.} (Ma sœur est plus âgée que moi.)
\item \textbf{La taille/importance abstraite :} \esp{El \textbf{mayor} problema es la falta de fondos.} (Le plus grand problème est le manque de fonds.) ; \esp{Una \textbf{menor} cuantía.} (Une moindre somme.)
\end{enumerate}
Pour la \textbf{taille physique} (hauteur, volume), on utilise les réguliers :
\esp{Este edificio es \textbf{más grande} que aquél.} (Cet immeuble est plus grand que celui-là.)
\esp{Esta caja es \textbf{más pequeña} que la otra.} (Cette boîte est plus petite que l'autre.)
\end{propriete}

\begin{exemplebox}[Irréguliers en contexte CV / Entretien]
\esp{Mi \textbf{mejor} cualidad es la puntualidad.} « Ma meilleure qualité est la ponctualité. » \\
\esp{Este ha sido mi \textbf{peor} año laboral.} « Ce fut ma pire année professionnelle. » \\
\esp{Buscamos al candidato de \textbf{mayor} experiencia.} « Nous cherchons le candidat ayant la plus grande expérience. » \\
\esp{Fui el \textbf{menor} de mis hermanos en incorporarme al mercado laboral.} « Je fus le plus jeune de mes frères à entrer sur le marché du travail. »
\end{exemplebox}

\courssub{Cas particulier : Devant un chiffre -- \esp{más de / menos de}}

Lorsque la comparaison porte sur un \textbf{nombre, une quantité chiffrée, une date ou un pourcentage}, la préposition change : \esp{que} devient \esp{de}.

\begin{propriete}[Comparatif de quantité chiffrée]
\textbf{Règle :} Devant un chiffre (nombre, date, \%), \esp{más} et \esp{menos} sont suivis de \esp{de} (et non \esp{que}).
\begin{itemize}
\item \esp{Más de cien candidatos.} (Plus de cent candidats.)
\item \esp{Menos de 20 años.} (Moins de 20 ans.)
\item \esp{Más del 50\% de los votos.} (Plus de 50\% des voix.)
\end{itemize}
\textbf{Contraste :}
\esp{Tengo \textbf{más de} diez años de experiencia.} (Quantité chiffrée : \esp{de})
\esp{Tengo \textbf{más} experiencia \textbf{que} tú.} (Comparaison standard : \esp{que})
\end{propriete}

\begin{attention}[Piège n°3 : \esp{más de} vs \esp{más que}]
C'est l'erreur n°1 en thème (français $\rightarrow$ espagnol) sur les comparatifs.
\begin{itemize}
\item « Plus de 10 ans » $\rightarrow$ \esp{\textbf{más de} diez años}.
\item « Plus que toi » $\rightarrow$ \esp{\textbf{más que} tú}.
\item « Plus d'expérience que toi » $\rightarrow$ \esp{\textbf{más} experiencia \textbf{que} tú} (nom non chiffré).
\end{itemize}
\end{attention}

\begin{exoresolu}[Traduction guidée : Comparatifs en contexte professionnel]
\textbf{Énoncé :} Traduisez en espagnol les phrases suivantes.
\begin{enumerate}
\item Ce poste est mieux rémunéré que le précédent.
\item Elle a moins d'expérience que moi.
\item Il parle anglais aussi couramment que son frère.
\item Nous avons reçu autant de CV que la semaine dernière.
\item C'est le meilleur candidat pour ce poste.
\item Ma sœur est plus jeune que moi.
\item L'entreprise a plus de 500 employés.
\item Ce problème est moins grave qu'il n'y paraît.
\end{enumerate}

\tcblower
\textbf{Solution détaillée :}
\begin{enumerate}
\item \esp{Este puesto está \textbf{mejor remunerado} que el anterior.} (\emph{Irrégulier : mejor} ; \emph{participe passé comme adjectif}).
\item \esp{Ella tiene \textbf{menos experiencia que} yo.} (\emph{Infériorité sur nom : menos ... que} ; \emph{pas de "de"}).
\item \esp{Él habla inglés \textbf{tan fluidamente como} su hermano.} (\emph{Égalité adverbe : tan ... como}).
\item \esp{Hemos recibido \textbf{tántos currículos como} la semana pasada.} (\emph{Égalité nom pluriel masc. : tantos ... como} ; \emph{accent sur tántos}).
\item \esp{Es el \textbf{mejor} candidato para este puesto.} (\emph{Irrégulier mejor devant nom} ; \emph{pas de "más"}).
\item \esp{Mi hermana es \textbf{menor} que yo.} (\emph{Irrégulier pour l'âge : menor} ; \emph{pas "más pequeña"}).
\item \esp{La empresa tiene \textbf{más de} quinientos empleados.} (\emph{Chiffre : más de}).
\item \esp{Este problema es \textbf{menos grave} de lo que parece.} (\emph{Infériorité adjectif : menos ... que} ; \emph{structure "de lo que" après adjectif pour clause}). \emph{Note :} \esp{menos grave que lo que parece} est aussi possible, mais \esp{de lo que} est plus élégant après comparatif.
\end{enumerate}
\end{exoresolu}

\begin{textebox}[Extrait d'un rapport RH : \emph{Análisis comparativo de perfiles}]
\textbf{Contexto :} Un responsable des ressources humaines d'une multinationale à Mexico compare deux profils pour un poste de chef de projet.

\emph{Tras revisar los dos currículos, el perfil de Ana presenta \textbf{mayor} coherencia con los requisitos del puesto. Aunque Carlos tiene \textbf{tanta} formación técnica \textbf{como} ella, Ana demuestra \textbf{mejor} capacidad de liderazgo: ha dirigido equipos \textbf{más grandes que} los de Carlos y con \textbf{mejores} resultados en plazos de entrega. En cuanto a idiomas, Carlos habla inglés \textbf{tan fluidamente como} Ana, pero ella posee \textbf{menor} dificultad para adaptarse a entornos multiculturales, pues ha vivido en tres países. El salario que solicita Ana es \textbf{más alto que} el de Carlos, pero está \textbf{menos} alejado de la media del sector de lo que parece; de hecho, pide \textbf{menos de} 5.000 pesos por encima del convenio. En conclusión, Ana es la candidata \textbf{mejor} valorada por el comité.}

\medskip
\textbf{Glossaire :}
\begin{itemize}
\item \esp{coherencia} : cohérence, adéquation
\item \esp{plazos de entrega} : délais de livraison
\item \esp{entornos multiculturales} : environnements multiculturels
\item \esp{media del sector} : moyenne du secteur
\item \esp{convenio} : convention collective
\item \esp{valorada} : évaluée, appréciée
\end{itemize}

\textbf{Questions d'observation (oral / écrit) :}
\begin{enumerate}
\item Relevez tous les comparatifs du texte et classez-les (supériorité, infériorité, égalité, irréguliers).
\item Pourquoi \esp{mayor} et non \esp{más grande} pour \esp{coherencia} ? Pourquoi \esp{menor} pour \esp{dificultad} ?
\item Justifiez l'accord de \esp{tanta} et l'absence d'accord de \esp{tan}.
\item Traduisez la phrase : \esp{El salario que solicita Ana es más alto que el de Carlos...}
\end{enumerate}
\end{textebox}

\coursec{Les superlatifs : relatif et absolu}

Dans la section précédente, nous avons appris à \emph{comparer} deux éléments : \esp{María es más alta que Ana}, \esp{Pedro trabaja tanto como su hermano}. Mais que se passe-t-il quand un élément se distingue de \emph{tous} les autres, quand il occupe la première (ou la dernière) place d'un groupe ? Et comment exprimer qu'une qualité est portée à son \emph{maximum}, sans comparaison avec personne ? C'est le rôle du \cle{superlatif}, qui existe sous deux formes en espagnol : le \cle{superlatif relatif} (\esp{el más alto de la clase}) et le \cle{superlatif absolu} (\esp{altísimo}). Observons d'abord ces deux mécanismes dans un texte authentique du monde du travail.

\courssub{Observation : le superlatif dans le monde du travail}

\begin{textebox}[Anuncio de prensa : la empresa del año]
La empresa camerunesa \emph{AgroCacao S.A.} ha sido elegida \textbf{la empresa más innovadora del país} por un jurado de expertos internacionales. Su directora general, la señora Mbarga, es \textbf{la más joven de los directivos} del sector, pero también \textbf{la mejor preparada}: habla cuatro idiomas y posee un currículum \textbf{buenísimo}, con experiencia en España y en Guinea Ecuatorial.

«Nuestro equipo es \textbf{sumamente competente}», declara la directora. «Tenemos los ingenieros \textbf{más cualificados de la región} y los obreros \textbf{más experimentados}. El trabajo aquí es \textbf{dificilísimo} a veces, pero el ambiente es \textbf{agradabilísimo}. Somos la fábrica \textbf{menos contaminante del continente} y queremos seguir siendo \textbf{la mejor}.»

Los salarios no son los \textbf{más altos del mercado}, pero los empleados afirman que la satisfacción es \textbf{altísima}: el noventa por ciento se declara \textbf{felicísimo} de trabajar allí. Para los jóvenes titulados de Yaoundé y Douala, conseguir un puesto en AgroCacao se ha convertido en el objetivo \textbf{más importante de su carrera}.
\end{textebox}

\textbf{Glossaire :} \esp{el jurado} = le jury ; \esp{los directivos} = les cadres dirigeants ; \esp{cualificado} = qualifié ; \esp{contaminante} = polluant ; \esp{los salarios} = les salaires ; \esp{un puesto} = un poste ; \esp{conseguir} = obtenir ; \esp{seguir siendo} = continuer à être ; \esp{la carrera} = la carrière professionnelle.

\textbf{Questions d'observation :}
\begin{enumerate}
\item Relevez toutes les formes superlatives du texte et classez-les en deux groupes.
\item Quelle préposition introduit le complément après \esp{más innovadora} ? Et en français ?
\item Quel suffixe revient dans \esp{buenísimo}, \esp{dificilísimo}, \esp{felicísimo} ? Que signifie-t-il ?
\end{enumerate}

\textbf{Observations.} On remarque deux constructions différentes :
\begin{itemize}
\item \esp{la empresa más innovadora \textbf{del} país}, \esp{la más joven \textbf{de} los directivos} : l'élément est comparé à un \emph{groupe} dont il fait partie → \textbf{superlatif relatif}.
\item \esp{buenísimo}, \esp{dificilísimo}, \esp{sumamente competente} : aucune comparaison, la qualité est simplement portée à un \emph{très haut degré} → \textbf{superlatif absolu}.
\end{itemize}

\courssub{Le superlatif relatif : \esp{el más... de}}

\begin{propriete}[Formation du superlatif relatif]
Le superlatif relatif exprime qu'un élément possède une qualité \emph{au maximum} par rapport à un groupe. Sa structure est :

\begin{center}
\textbf{article défini + (nom) + más / menos + adjectif + de + complément}
\end{center}

\esp{El candidato más serio de la empresa.} « Le candidat le plus sérieux de l'entreprise. »\\
\esp{La fábrica menos contaminante del continente.} « L'usine la moins polluante du continent. »

L'article s'\textbf{accorde} avec le nom : \esp{el más alto}, \esp{la más alta}, \esp{los más altos}, \esp{las más altas}. L'adjectif s'accorde aussi normalement. Quand le nom est sous-entendu, on garde l'article seul : \esp{Esta candidata es la más seria.} « Cette candidate est la plus sérieuse. »
\end{propriete}

\begin{attention}[\esp{de} et non \esp{que} !]
C'est l'erreur numéro un des francophones. En français, « le plus grand \emph{que}... » n'existe pas au superlatif : on dit « le plus grand \emph{de} la classe ». En espagnol, le complément du superlatif est TOUJOURS introduit par \esp{de}, jamais par \esp{que} :

\esp{Es la mejor estudiante \textbf{del} grupo.} « C'est la meilleure élève du groupe. » (et non \esp{que el grupo})

Attention à la contraction obligatoire : \esp{de + el = del}. « Le meilleur du monde » = \esp{el mejor \textbf{del} mundo}. En revanche, pas de contraction avec \esp{la} : \esp{la más alta \textbf{de la} clase}.
\end{attention}

\begin{propriete}[Superlatifs relatifs irréguliers]
Comme au comparatif, quatre adjectifs ont des formes spéciales :

\begin{center}
\begin{tabular}{|l|l|l|}
\hline
\rowcolor{popblueL} Adjectif & Superlatif relatif & Exemple \\
\hline
bueno (bon) & el mejor & \esp{el mejor candidato} \\
\hline
malo (mauvais) & el peor & \esp{el peor error} \\
\hline
grande (grand, âgé) & el mayor & \esp{el mayor de los hermanos} \\
\hline
pequeño (petit, jeune) & el menor & \esp{la menor de la clase} \\
\hline
\end{tabular}
\end{center}

\esp{Es el mejor del país.} « C'est le meilleur du pays. »\\
\esp{Fue el peor día de mi vida.} « Ce fut le pire jour de ma vie. »

\textbf{Remarque :} \esp{mayor} et \esp{menor} s'emploient surtout pour l'\emph{âge} (aîné / cadet) ; pour la taille ou la dimension, on dit \esp{el más grande}, \esp{el más pequeño} : \esp{la ciudad más grande del país} « la plus grande ville du pays ».
\end{propriete}

\begin{exemplebox}[Exemples traduits]
\begin{itemize}
\item \esp{Es la empresa más importante del país.} « C'est l'entreprise la plus importante du pays. »
\item \esp{Yaoundé es una de las ciudades más grandes de Camerún.} « Yaoundé est l'une des plus grandes villes du Cameroun. »
\item \esp{Este es el puesto menos pagado de la empresa.} « C'est le poste le moins payé de l'entreprise. »
\item \esp{Mi hermana es la mayor de la familia.} « Ma sœur est l'aînée de la famille. »
\item \esp{El español es la asignatura más fácil de mi horario.} « L'espagnol est la matière la plus facile de mon emploi du temps. »
\end{itemize}
\end{exemplebox}

\courssub{Le superlatif absolu : \esp{-ísimo}, \esp{muy}, \esp{sumamente}}

Le superlatif absolu ne compare \emph{rien} : il affirme simplement qu'une qualité atteint un degré \emph{très élevé}. Il correspond au français « très », « extrêmement », « tout à fait ».

\begin{propriete}[Formation du superlatif absolu en \esp{-ísimo}]
On ajoute le suffixe \esp{-ísimo / -ísima / -ísimos / -ísimas} à l'adjectif, qui s'accorde en genre et en nombre avec le nom.

\textbf{Règle de formation :}
\begin{itemize}
\item Adjectif terminé par une voyelle : on \textbf{supprime la voyelle finale} puis on ajoute \esp{-ísimo} : \esp{alto → altísimo}, \esp{grande → grandísimo}, \esp{bonita → bonitísima}.
\item Adjectif terminé par une consonne : on ajoute directement \esp{-ísimo} : \esp{fácil → facilísimo}, \esp{feliz → felicísimo}, \esp{joven → jovencísimo}.
\item Le suffixe porte TOUJOURS un \textbf{accent} sur le \esp{í}, et l'accent éventuel de l'adjectif simple disparaît : \esp{fácil → facilísimo}, \esp{difícil → dificilísimo}.
\end{itemize}

\esp{El examen fue facilísimo.} « L'examen était très facile / extrêmement facile. »\\
\esp{Su currículum es buenísimo.} « Son CV est excellent (très bon). »
\end{propriete}

\begin{attention}[Modifications orthographiques devant \esp{-ísimo}]
Pour conserver le son de la consonne finale, certains adjectifs changent d'orthographe :
\begin{itemize}
\item \textbf{-co → -quísimo} : \esp{rico → riquísimo}, \esp{poco → poquísimo}, \esp{chico → chiquísimo}.
\item \textbf{-go → -guísimo} : \esp{largo → larguísimo}, \esp{amargo → amarguísimo} (le \esp{u} garde le son [g]).
\item \textbf{-z → -císimo} : \esp{feliz → felicísimo}, \esp{feroz → ferocísimo}.
\item Quelques formes héritées du latin (registre soutenu ou culturel) : \esp{pobre → paupérrimo}, \esp{antiguo → antiquísimo}, \esp{fuerte → fortísimo}.
\end{itemize}
\end{attention}

\begin{propriete}[Les alternatives : \esp{muy} et \esp{sumamente}]
À l'écrit soutenu, et surtout dans les CV et lettres de motivation, on préfère souvent :
\begin{itemize}
\item \esp{muy + adjectif} : \esp{muy cualificado} « très qualifié » ;
\item \esp{sumamente + adjectif} : \esp{sumamente motivado} « extrêmement motivé » (registre professionnel, plus élégant).
\end{itemize}
\esp{Es una candidata sumamente competente.} « C'est une candidate extrêmement compétente. »
\end{propriete}

\begin{attention}[« Très meilleur » n'existe pas]
\esp{muy} ne peut jamais modifier un comparatif ou un superlatif irrégulier. Pour dire « bien meilleur / beaucoup meilleur », on utilise \esp{mucho} : \esp{Este método es \textbf{mucho mejor}.} « Cette méthode est bien meilleure. » De même : \esp{mucho peor} « bien pire », \esp{mucho mayor} « bien plus âgé ». Jamais \esp{muy mejor} !
\end{attention}

\begin{savaistu}[\esp{sumamente}, le mot qui fait professionnel]
Dans les lettres de motivation et les CV hispanophones, \esp{sumamente} est très apprécié : écrire \esp{Estoy sumamente motivado por este puesto} fait beaucoup plus professionnel que \esp{muy motivado}. À l'inverse, le suffixe \esp{-ísimo} est surtout oral et expressif : un Équato-Guinéen dira volontiers \esp{¡La fiesta estuvo buenísima!}, mais dans un document administratif il écrira \esp{muy satisfactorio}. Retenez la nuance : \esp{-ísimo} = expressif et chaleureux ; \esp{sumamente} = soutenu et professionnel.
\end{savaistu}

\courssub{Tableau récapitulatif et comparaison avec le français}

\begin{center}
\begin{tabularx}{\linewidth}{|X|X|X|}
\hline
\rowcolor{popblueL} Français & Español & Remarque \\
\hline
le plus grand de la classe & el más alto de la clase & \esp{de}, jamais \esp{que} \\
\hline
le meilleur du groupe & el mejor del grupo & contraction \esp{de + el = del} \\
\hline
très grand & muy alto / altísimo & « très » ne se suffixe jamais en français \\
\hline
extrêmement qualifié & sumamente cualificado & registre soutenu, CV \\
\hline
très facile & facilísimo & accent obligatoire sur \esp{í} \\
\hline
très riche & riquísimo & \esp{-co → -quísimo} \\
\hline
bien meilleur & mucho mejor & jamais \esp{muy mejor} \\
\hline
\end{tabularx}
\end{center}

\begin{attention}[Piège de traduction]
Le français « très » est un adverbe \emph{invariable} qui se place devant l'adjectif ; l'espagnol \esp{-ísimo} est un \emph{suffixe accordé}. Ne traduisez jamais « très bonne » par \esp{muy buena} si le contexte est expressif et oral : préférez \esp{buenísima}. Et n'oubliez pas l'accord : \esp{las clases son interesantísimas} (féminin pluriel).
\end{attention}

\begin{popfigure}
\begin{tikzpicture}[font=\small]
\draw[-{Stealth}, very thick, popblue] (0,0) -- (10.5,0);
\foreach \x/\txt in {0.8/{alto}, 3.6/{muy alto}, 6.4/{altísimo}, 9.4/{el más alto\\ de la clase}} {
  \fill[poporange] (\x,0) circle (0.09);
  \node[above, align=center, popdark] at (\x,0.15) {\txt};
}
\node[below, popmuted] at (5.2,-0.35) {intensité croissante};
\draw[popgreen, thick] (2.6,-0.9) -- (7.4,-0.9);
\draw[popgreen, thick] (2.6,-0.9) -- (2.6,-0.7);
\draw[popgreen, thick] (7.4,-0.9) -- (7.4,-0.7);
\node[below, popgreen] at (5,-1.15) {superlatif absolu (sans comparaison)};
\node[below, poppurple] at (9.4,-1.15) {superlatif relatif};
\end{tikzpicture}
\legende{Échelle de l'intensité : de l'adjectif simple au superlatif relatif.}
\end{popfigure}

\begin{popfigure}
\begin{tikzpicture}[mindmap, grow cyclic, every node/.style={concept, font=\small},
  root concept/.append style={concept color=popblue, text=white, font=\small\bfseries},
  level 1/.append style={level distance=3.4cm, sibling angle=90},
  level 2/.append style={level distance=2.6cm, sibling angle=50}]
\node[root concept]{SUPERLATIVO}
  child[concept color=poporange] { node {relativo}
    child[concept color=poporangeL] { node[align=center] {el más alto\\ de la clase} }
    child[concept color=poporangeL] { node[align=center] {el mejor\\ del grupo} } }
  child[concept color=popgreen] { node {absoluto}
    child[concept color=popgreenL] { node[align=center] {altísimo\\ buenísimo} }
    child[concept color=popgreenL] { node[align=center] {muy alto} }
    child[concept color=popgreenL] { node[align=center] {sumamente\\ cualificado} } };
\end{tikzpicture}
\legende{Carte mentale : les deux visages du superlatif espagnol.}
\end{popfigure}

\begin{exoresolu}[Traduction guidée (thème)]
Traduisez en espagnol : « María est la candidate la plus qualifiée de l'entreprise. Son dossier est très bon et elle est extrêmement motivée. C'est aussi la plus jeune du groupe. »

\tcblower
\textbf{Solution détaillée.}
\begin{enumerate}
\item « la candidate la plus qualifiée de l'entreprise » → superlatif relatif : article féminin \esp{la} + \esp{más} + adjectif accordé + \esp{de} : \esp{la candidata más cualificada de la empresa}.
\item « son dossier est très bon » → « dossier » se dit \esp{expediente} (attention au faux ami : \esp{carpeta} = chemise, classeur) ; superlatif absolu expressif : \esp{Su expediente es buenísimo} (on supprime la voyelle de \esp{bueno}).
\item « extrêmement motivée » → registre soutenu : \esp{sumamente motivada} (accord féminin).
\item « la plus jeune du groupe » → \esp{la más joven del grupo} (contraction \esp{de + el = del}).
\end{enumerate}
\textbf{Traduction complète :} \esp{María es la candidata más cualificada de la empresa. Su expediente es buenísimo y está sumamente motivada. Es también la más joven del grupo.}
\end{exoresolu}

\begin{exercice}[À vous de jouer]
Traduisez en espagnol :
\begin{enumerate}
\item Douala est la ville la plus peuplée du Cameroun.
\item Cet exercice est très facile, mais le suivant est très difficile.
\item C'est le pire restaurant de la ville, et les prix sont très élevés.
\item Mon frère aîné est le plus travailleur de la famille.
\end{enumerate}
\end{exercice}

\begin{corrige}[Exercice : superlatifs]
\begin{enumerate}
\item \esp{Douala es la ciudad más poblada de Camerún.} (superlatif relatif + \esp{de})
\item \esp{Este ejercicio es facilísimo, pero el siguiente es dificilísimo.} (accents obligatoires ; \esp{fácil} et \esp{difícil} perdent leur accent en gagnant le suffixe)
\item \esp{Es el peor restaurante de la ciudad, y los precios son altísimos.} (\esp{peor} irrégulier ; accord pluriel \esp{altísimos})
\item \esp{Mi hermano mayor es el más trabajador de la familia.} (\esp{mayor} = aîné)
\end{enumerate}
\end{corrige}

\begin{aretenir}[L'essentiel de la section]
\begin{itemize}
\item \textbf{Superlatif relatif} : \esp{el/la/los/las + más (menos) + adjectif + de} ; formes irrégulières \esp{el mejor, el peor, el mayor, el menor} ; complément introduit par \esp{de} (contraction \esp{del}).
\item \textbf{Superlatif absolu} : \esp{adjectif + -ísimo} (accordé, accentué) ou \esp{muy / sumamente + adjectif}.
\item Modifications orthographiques : \esp{rico → riquísimo}, \esp{largo → larguísimo}, \esp{feliz → felicísimo}, \esp{grande → grandísimo}.
\item « Bien meilleur » = \esp{mucho mejor}, jamais \esp{muy mejor}.
\item Registres : \esp{-ísimo} = expressif, oral ; \esp{sumamente} = soutenu, idéal pour le CV que nous rédigerons dans la section 5.
\end{itemize}
\end{aretenir}

\coursec{Traduction guidée : thème et version}

Dans la section précédente, nous avons étudié les superlatifs, relatif (\esp{el más alto de la clase}) et absolu (\esp{altísimo}, \esp{muy alto}, \esp{sumamente alto}). Nous savons donc désormais \emph{former} comparatifs et superlatifs. Mais un savoir grammatical ne suffit pas à l'examen : il faut savoir \emph{traduire}, c'est-à-dire passer du français à l'espagnol (thème) et de l'espagnol au français (version) sans faute de structure ni d'accord. Cette section vous donne une méthode systématique, puis vous entraîne sur une version et un thème complets, corrigés et justifiés phrase par phrase.

\courssub{Méthode : traduire une comparaison en cinq étapes}

Traduire une comparaison n'est pas une affaire d'intuition : c'est une \emph{décision} que l'on peut décomposer en étapes mécaniques. Suivez toujours le même protocole.

\begin{methode}[Traduire une comparaison : les cinq étapes]
\begin{enumerate}
\item \textbf{Souligner} dans la phrase française l'élément comparé : est-ce un adjectif (\emph{plus grand}), un adverbe (\emph{moins vite}), un nom (\emph{autant de travail}) ou un verbe (\emph{il travaille plus}) ?
\item \textbf{Choisir la structure} en fonction de la nature de cet élément : \esp{más/menos + adjectif ou adverbe + que} ; \esp{tan + adjectif ou adverbe + como} ; \esp{tanto/a/os/as + nom + como} ; \esp{verbe + más/menos + que}.
\item \textbf{Accorder} : \esp{tanto} s'accorde avec le nom qui suit (\esp{tanta paciencia}, \esp{tantos empleos}) ; l'adjectif s'accorde avec son nom (\esp{una crisis más grave}).
\item \textbf{Vérifier la particule} : \esp{que} après \esp{más/menos} ; \esp{como} après \esp{tan/tanto} ; \esp{de} (et non \esp{que}) devant un chiffre : \esp{más de veinte años}.
\item \textbf{Relire la phrase entière} : verbe conjugué, sujet placé, accents écrits, sens global respecté.
\end{enumerate}
\end{methode}

\begin{popfigure}
\begin{center}
\begin{tikzpicture}[
  node distance=0.55cm and 0.9cm,
  decision/.style={diamond, draw=poporange, fill=poporangeL, text width=3.4cm, align=center, inner sep=1pt, font=\small, aspect=2.2},
  action/.style={rectangle, rounded corners, draw=popblue, fill=popblueL, text width=4.6cm, align=center, font=\small, inner sep=4pt},
  start/.style={rectangle, rounded corners, draw=popgreen, fill=popgreenL, text width=5.5cm, align=center, font=\small\bfseries, inner sep=4pt},
  fleche/.style={-{Stealth[length=2.5mm]}, thick, popdark}
]
\node[start] (dep) {Quel est l'élément comparé ?};
\node[decision, below left=0.8cm and -1.2cm of dep] (adj) {Adjectif ou adverbe ?};
\node[decision, below right=0.8cm and -1.2cm of dep] (nom) {Un nom ?};
\node[action, below=0.8cm of adj, align=center] (resadj){\esp{más/menos... que}\\ \esp{tan... como}\\ (invariable)};
\node[action, below=0.8cm of nom, align=center] (resnom){\esp{más/menos + nom + que}\\ \esp{tanto/a/os/as + nom + como}};
\node[decision, below=1.5cm of dep, yshift=-3.4cm] (chiffre) {Un chiffre suit-il ?};
\node[action, below=0.7cm of chiffre, align=center] (reschiffre){\esp{más de / menos de + chiffre}\\ (jamais \esp{que} !)};
\draw[fleche] (dep) -- (adj) node[midway, above left, font=\footnotesize]{oui};
\draw[fleche] (dep) -- (nom) node[midway, above right, font=\footnotesize]{nom};
\draw[fleche] (adj) -- (resadj);
\draw[fleche] (nom) -- (resnom);
\draw[fleche] (resadj.south) |- (chiffre.west);
\draw[fleche] (resnom.south) |- (chiffre.east);
\draw[fleche] (chiffre) -- (reschiffre) node[midway, right, font=\footnotesize]{oui};
\end{tikzpicture}
\end{center}
\legende{Organigramme de décision : choisir la bonne structure comparative.}
\end{popfigure}

\begin{attention}[Le piège capital : « autant de » devant un nom]
En thème, la tentation est grande de traduire mécaniquement « autant... que » par \esp{tan... como}. C'est faux devant un nom : « autant de travail que » se traduit \esp{tanto trabajo como} et non \esp{tan trabajo como}. La règle : \esp{tan} ne précède jamais un nom ; devant un nom, on emploie \esp{tanto} (accordé). De même, « plus de vingt ans » se dit \esp{más de veinte años} : devant un chiffre, \esp{que} devient \esp{de}.
\end{attention}

\courssub{Version : l'emploi des jeunes au Cameroun et en Espagne}

Lisez d'abord le texte en entier sans traduire, pour en saisir le sens global. Repérez ensuite les huit structures à traduire : six comparaisons et deux superlatifs.

\begin{textebox}[El empleo de los jóvenes : dos realidades]
En España, los jóvenes encuentran menos empleos que sus padres a la misma edad. La situación es peor que hace diez años, pero los salarios son más altos que en otros países europeos. El sector más dinámico es el de la tecnología, que ofrece tantos puestos como el comercio.

En Camerún, el problema es diferente. Los jóvenes de Yaoundé y Duala trabajan más horas que los funcionarios, pero ganan menos dinero que ellos. El sector informal es tan importante como el sector público : los vendedores del mercado, los conductores de mototaxis y los artesanos mantienen a tantas familias como las grandes empresas. Sin embargo, la formación profesional es menos desarrollada que la educación general, y muchos diplomados aceptan trabajos peor pagados que sus estudios.

Para los economistas, el desafío más urgente es el de crear empleos estables. « El peor enemigo del joven no es el trabajo duro, sino la falta de oportunidades », escribe un periodista de Duala. Es una verdad tan simple como olvidada.
\end{textebox}

\textbf{Glossaire}

\begin{center}
\begin{tabularx}{\linewidth}{|X|X|X|}\hline
\rowcolor{popblueL} Español & Français & Remarque \\ \hline
a la misma edad & au même âge & \esp{edad} est féminin \\ \hline
los puestos & les postes (emplois) & faux ami : pas « postes » de police \\ \hline
los funcionarios & les fonctionnaires & \\ \hline
el sector informal & le secteur informel & économie non déclarée \\ \hline
los mototaxis & les motos-taxis & les « bendskins » \\ \hline
mantener a & faire vivre, subvenir aux besoins de & irrégulier comme \esp{tener} \\ \hline
los diplomados & les diplômés & \\ \hline
peor pagados & moins bien payés & \esp{peor} = pire / plus mal \\ \hline
el desafío & le défi & accent sur le í \\ \hline
la falta de & le manque de & \\ \hline
\end{tabularx}
\end{center}

\begin{exoresolu}[Traduisez les huit structures soulignées du texte]
Énoncé. Traduisez en français : 1. \esp{menos empleos que sus padres} ; 2. \esp{peor que hace diez años} ; 3. \esp{más altos que en otros países} ; 4. \esp{el sector más dinámico} ; 5. \esp{tantos puestos como el comercio} ; 6. \esp{más horas que los funcionarios} ; 7. \esp{el desafío más urgente} ; 8. \esp{tan simple como olvidada}.
\tcblower
Solution détaillée. 1. « moins d'emplois que leurs parents » : comparaison d'infériorité portant sur un nom, \esp{menos + nom + que}. 2. « pire qu'il y a dix ans » : comparatif irrégulier \esp{peor} (= plus mauvais, pire). 3. « plus élevés que dans d'autres pays » : supériorité sur adjectif, \esp{más + adj. + que}, avec accord (\esp{altos} s'accorde avec \esp{salarios}). 4. « le secteur le plus dynamique » : superlatif relatif, article + \esp{más} + adjectif. 5. « autant de postes que le commerce » : égalité sur un nom, \esp{tantos} (masculin pluriel, accord avec \esp{puestos}). 6. « plus d'heures que les fonctionnaires » : supériorité sur un nom. 7. « le défi le plus urgent » : superlatif relatif. 8. « aussi simple qu'oubliée » : égalité sur un adjectif, \esp{tan... como}, \esp{tan} invariable.
\end{exoresolu}

\courssub{Thème : cinq phrases à traduire}

\begin{exercice}[Thème guidé : l'emploi des jeunes]
Traduisez en espagnol :
\begin{enumerate}
\item Le chômage des jeunes est plus élevé que celui des adultes.
\item Les jeunes de Douala ont autant de diplômes que les jeunes de Yaoundé.
\item C'est la pire année de la décennie pour l'emploi.
\item Plus de mille jeunes ont participé au forum de l'emploi.
\item María parle espagnol mieux que son frère, mais elle écrit moins vite que lui.
\end{enumerate}
\end{exercice}

\begin{corrige}[Thème guidé]
\begin{enumerate}
\item \esp{El desempleo de los jóvenes es más alto que el de los adultos.} « Celui des adultes » se rend par l'article \esp{el} reprenant \esp{desempleo} : ne répétez pas le nom. Supériorité sur adjectif : \esp{más alto que}.
\item \esp{Los jóvenes de Duala tienen tantos diplomas como los jóvenes de Yaoundé.} Égalité sur un nom : \esp{tantos} (accord masculin pluriel avec \esp{diplomas}) et non \esp{tan}.
\item \esp{Es el peor año de la década para el empleo.} Comparatif irrégulier \esp{peor} avec valeur de superlatif relatif grâce à l'article \esp{el} : « la pire année ».
\item \esp{Más de mil jóvenes participaron en el foro del empleo.} Devant le chiffre \esp{mil}, on dit \esp{más de} et jamais \esp{más que}. Passé simple de \esp{participar} : \esp{participaron}.
\item \esp{María habla español mejor que su hermano, pero escribe menos rápido que él.} « Mieux » (adverbe) = \esp{mejor} ; « moins vite » = \esp{menos rápido} (adverbe invariable). Après \esp{que}, le pronom tonique \esp{él} avec accent.
\end{enumerate}
\end{corrige}

\courssub{Correction collective : justification grammaticale de chaque choix}

En classe, la correction se fait au tableau, phrase par phrase, et chaque élève doit pouvoir \emph{justifier} sa structure. Voici la grille de justification attendue.

\begin{center}
\begin{tabularx}{\linewidth}{|X|X|X|}\hline
\rowcolor{popblueL} Français & Español & Justification \\ \hline
plus élevé que & \esp{más alto que} & supériorité + adjectif \\ \hline
autant de diplômes que & \esp{tantos diplomas como} & égalité + nom : accord de \esp{tanto} \\ \hline
la pire année & \esp{el peor año} & irrégulier \esp{peor} + article \\ \hline
plus de mille jeunes & \esp{más de mil jóvenes} & chiffre : \esp{de}, pas \esp{que} \\ \hline
mieux que & \esp{mejor que} & adverbe irrégulier (\esp{bien} $\to$ \esp{mejor}) \\ \hline
aussi simple que & \esp{tan simple como} & égalité + adjectif : \esp{tan} invariable \\ \hline
le secteur le plus dynamique & \esp{el sector más dinámico} & superlatif relatif : article + \esp{más} \\ \hline
moins d'emplois que & \esp{menos empleos que} & infériorité + nom \\ \hline
\end{tabularx}
\end{center}

\begin{aretenir}[Bilan : les cinq pièges récurrents relevés dans les productions]
\begin{enumerate}
\item \esp{tan trabajo como} au lieu de \esp{tanto trabajo como} : \esp{tan} ne précède jamais un nom.
\item \esp{más que veinte} au lieu de \esp{más de veinte} : devant un chiffre, toujours \esp{de}.
\item Oublier l'accord de \esp{tanto} : \esp{tanta gente}, \esp{tantos empleos}, \esp{tantas familias}.
\item Traduire « meilleur » et « mieux » par \esp{más bueno} / \esp{más bien} : l'espagnol impose \esp{mejor} dans les deux cas (adjectif et adverbe).
\item Oublier l'article dans le superlatif relatif : « le plus grand marché » = \esp{el mercado más grande}, jamais \esp{mercado más grande} seul.
\end{enumerate}
\end{aretenir}

\begin{experience}[Correction en ateliers]
Par groupes de quatre, échangez vos copies du thème. Chaque groupe relève les erreurs, les classe selon les cinq pièges du bilan ci-dessus, puis propose une version corrigée collective. Un rapporteur présente au tableau la phrase la plus difficile du groupe et justifie la structure choisie devant la classe.
\end{experience}

Vous êtes maintenant capable de traduire toute comparaison et tout superlatif, en thème comme en version, en justifiant chaque choix. La section suivante mettra cette précision linguistique au service d'un document concret du monde du travail : le \esp{currículum vítae} en espagnol.

\coursec{Leçon E17 : Traduction : comparatifs et superlatifs ; Redacción : el currículum vítae}

\courssub{Observation : la comparaison en contexte}

Avant d'étudier les règles, lisons un extrait authentique d'une annonce de stage publiée sur le site d'une entreprise de télécommunications à Douala, puis un échange entre deux candidats. Repérez toutes les structures qui expriment une comparaison.

\begin{textebox}[Annonce de stage et dialogue entre candidats]
\textbf{OFERTA DE PRÁCTICAS : CAMTEL DOUALA}\par
Buscamos un estudiante de Bachillerato o Universidad para prácticas de secretariado bilingüe (francés-español).\par
Requisitos :\par
- Nivel de español \textbf{más alto que} B1 (preferiblemente B2 o \textbf{superior}).\par
- Disponibilidad \textbf{tan inmediata como} sea posible.\par
- Capacidad de organización \textbf{mejor que} la media.\par
- Se valorará experiencia \textbf{menor a} 6 meses, pero \textbf{mayor} motivación.\par
\medskip
\textbf{Diálogo : Aminatou et Jean}\par
— \esp{¿Has visto la oferta? Piden un nivel más alto que el mío.} (Aminatou)\par
— \esp{Tranquila. Tu español es tan bueno como el de un nativo. Y tienes más experiencia que yo.} (Jean)\par
— \esp{Es cierto, pero mi francés es menos fluido que el tuyo.} (Aminatou)\par
— \esp{Da igual. Lo importante es que somos los \textbf{mejores} candidatos. ¡Vamos a enviar el CV hoy mismo!} (Jean)
\end{textebox}

\textbf{Glossaire :} \esp{buscamos} = nous cherchons ; \esp{estudiante} = étudiant(e) ; \esp{requisitos} = prérequis ; \esp{se valorará} = sera valorisé (passif réfléchi) ; \esp{media} = moyenne ; \esp{tranquila} = calme / ne t'inquiète pas ; \esp{da igual} = peu importe / ça ne fait rien.

\textbf{Questions d'observation :}
\begin{enumerate}
\item Relevez quatre comparatifs de supériorité ou d'infériorité dans l'annonce.
\item Quelle structure exprime l'égalité dans « \esp{tan inmediata como sea posible} » ?
\item Identifiez les deux adjectifs irréguliers utilisés dans l'annonce (\esp{mejor}, \esp{mayor}, \esp{menor}...).
\item Dans le dialogue, comment Jean exprime-t-il le superlatif de supériorité ?
\end{enumerate}

\begin{corrige}[Questions d'observation]
\begin{enumerate}
\item \esp{más alto que} (supériorité), \esp{mejor que} (supériorité irrégulière), \esp{menor a} (infériorité irrégulière), \esp{menos fluido que} (infériorité).
\item \esp{tan... como} (égalité d'adjectif/adverbe).
\item \esp{mejor} (meilleur), \esp{mayor} (plus grand/âgé/supérieur), \esp{menor} (plus petit/jeune/inférieur).
\item \esp{los mejores candidatos} (superlatif relatif : article défini + \esp{mejor} + nom).
\end{enumerate}
\end{corrige}

\courssub{Les comparatifs réguliers}

\begin{propriete}[Comparatifs de supériorité, d'infériorité et d'égalité]
En espagnol, la comparaison s'établit à l'aide de deux termes corrélatifs encadrant l'adjectif, l'adverbe ou le nom.
\begin{enumerate}
\item \textbf{Supériorité} : \esp{más + {adjectif/adverbe/nom} + que} \\
\esp{Este trabajo es más interesante que el otro.} (Ce travail est plus intéressant que l'autre.) \\
\esp{Gana más dinero que yo.} (Il gagne plus d'argent que moi.)
\item \textbf{Infériorité} : \esp{menos + {adjectif/adverbe/nom} + que} \\
\esp{Este puesto es menos estresante que aquél.} (Ce poste est moins stressant que celui-là.) \\
\esp{Trabaja menos horas que su jefe.} (Il travaille moins d'heures que son patron.)
\item \textbf{Égalité (adjectifs/adverbes)} : \esp{tan + {adjectif/adverbe} + como} \\
\esp{Mi jefa es tan exigente como la tuya.} (Ma chef est aussi exigeante que la tienne.) \\
\esp{Habla tan rápido como un locutor.} (Il parle aussi vite qu'un speaker.)
\item \textbf{Égalité (noms)} : \esp{tanto(s)/tanta(s) + {nom} + como} (accord en genre et nombre) \\
\esp{Tengo tanta experiencia como tú.} (J'ai autant d'expérience que toi.) \\
\esp{Tiene tantos títulos como su hermano.} (Il a autant de diplômes que son frère.)
\end{enumerate}
\end{propriete}

\begin{attention}[Pièges francophones : \esp{que} vs \emph{que}, \esp{como} vs \emph{comme}]
\begin{itemize}
\item Après \esp{más/menos}, on utilise toujours \esp{que} (jamais \esp{de}). \textbf{Erreur :} \esp{*más de 18 años} (sauf pour « plus de » quantité chiffrée : \esp{más de 18 años} = plus de 18 ans). \textbf{Correct (comparaison) :} \esp{Tengo más años que tú.}
\item Après \esp{tan/tanto}, on utilise \esp{como} (jamais \esp{que}). \textbf{Erreur :} \esp{*tan bueno que yo}. \textbf{Correct :} \esp{tan bueno como yo}.
\item \esp{Tanto} s'accorde : \esp{tanto tiempo, tanta paciencia, tantos libros, tantas ganas}.
\item \esp{Que} / \esp{como} sont suivis du \textbf{sujet} (pronom tonique ou nom), pas du COD. \esp{Más alto que YO} (et non *que me).
\end{itemize}
\end{attention}

\begin{exemplebox}[Comparatifs avec pronoms toniques]
\begin{center}
\begin{tabular}{|l|l|}
\hline
\textbf{Espagnol} & \textbf{Français} \\
\hline
\esp{Él es más alto que yo.} & Il est plus grand que moi. \\
\hline
\esp{Nosotros trabajamos menos que ellos.} & Nous travaillons moins qu'eux. \\
\hline
\esp{Ella tiene tanta paciencia como tú.} & Elle a autant de patience que toi. \\
\hline
\end{tabular}
\end{center}
\end{exemplebox}

\courssub{Les comparatifs irréguliers}

\begin{propriete}[Mejor, peor, mayor, menor]
Quatre adjectifs/adverbes ont des formes comparatives irrégulières, héritées du latin. Elles remplacent \esp{más bueno, más malo, más grande, más pequeño} dans leurs sens principaux.
\begin{center}
\begin{tabular}{|c|c|c|c|}
\hline
\rowcolor{popblueL} \textbf{Positif} & \textbf{Comparatif irrégulier} & \textbf{Sens principal} & \textbf{Forme régulière (autres sens)} \\
\hline
\esp{bueno / bien} & \esp{mejor} & meilleur / mieux & \esp{más bueno} = plus brave, plus gentil \\
\hline
\esp{malo / mal} & \esp{peor} & pire / plus mal & \esp{más malo} = plus méchant, plus grave \\
\hline
\esp{grande} & \esp{mayor} & plus âgé, supérieur (rang/taille) & \esp{más grande} = plus grand (taille physique) \\
\hline
\esp{pequeño} & \esp{menor} & plus jeune, inférieur (rang/taille) & \esp{más pequeño} = plus petit (taille physique) \\
\hline
\end{tabular}
\end{center}
\end{propriete}

\begin{aretenir}[Emploi de \esp{mayor/menor} vs \esp{más grande/más pequeño}]
\begin{itemize}
\item \textbf{Âge / Antériorité / Rang :} \esp{mayor / menor}. \esp{Mi hermano mayor} (mon frère aîné), \esp{un problema mayor} (un problème plus grave/majeur), \esp{menor de edad} (mineur).
\item \textbf{Taille physique / Volume :} \esp{más grande / más pequeño}. \esp{Esta caja es más grande que esa.}
\item \textbf{Bon / Mauvais (qualité morale, goût) :} \esp{más bueno / más malo}. \esp{Un niño más bueno} (un enfant plus sage/gentil), \esp{una manzana más mala} (une pomme plus pourrie).
\item \textbf{Bon / Mauvais (qualité, compétence, santé) :} \esp{mejor / peor}. \esp{Un estudiante mejor} (un meilleur élève), \esp{Me siento mejor} (je me sens mieux).
\end{itemize}
\end{aretenir}

\begin{exoresolu}[Choisir la forme correcte]
\textbf{Énoncé.} Complétez avec la forme comparative appropriée (irrégulière ou régulière) :\par
1. Mi abuelo tiene 80 años, mi abuela 78. Mi abuelo es \_\_\_\_\_ que mi abuela. (âgé)\par
2. Esta caja pesa 10 kg, aquella 5 kg. Esta caja es \_\_\_\_\_ que aquella. (lourde/grande physiquement)\par
3. Juan estudia mucho, Pedro estudia poco. Juan es \_\_\_\_\_ estudiante que Pedro. (bon)\par
4. Este examen es \_\_\_\_\_ que el del año pasado. (difficile / mauvais)\par
5. Mi salario es 200.000 FCFA, el tuyo 150.000. Mi salario es \_\_\_\_\_ que el tuyo. (supérieur)\par
\tcblower
\textbf{Solution.}\par
1. \esp{mayor} (âge). 2. \esp{más grande} ou \esp{más pesada} (taille/poids physique, pas rang). 3. \esp{mejor} (compétence/qualité). 4. \esp{peor} (qualité négative) ou \esp{más difícil} (degré de difficulté). 5. \esp{mayor} (quantité, rang, supériorité numérique) ou \esp{más alto}.
\end{exoresolu}

\courssub{Les superlatifs : relatif et absolu}

\begin{propriete}[Superlatif relatif : le plus / le moins \emph{de/dans}]
Il identifie un élément au sein d'un groupe. Structure : \esp{el/la/los/las + más/menos + adjectif + {de/en}}.
\begin{itemize}
\item \esp{de} pour les groupes de personnes/objets définis : \esp{El mejor estudiante \textbf{de} la clase.} (Le meilleur élève de la classe.)
\item \esp{en} pour les lieux/espaces : \esp{El edificio más alto \textbf{en} Yaoundé.} (Le plus haut bâtiment de Yaoundé.)
\item Accord de l'article et de l'adjectif avec le nom noyau.
\item Irréguliers : \esp{el mejor, la peor, los mayores, las menores}.
\end{itemize}
\end{propriete}

\begin{propriete}[Superlatif absolu : très, extrêmement (sans comparaison)]
Il exprime un degré très élevé sans comparer à un groupe. Trois formes principales :
\begin{enumerate}
\item \textbf{Suffixation en \esp{-ísimo/a/os/as}} (morphologique, expressif, littéraire) : \esp{rápido $\rightarrow$ rapidísimo}, \esp{fuerte $\rightarrow$ fortísimo}, \esp{bueno $\rightarrow$ bonísimo} (irrégulier), \esp{malo $\rightarrow$ malísimo}, \esp{grande $\rightarrow$ grandísimo}, \esp{pequeño $\rightarrow$ pequeñísimo}. \textbf{Accent obligatoire} sur le \esp{í}.
\item \textbf{Adverbe \esp{muy} + adjectif} (standard, neutre) : \esp{muy rápido, muy buena, muy interesantes}.
\item \textbf{Adverbes intensifs} : \esp{sumamente, extremadamente, altamente, profundamente} + adjectif (registre formel/technique). \esp{Sumamente urgente, altamente cualificado}.
\end{enumerate}
\end{propriete}

\begin{attention}[Accords et pièges du superlatif absolu en \esp{-ísimo}]
\begin{itemize}
\item L'adjectif perd son accent s'il en a un : \esp{fácil $\rightarrow$ facilísimo} (pas *fácilísimo), \esp{inglés $\rightarrow$ inglesísimo}.
\item Changement orthographique pour garder le son : \esp{fuerte $\rightarrow$ fortísimo} (gu $\rightarrow$ g), \esp{grande $\rightarrow$ grandísimo} (d $\rightarrow$ d), \esp{rico $\rightarrow$ riquísimo} (c $\rightarrow$ qu).
\item \esp{Muy} est invariable. \esp{Muy buenos amigos} (pas *muy buenas amigos).
\item \textbf{Faux ami :} \esp{actualmente} = actuellement (en ce moment), \textbf{pas} « en fait ». Pour « en fait » : \esp{en realidad, de hecho}.
\end{itemize}
\end{attention}

\begin{exemplebox}[Superlatifs en contexte professionnel]
\begin{itemize}
\item \esp{Es el candidato \textbf{más capacitado de todos}.} (Relatif : le plus qualifié de tous.)
\item \esp{La entrevista fue \textbf{sumamente difícil}.} (Absolu formel : extrêmement difficile.)
\item \esp{Tienes un currículum \textbf{riquísimo en experiencias}.} (Absolu \esp{-ísimo} : très riche en expériences.)
\item \esp{Es \textbf{la mejor oferta del mercado}.} (Relatif irrégulier : la meilleure offre du marché.)
\end{itemize}
\end{exemplebox}

\courssub{Traduction guidée : thème et version}

\begin{methode}[Traduire les comparatifs et superlatifs : 4 étapes]
\begin{enumerate}
\item \textbf{Identifier la structure} : comparatif (sup/inf/égalité) ou superlatif (relatif/absolu) ? Régulier ou irrégulier ?
\item \textbf{Repérer le support} : adjectif, adverbe, nom, verbe ? (Détermine \esp{tan/tanto}, accord de \esp{tanto}).
\item \textbf{Choisir le bon connecteur} : \esp{que} (sup/inf), \esp{como} (égalité), \esp{de/en} (superlatif relatif).
\item \textbf{Vérifier les accords} (genre/nombre pour \esp{tanto, el más, -ísimo}) et les irréguliers (\esp{mejor, mayor}...).
\end{enumerate}
\end{methode}

\begin{exoresolu}[Thème : Français $\rightarrow$ Espagnol]
\textbf{Énoncé.} Traduisez en espagnol :\par
1. Ce stage est plus intéressant que le précédent.\par
2. Elle a autant de diplômes que son collègue.\par
3. C'est le pire candidat que nous ayons vu. (Subjonctif après superlatif + \esp{que} + verbe).\par
4. Il parle espagnol aussi bien que moi.\par
5. Nous avons une disponibilité immédiate, très supérieure à la moyenne.\par
6. Ce logiciel est extrêmement facile à utiliser. (Deux versions possibles).\par
\tcblower
\textbf{Solution détaillée.}\par
1. \esp{Este período de prácticas es más interesante que el anterior.} (\esp{prácticas} = stage, \esp{anterior} = précédent).\par
2. \esp{Ella tiene tantos títulos como su compañero.} (\esp{títulos} masc. pl. $\rightarrow$ \esp{tantos}).\par
3. \esp{Es el peor candidato que hayamos visto.} (Superlatif relatif \esp{el peor} + subjonctif \esp{hayamos visto} car opinion/négation implicite « que nous ayons vu » = relativité).\par
4. \esp{Habla español tan bien como yo.} (Adverbe \esp{bien} $\rightarrow$ \esp{tan bien como}).\par
5. \esp{Tenemos disponibilidad inmediata, muy superior a la media.} (\esp{superior a} = supérieur à, structure adjectif + préposition).\par
6. \esp{Este software es facilísimo de usar.} (ou \esp{sumamente fácil de usar} / \esp{extremadamente fácil}). \esp{Fácil} perd l'accent $\rightarrow$ \esp{facilísimo}.
\end{exoresolu}

\begin{exoresolu}[Version : Espagnol $\rightarrow$ Français]
\textbf{Énoncé.} Traduisez en français :\par
1. \esp{Mi jefa es menor que yo, pero tiene más experiencia.}\par
2. \esp{No hay nadie tan puntual como él en la oficina.}\par
3. \esp{El salario ofrecido es inferior a 200.000 francos.}\par
4. \esp{Esta tarea es riquísima en aprendizajes.}\par
5. \esp{Llegaron los mayores primero, luego los menores.}\par
\tcblower
\textbf{Solution.}\par
1. Ma chef est plus jeune que moi, mais elle a plus d'expérience. (\esp{menor} = plus jeune).\par
2. Il n'y a personne d'aussi ponctuel que lui au bureau. (\esp{nadie tan... como} = personne n'est aussi... que).\par
3. Le salaire offert est inférieur à 200.000 francs. (\esp{inferior a} = inférieur à).\par
4. Cette tâche est extrêmement riche en apprentissages. (\esp{riquísima} = superlatif absolu \esp{-ísima} de \esp{rica}).\par
5. Les aînés sont arrivés en premier, puis les cadets. (\esp{los mayores/menores} = les aînés / les cadets / les majeurs / les mineurs selon contexte).
\end{exoresolu}

\begin{exercice}[Entraînement : Thème et Version]
\textbf{A. Thème (Fr $\rightarrow$ Es).}
1. Mon frère cadet est plus grand que moi. (Taille physique).\par
2. Cette entreprise paie moins bien que l'autre.\par
3. Nous cherchons le candidat le plus qualifié de la région.\par
4. Tu as tort : mon espagnol est aussi bon que le tien.\par
5. C'est une décision sumamente importante.\par
\medskip
\textbf{B. Version (Es $\rightarrow$ Fr).}
1. \esp{Tengo menos tiempo que tú para hacer el CV.}\par
2. \esp{Es la mejor opción para tu futuro profesional.}\par
3. \esp{Los requisitos son tan claros como el agua.}\par
4. \esp{Su currículum es malísimo: tiene muchas faltas de ortografía.}\par
5. \esp{Ana es mayor que Luis, pero Luis es más alto.}
\end{exercice}

\begin{corrige}[Exercice : Thème et Version]
\textbf{A. Thème.}
1. \esp{Mi hermano menor es más grande que yo.} (\esp{menor} = cadet, \esp{más grande} = taille physique).\par
2. \esp{Esta empresa paga peor que la otra.} (Irrégulier \esp{peor} pour \esp{mal}).\par
3. \esp{Buscamos al candidato más cualificado de la región.} (Relatif \esp{el más... de}).\par
4. \esp{Te equivocas: mi español es tan bueno como el tuyo.} (\esp{tan bueno como}).\par
5. \esp{Es una decisión sumamente importante.} (Adverbe formel).\par
\medskip
\textbf{B. Version.}
1. J'ai moins de temps que toi pour faire le CV.\par
2. C'est la meilleure option pour ton avenir professionnel.\par
3. Les prérequis sont aussi clairs que de l'eau de roche (très clairs).\par
4. Son CV est très mauvais / exécrable : il a beaucoup de fautes d'orthographe. (\esp{malísimo} = superlatif absolu péjoratif).\par
5. Ana est plus âgée que Luis, mais Luis est plus grand. (\esp{mayor} = âge, \esp{más alto} = taille).
\end{corrige}

\courssub{Le currículum vítae en espagnol}

Dans la section précédente, nous avons traduit des phrases contenant des comparatifs et des superlatifs, en français comme en espagnol. Ces outils de la comparaison nous seront d'ailleurs utiles ici : dans un CV, on présente souvent ses compétences de manière nuancée (\esp{domino el francés mejor que el alemán}). Nous passons maintenant de la traduction à l'\cle{expression écrite professionnelle} : savoir rédiger un \cle{currículum vítae} en espagnol est une compétence concrète du monde du travail, au cœur de ce module 4. Un jour, vous postulerez peut-être à un emploi dans une entreprise hispanophone de Douala, à un stage en Guinée équatoriale ou à une université espagnole : il vous faudra un CV correct, clair et conforme aux usages hispaniques.

\courssub{Qu'est-ce qu'un currículum vítae ?}

\begin{definition}[El currículum vítae]
Le \cle{currículum vítae} (en abrégé \esp{CV}, parfois appelé simplement \esp{currículum}) est un document écrit qui résume, de manière claire et ordonnée, le \cle{parcours}, les \cle{diplômes}, l'\cle{expérience professionnelle} et les \cle{compétences} d'une personne qui postule à un emploi (\esp{solicita un empleo}). En espagnol, l'expression vient du latin et signifie littéralement « course de la vie ».
\end{definition}

\begin{propriete}[Les trois qualités d'un bon CV]
Un CV efficace, en espagnol comme en français, respecte trois principes :
\begin{enumerate}
\item \textbf{La brièveté} : une seule page (\esp{una sola página}), surtout pour un jeune candidat.
\item \textbf{La clarté} : des rubriques (\esp{apartados}) bien séparées, avec des titres en évidence.
\item \textbf{L'ordre chronologique inverse} (\esp{orden cronológico inverso}) : on présente toujours l'élément le plus récent en premier, puis on remonte vers le plus ancien.
\end{enumerate}
\end{propriete}

\courssub{Observation : un extrait de CV modèle}

Avant d'énoncer les règles, observons un extrait de CV rédigé selon les conventions hispaniques. Lisez-le attentivement et repérez la structure, le vocabulaire et la forme des phrases.

\begin{textebox}[Extrait de CV modèle]
\textbf{CURRÍCULUM VÍTAE}\par
\textbf{DATOS PERSONALES}\par
Nombre y apellidos : Aminatou Mbarga Essomba\par
Fecha de nacimiento : 12/03/2007\par
Dirección : B.P. 4025, Yaoundé (Camerún)\par
Teléfono : (+237) 690 00 00 00\par
Correo electrónico : aminatou.mbarga@correo.com\par
Nacionalidad : camerunesa\par
\textbf{FORMACIÓN ACADÉMICA}\par
2023-2024 : Bachillerato, serie A4, Lycée de Nkol-Eton, Yaoundé.\par
2020-2023 : Primer ciclo de secundaria, Collège Bilingue de Mfoundi.\par
\textbf{EXPERIENCIA PROFESIONAL}\par
Julio 2024 : prácticas de secretariado, empresa Camtel, Douala. Tareas : atender a los clientes, organizar archivos, redactar cartas comerciales.\par
\textbf{IDIOMAS}\par
Francés : lengua materna. Español : nivel avanzado. Inglés : nivel intermedio.\par
\textbf{OTROS DATOS DE INTERÉS}\par
Permiso de conducir : categoría B. Disponibilidad inmediata.
\end{textebox}

\textbf{Glossaire :} \esp{apellidos} = noms de famille ; \esp{fecha de nacimiento} = date de naissance ; \esp{dirección} = adresse ; \esp{prácticas} = stage ; \esp{atender a los clientes} = accueillir/servir les clients ; \esp{archivos} = dossiers, fichiers ; \esp{redactar} = rédiger ; \esp{lengua materna} = langue maternelle ; \esp{nivel avanzado / intermedio} = niveau avancé / intermédiaire ; \esp{permiso de conducir} = permis de conduire ; \esp{disponibilidad} = disponibilité.

\textbf{Questions d'observation :}
\begin{enumerate}
\item Combien de rubriques compte cet extrait ? Citez-les dans l'ordre.
\item Dans la rubrique \esp{Formación académica}, quel diplôme est mentionné en premier ? Pourquoi ?
\item Quelle forme verbale est utilisée pour décrire les tâches du stage (\esp{atender, organizar, redactar}) ?
\item La date de naissance est-elle exprimée par une phrase complète ? Comment est-elle présentée ?
\end{enumerate}

\begin{corrige}[Questions d'observation]
\begin{enumerate}
\item Cinq rubriques : \esp{Datos personales}, \esp{Formación académica}, \esp{Experiencia profesional}, \esp{Idiomas}, \esp{Otros datos de interés}.
\item Le \esp{Bachillerato} (2023-2024) est mentionné en premier car c'est le diplôme le plus récent : c'est l'\esp{orden cronológico inverso}.
\item L'infinitif : \esp{atender, organizar, redactar}. C'est la convention pour décrire des tâches dans un CV.
\item Non. On écrit simplement \esp{Fecha de nacimiento : 12/03/2007}, sans phrase comme « je suis né le... ».
\end{enumerate}
\end{corrige}

\courssub{Les six rubriques du CV en espagnol}

\begin{propriete}[Structure standard d'un CV hispanophone]
Un CV en espagnol comporte généralement six rubriques, dans cet ordre :
\begin{enumerate}
\item \esp{Datos personales} : \esp{nombre y apellidos} (prénom et noms), \esp{fecha y lugar de nacimiento} (date et lieu de naissance), \esp{dirección} (adresse), \esp{teléfono}, \esp{correo electrónico}, \esp{nacionalidad}, et éventuellement \esp{estado civil} (état civil : \esp{soltero/a} = célibataire, \esp{casado/a} = marié(e)).
\item \esp{Formación académica} : les diplômes, du plus récent au plus ancien : \esp{Bachillerato}, \esp{Universidad}, \esp{títulos} (titres, diplômes), avec les dates et les établissements.
\item \esp{Experiencia profesional} : les emplois (\esp{empleos}) et stages (\esp{prácticas}), avec le \esp{puesto ocupado} (poste occupé), les \esp{fechas} (dates) et les \esp{tareas realizadas} (tâches accomplies).
\item \esp{Idiomas} : la \esp{lengua materna} puis les langues étrangères avec leur niveau : \esp{básico} (élémentaire), \esp{intermedio} (intermédiaire), \esp{avanzado} (avancé), \esp{bilingüe} (bilingue), et les certifications éventuelles (DELE, par exemple).
\item \esp{Competencias / Habilidades} : l'informatique (\esp{informática} : \esp{manejo de Word y Excel} = maîtrise de Word et Excel), les qualités (\esp{capacidad de trabajo en equipo} = capacité de travail en équipe, \esp{puntualidad} = ponctualité, \esp{responsabilidad} = sens des responsabilités).
\item \esp{Otros datos de interés} : \esp{permiso de conducir}, \esp{aficiones} (loisirs), \esp{referencias} (références), \esp{disponibilidad} (disponibilité).
\end{enumerate}
\end{propriete}

\begin{popfigure}
\begin{center}
\begin{tikzpicture}[x=1cm,y=1cm, font=\small]
\fill[popblueL, draw=popblue, rounded corners=2pt] (0,10.2) rectangle (8,11.2);
\node[align=center, font=\bfseries] at (4,10.7) {DATOS PERSONALES};
\fill[poporangeL, draw=poporange, rounded corners=2pt] (0,8.4) rectangle (8,10.0);
\node[align=center, font=\bfseries] at (4,9.4) {FORMACIÓN ACADÉMICA};
\node[align=center] at (4,8.8) {del más reciente al más antiguo};
\fill[popgreenL, draw=popgreen, rounded corners=2pt] (0,6.4) rectangle (8,8.2);
\node[align=center, font=\bfseries] at (4,7.6) {EXPERIENCIA PROFESIONAL};
\node[align=center] at (4,6.95) {puesto, fechas, tareas (infinitivos)};
\fill[poppurpleL, draw=poppurple, rounded corners=2pt] (0,4.8) rectangle (8,6.2);
\node[align=center, font=\bfseries] at (4,5.7) {IDIOMAS};
\node[align=center] at (4,5.15) {lengua materna + niveles};
\fill[poppinkL, draw=poppink, rounded corners=2pt] (0,3.2) rectangle (8,4.6);
\node[align=center, font=\bfseries] at (4,4.1) {COMPETENCIAS / HABILIDADES};
\node[align=center] at (4,3.55) {informática, cualidades};
\fill[popgoldL, draw=popgold, rounded corners=2pt] (0,1.8) rectangle (8,3.0);
\node[align=center, font=\bfseries] at (4,2.55) {OTROS DATOS DE INTERÉS};
\node[align=center] at (4,2.1) {permiso, aficiones, disponibilidad};
\draw[-{Stealth[length=3mm]}, very thick, popdark] (8.6,10.6) -- (8.6,2.4);
\node[rotate=90, align=center, font=\small\itshape] at (9.15,6.5) {una sola página};
\end{tikzpicture}
\end{center}
\legende{Maquette schématique d'un CV espagnol : les six rubriques empilées, de haut en bas, sur une seule page.}
\end{popfigure}

\begin{center}
\begin{tabularx}{\linewidth}{|X|X|X|}
\hline
\rowcolor{popblueL} \textbf{Español} & \textbf{Français} & \textbf{Remarque} \\
\hline
\esp{formación (académica)} & formation (scolaire) & rubrique des diplômes \\
\hline
\esp{puesto (de trabajo)} & poste (de travail) & \esp{ocupar un puesto} = occuper un poste \\
\hline
\esp{prácticas} & stage & \esp{hacer unas prácticas} = faire un stage \\
\hline
\esp{solicitud (de empleo)} & candidature, demande d'emploi & \esp{solicitar un empleo} = postuler \\
\hline
\esp{entrevista (de trabajo)} & entretien (d'embauche) & \esp{pasar una entrevista} = passer un entretien \\
\hline
\esp{datos personales} & données personnelles & \esp{datos} = données, informations \\
\hline
\esp{fecha de nacimiento} & date de naissance & format : jour/mois/année \\
\hline
\esp{correo electrónico} & courriel, adresse e-mail & on dit aussi \esp{e-mail} \\
\hline
\esp{estado civil} & état civil & \esp{soltero/a}, \esp{casado/a} \\
\hline
\esp{lengua materna} & langue maternelle & première langue apprise \\
\hline
\esp{nivel básico / intermedio / avanzado} & niveau élémentaire / intermédiaire / avancé & pour les langues \\
\hline
\esp{manejo de Word y Excel} & maîtrise de Word et Excel & \esp{manejar} = maîtriser, manier \\
\hline
\esp{capacidad de trabajo en equipo} & capacité de travail en équipe & qualité très demandée \\
\hline
\esp{permiso de conducir} & permis de conduire & \esp{categoría B} = permis B \\
\hline
\esp{disponibilidad inmediata} & disponibilité immédiate & atout pour le candidat \\
\hline
\esp{referencias} & références & personnes à contacter \\
\hline
\end{tabularx}
\end{center}

\courssub{Les conventions de rédaction}

\begin{propriete}[Comment rédiger les tâches et les descriptions]
Dans un CV espagnol, on n'écrit pas de phrases longues avec sujet et verbe conjugué. Pour décrire les tâches accomplies, on utilise :
\begin{itemize}
\item soit des \cle{infinitifs}, de préférence des \cle{verbes d'action} : \esp{atender a los clientes} (accueillir les clients), \esp{organizar archivos} (organiser les dossiers), \esp{gestionar la correspondencia} (gérer le courrier), \esp{redactar informes} (rédiger des rapports), \esp{vender productos} (vendre des produits) ;
\item soit des \cle{noms} : \esp{atención al cliente} (accueil des clients), \esp{organización de archivos}, \esp{gestión de correspondencia}.
\end{itemize}
Chaque rubrique se présente en lignes courtes, précédées de dates. On évite absolument les paragraphes rédigés.
\end{propriete}

\begin{exemplebox}[Verbes d'action utiles pour un CV]
\begin{itemize}
\item \esp{gestionar} = gérer $\rightarrow$ \esp{Gestionar la agenda del director.} « Gérer l'agenda du directeur. »
\item \esp{organizar} = organiser $\rightarrow$ \esp{Organizar reuniones y archivos.} « Organiser des réunions et des dossiers. »
\item \esp{atender (a)} = accueillir, s'occuper de $\rightarrow$ \esp{Atender al público.} « Accueillir le public. »
\item \esp{redactar} = rédiger $\rightarrow$ \esp{Redactar cartas comerciales.} « Rédiger des lettres commerciales. »
\item \esp{colaborar} = collaborer $\rightarrow$ \esp{Colaborar con el equipo de ventas.} « Collaborer avec l'équipe de vente. »
\item \esp{traducir} = traduire $\rightarrow$ \esp{Traducir documentos del francés al español.} « Traduire des documents du français vers l'espagnol. »
\item \esp{supervisar} = superviser $\rightarrow$ \esp{Supervisar el stock.} « Superviser le stock. »
\item \esp{resolver} = résoudre $\rightarrow$ \esp{Resolver incidencias.} « Résoudre des incidents. »
\end{itemize}
\end{exemplebox}

\begin{methode}[Rédiger un CV en espagnol, étape par étape]
\begin{enumerate}
\item \textbf{Préparer une fiche brouillon} en français : liste de vos diplômes, expériences, langues, qualités, avec les dates exactes.
\item \textbf{Tracer les six rubriques} sur la page, dans l'ordre : \esp{Datos personales}, \esp{Formación académica}, \esp{Experiencia profesional}, \esp{Idiomas}, \esp{Competencias}, \esp{Otros datos de interés}.
\item \textbf{Remplir chaque rubrique du plus récent au plus ancien} (\esp{orden cronológico inverso}), avec des lignes courtes : date, intitulé, lieu.
\item \textbf{Décrire les tâches à l'infinitif} avec des verbes d'action : \esp{atender a los clientes}, \esp{organizar archivos}. Jamais de phrases longues.
\item \textbf{Vérifier la longueur} : une seule page. Supprimer tout ce qui n'est pas utile au poste visé.
\item \textbf{Relire} : accents, accords (\esp{camerunesa} si la candidate est une femme), format des dates (jour/mois/année).
\end{enumerate}
\end{methode}

\begin{attention}[Faux amis et calques fréquents]
\begin{itemize}
\item Pour la rubrique des diplômes, on dit \esp{formación (académica)} et non \esp{educación} : \esp{educación} désigne plutôt l'éducation au sens large (manners, valeurs). Un \esp{hombre bien educado} est un homme poli, bien élevé.
\item \esp{Actualmente} signifie « actuellement, en ce moment », et non « actuellement » au sens de « de nos jours » dans tous les contextes ; surtout, ne jamais l'utiliser pour « en fait » (\esp{en realidad}).
\item Ne traduisez pas mot à mot « je suis né le 12 mars 2007 » : dans un CV, on écrit simplement \esp{Fecha de nacimiento : 12/03/2007}. La formule \esp{Nací el 12 de marzo de 2007} n'appartient pas au style du CV.
\item Attention au format des dates : en espagnol, on écrit le jour d'abord : \esp{12/03/2007} = le 12 mars 2007 (et non le 3 décembre).
\item \esp{Asistir a} signifie « assister à, être présent à », et non « aider » (qui se dit \esp{ayudar}).
\item \esp{Una entrevista} est un entretien ; « une interview » journalistique se dit aussi \esp{entrevista}, mais « passer un entretien d'embauche » = \esp{pasar una entrevista de trabajo}.
\item \esp{Carrera} signifie « carrière universitaire » (ex: \esp{carrera de Derecho} = licence de droit) ou « course », pas « carrière professionnelle » (\esp{trayectoria profesional}).
\end{itemize}
\end{attention}

\begin{savaistu}[CV et lettre de motivation dans le monde hispanophone]
Dans le monde hispanophone, on joint presque toujours au CV une \esp{carta de presentación}, c'est-à-dire une lettre de motivation dans laquelle le candidat explique pourquoi il postule. En Espagne, le modèle \textbf{Europass} est le format européen officiel reconnu par les employeurs et les administrations : il propose des rubriques normalisées, très proches de celles que nous venons d'étudier. Autre différence culturelle avec la France : en Espagne et en Amérique latine, il est encore courant de joindre une \esp{fotografía} (photo d'identité) au CV et d'indiquer son \esp{estado civil}, alors qu'en France ces éléments tendent à disparaître pour éviter les discriminations. Au Cameroun, où les entreprises hispanophones et les échanges avec la Guinée équatoriale se développent, un jeune diplômé capable de présenter un CV en espagnol possède un atout réel : l'espagnol est la langue officielle du seul pays hispanophone d'Afrique subsaharienne, voisin direct du Cameroun.
\end{savaistu}

\begin{exoresolu}[Compléter et corriger un CV]
\textbf{Énoncé.} Voici la rubrique « expérience » rédigée par un élève de Terminale. Corrigez les trois erreurs et réécrivez-la correctement :\par
\esp{EXPERIENCIA PROFESIONAL : Agosto 2024 : Yo trabajé como vendedor en una tienda de Douala. Mis tareas eran vender productos y atendí a los clients. Junio 2023 : prácticas en un banco de Yaoundé.}
\tcblower
\textbf{Solution détaillée.}\par
Erreur 1 : \esp{Yo trabajé} est une phrase conjuguée avec sujet ; dans un CV, on utilise un nom ou une structure courte : \esp{vendedor en una tienda de Douala}.\par
Erreur 2 : mélange de formes (\esp{eran} + \esp{atendí}) ; les tâches se mettent à l'infinitif : \esp{vender productos y atender a los clientes}.\par
Erreur 3 : l'ordre chronologique est correct (2024 avant 2023), mais il manque le poste occupé pour le stage.\par
Erreur 4 (bonus) : \esp{clients} s'écrit \esp{clientes} en espagnol.\par
\textbf{Version corrigée :}\par
\esp{EXPERIENCIA PROFESIONAL}\par
\esp{Agosto 2024 : vendedor, tienda de ropa, Douala. Tareas : vender productos, atender a los clientes.}\par
\esp{Junio 2023 : prácticas de auxiliar administrativo, banco de Yaoundé. Tareas : organizar archivos, gestionar la correspondencia.}
\end{exoresolu}

\begin{exercice}[À vous de rédiger]
Rédigez votre propre CV en espagnol, en suivant la maquette étudiée : les six rubriques, l'ordre chronologique inverse, les tâches à l'infinitif, le tout sur une page. Vous indiquerez au moins deux langues avec leur niveau, une expérience (réelle ou imaginée : stage, petit travail, activité associative) et deux qualités. Votre professeur vérifiera : le vocabulaire des rubriques, le format des dates et l'emploi des verbes d'action.
\end{exercice}

\begin{aretenir}[L'essentiel de la leçon E17]
\textbf{Comparatifs :} \esp{más/menos + adj/adv + que} (sup/inf) ; \esp{tan + adj/adv + como} (égalité) ; \esp{tanto(s)/tanta(s) + nom + como} (égalité nom). \textbf{Irréguliers :} \esp{mejor/peor} (qualité), \esp{mayor/menor} (âge/rang). \textbf{Superlatif relatif :} \esp{el/la/los/las + más/menos + adj + de/en}. \textbf{Superlatif absolu :} \esp{-ísimo} (accent sur í, accords), \esp{muy + adj}, \esp{sumamente + adj}. \textbf{CV :} 6 rubriques (\esp{Datos personales, Formación académica, Experiencia profesional, Idiomas, Competencias, Otros datos}), \esp{orden cronológico inverso}, infinitifs d'action, une page, pas de phrases conjuguées.
\end{aretenir}

\coursec{Production guidée : mon CV et mon autoportrait}

Dans la section précédente, nous avons découvert les rubriques du \esp{currículum vítae} en espagnol (\esp{datos personales}, \esp{formación académica}, \esp{experiencia profesional}, \esp{idiomas}, \esp{otros datos de interés}) et le vocabulaire qui leur est associé. Nous avons aussi appris, dans les sections 2 et 3, à construire les comparatifs (\esp{más… que}, \esp{tan… como}, \esp{mejor}, \esp{peor}) et les superlatifs (\esp{el más… de}, \esp{-ísimo}). Le moment est venu de réunir ces deux savoir-faire dans une production complète et authentique : rédiger votre propre CV en espagnol, puis vous présenter à l'oral comme dans un véritable \esp{entrevista de trabajo}. Suivez les cinq étapes de la frise ci-dessous : chacune correspond à une sous-partie de cette leçon.

\begin{popfigure}
\begin{tikzpicture}[node distance=0.4cm]
\tikzset{etape/.style={rectangle, rounded corners=4pt, draw=popblue, fill=popblueL, text=popdark, align=center, font=\small, minimum height=1.2cm, minimum width=3cm},
fleche/.style={-{Stealth[length=3mm]}, thick, poporange}}
\node[etape, align=center] (e1){\textbf{Étape 1}\\ Fiche d'identité\\ (en français)};
\node[etape, right=of e1, align=center] (e2){\textbf{Étape 2}\\ Transposition\\ (rubriques espagnoles)};
\node[etape, right=of e2, align=center] (e3){\textbf{Étape 3}\\ Rédaction\\ du CV};
\node[etape, right=of e3, align=center] (e4){\textbf{Étape 4}\\ Autoportrait\\ oral (4--5 phrases)};
\node[etape, right=of e4, align=center] (e5){\textbf{Étape 5}\\ Simulation\\ d'entretien};
\draw[fleche] (e1) -- (e2);
\draw[fleche] (e2) -- (e3);
\draw[fleche] (e3) -- (e4);
\draw[fleche] (e4) -- (e5);
\node[below=0.35cm of e3, font=\small\itshape, text=popmuted] {Relecture croisée à chaque étape (grille d'évaluation)};
\end{tikzpicture}
\legende{Frise des étapes de production : de la fiche d'identité à la simulation d'entretien}
\end{popfigure}

\courssub{Étape 1 : la fiche d'identité professionnelle en français}

Avant d'écrire un seul mot en espagnol, il faut rassembler ses informations \textbf{en français}. C'est une règle d'or de la traduction : on ne cherche jamais ses idées et sa langue en même temps. Remplissez la fiche suivante avec vos propres données (réelles ou imaginaires si vous souhaitez vous projeter dans un profil professionnel).

\begin{center}
\begin{tabularx}{\linewidth}{|l|X|}
\hline
\rowcolor{popblueL} \textbf{Rubrique (français)} & \textbf{À compléter} \\
\hline
Nom et prénom & \\
\hline
Date et lieu de naissance & \\
\hline
Adresse, téléphone, courriel & \\
\hline
Diplômes (du plus récent au plus ancien) & \\
\hline
Établissements fréquentés & \\
\hline
Expériences (stages, petits boulots, associations) & \\
\hline
Langues parlées et niveau & \\
\hline
Compétences (informatique, conduite, qualités) & \\
\hline
Centres d'intérêt & \\
\hline
Objectif professionnel & \\
\hline
\end{tabularx}
\end{center}

\begin{exemplebox}[Exemple de fiche remplie (profil fictif)]
\textbf{Nom} : Mbarga Etoundi, Laure. \textbf{Née le} : 12 mars 2007 à Yaoundé. \textbf{Diplômes} : \esp{Probatoire} A (2024), BEPC (2022). \textbf{Établissement} : Lycée bilingue d'Essos. \textbf{Expérience} : vendeuse au marché Mokolo pendant les vacances (2023) ; secrétaire du club d'espagnol. \textbf{Langues} : français (langue maternelle), \esp{español} (bon niveau), anglais (niveau intermédiaire), ewondo (langue maternelle). \textbf{Compétences} : traitement de texte, ponctualité, sens du contact. \textbf{Objectif} : devenir guide touristique bilingue.
\end{exemplebox}

\courssub{Étape 2 : la transposition dans les rubriques espagnoles}

Chaque information de la fiche française trouve sa place dans une rubrique espagnole. Le tableau ci-dessous récapitule les correspondances essentielles vues à la section 5, avec les expressions-clés à réutiliser.

\begin{center}
\begin{tabularx}{\linewidth}{|X|X|X|}
\hline
\rowcolor{popblueL} \textbf{Français} & \textbf{Español} & \textbf{Remarque} \\
\hline
État civil / Coordonnées & Datos personales & \esp{Fecha de nacimiento}, \esp{dirección}, \esp{teléfono}, \esp{correo electrónico} \\
\hline
Formation / Études & Formación académica & Ordre antichronologique : du diplôme le plus récent au plus ancien \\
\hline
Expérience professionnelle & Experiencia profesional & Verbes au passé : \esp{trabajé}, \esp{he trabajado}, \esp{colaboré} \\
\hline
Langues & Idiomas & \esp{nivel básico / intermedio / avanzado}, \esp{lengua materna} \\
\hline
Compétences & Competencias / Conocimientos & \esp{informática}, \esp{permiso de conducir} \\
\hline
Centres d'intérêt & Otros datos de interés & \esp{aficiones}: \esp{el fútbol}, \esp{la lectura}, \esp{la música} \\
\hline
\end{tabularx}
\end{center}

\begin{propriete}[Les formules de transposition]
Pour passer du français à l'espagnol dans un CV, on utilise des tournures nominales brèves, sans phrase complète :
\begin{itemize}
\item « Baccalauréat série A » $\rightarrow$ \esp{Bachillerato, serie A} ;
\item « niveau intermédiaire en anglais » $\rightarrow$ \esp{inglés: nivel intermedio} ;
\item « stage de deux mois à la SABC » $\rightarrow$ \esp{prácticas de dos meses en la SABC} ;
\item « permis de conduire catégorie B » $\rightarrow$ \esp{permiso de conducir, categoría B}.
\end{itemize}
\end{propriete}

\begin{attention}[Faux amis et pièges de transposition]
\begin{itemize}
\item « Une formation » se dit \esp{formación}, jamais \esp{una formación} dans le sens de « diplôme » seul ; on précise : \esp{formación académica}.
\item « Un stage » se traduit par \esp{unas prácticas} ou \esp{un período de prácticas}, et non par \esp{un estadio} (qui signifie « un stade »).
\item « Actuellement » se dit \esp{actualmente}, mais attention : cela signifie « en ce moment », pas « en réalité ».
\item Les mois et les langues ne prennent \textbf{pas de majuscule} en espagnol : \esp{marzo}, \esp{español}, \esp{inglés}.
\item « Né le 12 mars 2007 » $\rightarrow$ \esp{Nacida el 12 de marzo de 2007} : n'oubliez pas l'accord (\esp{nacido/nacida}) et la préposition \esp{de} entre les éléments de la date.
\end{itemize}
\end{attention}

\courssub{Étape 3 : la rédaction du CV complet}

Vous rédigez maintenant votre CV sur \textbf{une page}, en respectant les conventions : rubriques en gras, informations brèves, ordre antichronologique pour la formation et l'expérience, aucune faute d'accord ni d'accent.

\begin{exoresolu}[Rédiger les rubriques « Formación » et « Idiomas » du CV de Laure]
Énoncé. À partir de la fiche d'identité de Laure Mbarga (exemple de l'étape 1), rédigez en espagnol les rubriques \esp{Formación académica} et \esp{Idiomas} de son CV.
\tcblower
\textbf{Solution détaillée.}

\textbf{Formación académica} (ordre antichronologique, du plus récent au plus ancien) :

\esp{2024 : Probatorio, serie A, Liceo bilingüe de Essos, Yaundé.} $\rightarrow$ \esp{2024 : Probatorio, serie A, Liceo bilingüe de Essos, Yaundé.}

\esp{2022 : BEPC, Liceo bilingüe de Essos, Yaundé.} $\rightarrow$ \esp{2022 : BEPC, Liceo bilingüe de Essos, Yaundé.}

\textbf{Idiomas} (langue + deux-points + niveau) :

\esp{Francés: lengua materna.} « Français : langue maternelle. »

\esp{Español: nivel avanzado.} « Espagnol : niveau avancé. »

\esp{Inglés: nivel intermedio.} « Anglais : niveau intermédiaire. »

\esp{Ewondo: lengua materna.} « Ewondo : langue maternelle. »

Remarques : \esp{Yaundé} est la forme espagnole de Yaoundé ; \esp{bilingüe} porte une tréma sur le \emph{u} ; les noms de langues s'écrivent sans majuscule.
\end{exoresolu}

\begin{exercice}[À vous : le CV complet]
Rédigez votre propre \esp{currículum vítae} complet en espagnol (une page maximum), à partir de votre fiche d'identité de l'étape 1. Votre CV doit contenir obligatoirement les cinq rubriques : \esp{Datos personales}, \esp{Formación académica}, \esp{Experiencia profesional}, \esp{Idiomas}, \esp{Otros datos de interés}. Ajoutez une ligne \esp{Objetivo profesional}.
\end{exercice}

\courssub{Étape 4 : l'autoportrait oral en quatre phrases}

Le CV est un document écrit ; mais dans un entretien, il faut aussi savoir \textbf{se présenter à l'oral}. Nous allons réutiliser les comparatifs et superlatifs étudiés dans cette leçon pour construire un autoportrait dynamique et valorisant.

Observez d'abord ces modèles :

\begin{exemplebox}[Phrases de présentation avec comparatifs et superlatifs]
\esp{Hablo español mejor que el año pasado.} « Je parle espagnol mieux que l'année dernière. » (comparatif irrégulier \esp{mejor que})

\esp{El español es la asignatura que más me gusta.} « L'espagnol est la matière que je préfère / qui me plaît le plus. » (superlatif relatif)

\esp{Soy una persona puntual y trabajadora.} « Je suis quelqu'un de ponctuel et travailleur. » (accord : \esp{trabajadora} au féminin pour Laure)

\esp{Tengo un nivel intermedio de inglés.} « J'ai un niveau intermédiaire en anglais. »

\esp{Soy tan responsable como mis compañeros mayores.} « Je suis aussi responsable que mes camarades plus âgés. » (comparatif d'égalité)

\esp{Mi experiencia más importante es mi trabajo en el mercado.} « Mon expérience la plus importante est mon travail au marché. » (superlatif relatif)
\end{exemplebox}

\begin{methode}[Se présenter à l'oral : le schéma en cinq temps]
Pour une présentation orale réussie (une minute environ), suivez toujours le même schéma :
\begin{enumerate}
\item \textbf{Identité} : \esp{Me llamo…, tengo … años, soy de…}
\item \textbf{Formación} : \esp{He estudiado en…, tengo el Probatorio de la serie A.}
\item \textbf{Experiencia} : \esp{He trabajado como…, mi experiencia más importante es…}
\item \textbf{Competencias} : \esp{Hablo francés, español e inglés. Sé utilizar el ordenador.}
\item \textbf{Objetivo profesional} : \esp{Quiero trabajar como… porque soy la persona más indicada para este puesto.}
\end{enumerate}
\end{methode}

\begin{textebox}[Trame de présentation orale à compléter]
Me llamo \ldots{} y tengo \ldots{} años. Soy de Yaundé, la capital de Camerún. He estudiado en el Liceo bilingüe de \ldots{} y tengo el Probatorio de la serie A. Mi experiencia más importante es \ldots{} porque aprendí a trabajar con el público. Hablo francés, español e inglés: hablo español mejor que el año pasado y quiero seguir mejorando. Soy una persona puntual, trabajadora y muy motivada. Creo que soy la persona más indicada para este puesto porque me gusta el contacto con la gente y aprendo más rápido que nadie. ¡Muchas gracias por su atención!
\end{textebox}

\textbf{Glossaire} : \esp{aprendí} = j'ai appris ; \esp{seguir mejorando} = continuer à progresser ; \esp{motivada} = motivée ; \esp{la persona más indicada} = la personne la plus qualifiée/indiquée ; \esp{más rápido que nadie} = plus vite que quiconque ; \esp{por su atención} = pour votre attention.

\begin{attention}[À l'oral : ne lisez pas votre CV !]
L'erreur classique du candidat débutant est de réciter son CV ligne par ligne. Un entretien n'est pas une lecture :
\begin{itemize}
\item \textbf{Parcourez} votre parcours avec des phrases complètes, au \textbf{présent} (\esp{soy}, \esp{tengo}, \esp{hablo}) et au \textbf{pretérito perfecto} pour le parcours accompli : \esp{He estudiado en el Liceo de Essos.} « J'ai étudié au Lycée d'Essos. » ; \esp{He trabajado como vendedora.} « J'ai travaillé comme vendeuse. »
\item N'oubliez pas le participe passé irrégulier : \esp{he hecho} (j'ai fait), jamais \esp{he hacido}.
\item Regardez votre interlocuteur, parlez lentement, et n'oubliez pas l'intonation montante de la politesse finale : \esp{¡Muchas gracias!}
\end{itemize}
\end{attention}

\courssub{Étape 5 : la simulation d'entretien par binômes}

\begin{experience}[La entrevista de trabajo]
Travail par binômes, en deux manches pour que chacun joue les deux rôles.
\begin{enumerate}
\item \textbf{Préparation (5 min)} : chaque élève relit son CV et mémorise son autoportrait de l'étape 4.
\item \textbf{Présentation (1 min)} : l'élève A présente son CV à l'oral, sans le lire, en suivant le schéma identité $\rightarrow$ formación $\rightarrow$ experiencia $\rightarrow$ competencias $\rightarrow$ objetivo.
\item \textbf{Questions (2 min)} : l'élève B joue le rôle du \esp{jefe de personal} et pose au moins trois questions, par exemple : \esp{¿Por qué quiere usted este puesto?} « Pourquoi voulez-vous ce poste ? » ; \esp{¿Qué idiomas habla?} « Quelles langues parlez-vous ? » ; \esp{¿Cuál es su mejor cualidad?} « Quelle est votre meilleure qualité ? » ; \esp{¿Ha trabajado alguna vez con el público?} « Avez-vous déjà travaillé avec le public ? »
\item \textbf{Échange des rôles}, puis brève évaluation mutuelle à l'aide de la grille ci-dessous.
\end{enumerate}
\end{experience}

\begin{savaistu}[Le CV en Espagne, en Amérique latine et en Guinée équatoriale]
En \esp{España}, le CV inclut presque toujours une \textbf{photo} (format passeport, en haut à droite), la date de naissance et le \esp{DNI} (numéro d'identité) ou \esp{NIE} pour les étrangers. En \esp{Guinea Ecuatorial}, pays hispanophone d'Afrique centrale, les usages sont très proches de l'Espagne : photo, \esp{DNI} équato-guinéen, et mention de la \esp{nacionalidad}. En \esp{Amérique latine}, la photo devient facultative (souvent déconseillée pour éviter la discrimination), mais les \esp{datos personales} restent détaillés. Au Cameroun, pour un CV destiné à une entreprise hispanophone ou à une ONG internationale, suivez le modèle \textbf{espagnol} (photo + \esp{nacionalidad} + \esp{fecha de nacimiento}) : c'est ce que attendent les recruteurs hispanophones.
\end{savaistu}

\begin{exemplebox}[Réponses modèles aux questions de l'entretien]
\esp{Quiero este puesto porque me gusta el turismo y hablo tres idiomas.} « Je veux ce poste parce que j'aime le tourisme et je parle trois langues. »

\esp{Mi mejor cualidad es la puntualidad: soy la más puntual de mi clase.} « Ma meilleure qualité est la ponctualité : je suis la plus ponctuelle de ma classe. »

\esp{Sí, he trabajado con el público en el mercado de Mokolo.} « Oui, j'ai travaillé avec le public au marché Mokolo. »
\end{exemplebox}

\courssub{Relecture croisée : la grille d'évaluation}

Après la rédaction et la simulation, échangez votre CV avec votre voisin et évaluez-le selon la grille suivante. Chaque critère vaut un point ; visez le maximum.

\begin{center}
\begin{tabularx}{\linewidth}{|X|c|}
\hline
\rowcolor{popblueL} \textbf{Critère} & \textbf{Point} \\
\hline
Les cinq rubriques sont présentes et complètes & /1 \\
\hline
Ordre antichronologique respecté (formación, experiencia) & /1 \\
\hline
Accords corrects (\esp{nacida}, \esp{trabajadora}, \esp{puntual}) & /1 \\
\hline
Accents et orthographe (\esp{bilingüe}, \esp{informática}, \esp{Yaundé}) & /1 \\
\hline
Vocabulaire des rubriques correctement utilisé & /1 \\
\hline
Au moins un comparatif correct à l'oral (\esp{mejor que}, \esp{tan… como}) & /1 \\
\hline
Au moins un superlatif correct à l'oral (\esp{la más… de}, \esp{que más}) & /1 \\
\hline
Présentation orale fluide, sans lecture, avec \esp{he estudiado / he trabajado} & /1 \\
\hline
\end{tabularx}
\end{center}

\begin{exercice}[Bilan de la leçon]
\begin{enumerate}
\item Traduisez en espagnol : « Je parle espagnol mieux que l'année dernière, mais mon anglais est moins bon que mon espagnol. »
\item Complétez avec la forme correcte : \esp{Soy la persona más \ldots{} (indicado) para este puesto.}
\item Rédigez les quatre phrases de votre autoportrait oral en utilisant au moins un comparatif et un superlatif.
\end{enumerate}
\end{exercice}

\begin{corrige}[Bilan de la leçon]
\textbf{1.} \esp{Hablo español mejor que el año pasado, pero mi inglés es peor que mi español.} — On utilise les comparatifs irréguliers \esp{mejor que} (« mieux que ») et \esp{peor que} (« pire/moins bon que ») ; jamais \esp{más bueno} ni \esp{más malo} dans ce sens.

\textbf{2.} \esp{Soy la persona más indicada para este puesto.} — Superlatif relatif : \esp{más} + adjectif, avec accord de l'adjectif (\esp{indicada} au féminin si la locutrice est une femme, \esp{indicado} au masculin).

\textbf{3.} Exemple de production attendue : \esp{Me llamo Laure y tengo diecisiete años. Hablo español mejor que el año pasado. El español es la asignatura que más me gusta. Soy la alumna más puntual de mi clase y quiero trabajar como guía turística.} — Vérifier : présent (\esp{hablo}, \esp{soy}), comparatif irrégulier (\esp{mejor que}), superlatif relatif (\esp{la más puntual de}), accord des adjectifs.
\end{corrige}

\begin{aretenir}[L'essentiel de la production guidée]
\begin{itemize}
\item On prépare toujours ses idées \textbf{en français} (fiche d'identité) avant de transposer dans les rubriques espagnoles.
\item Le CV espagnol est bref, nominal, en ordre antichronologique ; les langues et les mois sans majuscule.
\item À l'oral, on ne lit pas : on suit le schéma identité $\rightarrow$ formación $\rightarrow$ experiencia $\rightarrow$ competencias $\rightarrow$ objetivo, au présent et au pretérito perfecto (\esp{he estudiado}, \esp{he trabajado}).
\item Les comparatifs (\esp{mejor que}, \esp{tan… como}) et superlatifs (\esp{la más… de}, \esp{que más}) valorisent l'autoportrait : \esp{Soy la persona más indicada para este puesto.}
\end{itemize}
\end{aretenir}

\newpage

\coursec{Activité d'intégration}

\coursec{Leçon E17 : Traduction — Comparatifs et superlatifs ; Redacción : El currículum vítae}

\courssub{Objectifs de la leçon}
Cette leçon de 2 heures vise deux compétences majeures du programme de Terminale A (Module 4 : Activités économiques et monde du travail) :
\begin{itemize}
\item \textbf{Traduire} (thème et version) des énoncés mobilisant les \cle{comparatifs} (supériorité, infériorité, égalité, irréguliers) et les \cle{superlatifs} (relatif, absolu), en évitant les pièges classiques des francophones (\esp{más mejor}, omission de l'article, accord de \esp{tanto}).
\item \textbf{Rédiger} un \cle{currículum vítae} complet en espagnol, conforme aux standards hispanophones (rubriques, lexique, format), et le défendre oralement lors d'un entretien simulé.
\end{itemize}
À l'issue de la leçon, l'élève doit être capable de décrire son parcours, de comparer ses atouts à un profil de poste et de produire un document professionnel utilisable hors de la classe.

\courssub{Situation déclenchante : La Feria de Empleo}
La ville de Douala accueille, au lycée bilingue, une \esp{Feria de Empleo} organisée par l'entreprise agroalimentaire hispano-camerounaise \emph{Cacao del Golfo S.A.}, qui exporte le cacao camerounais vers l'Espagne et l'Amérique latine. L'entreprise recrute de jeunes diplômés hispanophones pour des stages rémunérés (accueil clients, traduction commerciale, liaison Guinée équatoriale). La classe devient un salon de l'emploi : chaque élève alterne les rôles de \emph{candidat} et de \emph{responsable des ressources humaines} (\esp{recursos humanos}).

\begin{textebox}[Anuncio de empleo — Cacao del Golfo S.A.]
\textbf{¡TRABAJA CON NOSOTROS!}

Cacao del Golfo S.A., empresa camerunesa con sede en Douala, busca jóvenes dinámicos para prácticas remuneradas de tres meses.

\textbf{Perfil buscado :}
\begin{itemize}
\item Nivel de español igual o superior al nivel bachillerato.
\item Persona más organizada que la media, puntual y seria.
\item Experiencia previa: no es obligatoria, pero es una ventaja.
\item Disponibilidad inmediata.
\end{itemize}

\textbf{¿Cómo postular?}
Envía tu currículum vítae en español antes del viernes y prepárate para una entrevista de cinco minutos con nuestro equipo de recursos humanos.

\textit{El mejor candidato recibirá una carta de recomendación y una beca de formación.}
\end{textebox}

\textbf{Glossaire :} \esp{empresa} = entreprise ; \esp{con sede en} = dont le siège est à ; \esp{prácticas remuneradas} = stage rémunéré ; \esp{la media} = la moyenne ; \esp{una ventaja} = un avantage ; \esp{postular} = postuler ; \esp{una beca de formación} = une bourse de formation ; \esp{igual o superior a} = égal ou supérieur à.

\begin{experience}[Activité 1 : Analyse de l'annonce (10 min)]
Lisez l'annonce. Répondez par écrit en espagnol :
\begin{enumerate}
\item ¿Qué empresa ofrece el trabajo y dónde está su sede?
\item ¿Qué nivel de español se exige?
\item ¿Es la experiencia obligatoria?
\item Soulignez \textbf{un comparatif} et \textbf{un superlatif} dans le texte. Précisez pour chacun : type (régulier/irrégulier) et degré (supériorité/relatif).
\end{enumerate}
\textit{Corrigé attendu :} 1. Cacao del Golfo S.A., Douala. 2. Nivel bachillerato. 3. No, es una ventaja. 4. Comparatif : \esp{más organizada que} (régulier, supériorité). Superlatif : \esp{el mejor candidato} (irrégulier \esp{mejor}, relatif).
\end{experience}

\courssub{Observation : la comparaison en contexte}
Avant d'énoncer les règles, observons comment l'espagnol structure la comparaison dans des phrases authentiques issues du monde professionnel.

\begin{exemplebox}[Exemples authentiques]
\begin{itemize}
\item \esp{Mi salario es \textbf{más alto que} el de mi compañero.} (Supériorité)
\item \esp{Esta tarea es \textbf{menos difícil que} la anterior.} (Infériorité)
\item \esp{Trabajas \textbf{tan rápido como} yo.} (Égalité adverbe)
\item \esp{Tengo \textbf{tanta experiencia como} tú.} (Égalité nom — accord de \esp{tanta})
\item \esp{Eres \textbf{mejor} candidato que Juan.} (Irrégulier \esp{mejor})
\item \esp{Es \textbf{el más competente} del equipo.} (Superlatif relatif)
\item \esp{Estoy \textbf{contentísimo} con la oferta.} (Superlatif absolu en \esp{-ísimo})
\end{itemize}
\end{exemplebox}

\courssub{Les comparatifs : règles et emplois}

\begin{propriete}[Comparatifs réguliers]
\textbf{Forme} : \esp{más / menos / tan} + \emph{adjectif / adverbe / nom} + \esp{que}.
\begin{center}
\begin{tabularx}{\linewidth}{|X|X|X|}
\hline
\rowcolor{popblueL} \textbf{Type} & \textbf{Structure} & \textbf{Exemple + Traduction} \\
\hline
Supériorité & \esp{más + adj./adv. + que} & \esp{Este puesto es más interesante que el otro.} (Ce poste est plus intéressant que l'autre.) \\
\hline
Infériorité & \esp{menos + adj./adv. + que} & \esp{Gano menos dinero que mi jefe.} (Je gagne moins d'argent que mon patron.) \\
\hline
Égalité (adj./adv.) & \esp{tan + adj./adv. + como} & \esp{Eres tan puntual como yo.} (Tu es aussi ponctuel que moi.) \\
\hline
Égalité (nom) & \esp{tanto(s)/tanta(s) + nom + como} & \esp{Tengo tantas ganas como tú.} (J'ai autant d'envie que toi.) \\
\hline
\end{tabularx}
\end{center}
\textbf{Règles clés :}
\begin{itemize}
\item \esp{que} introduit le second terme (jamais \esp{de}).
\item \esp{tan} s'emploie devant adjectif/adverbe (invariable) ; \esp{tanto} s'accorde en genre et nombre avec le nom qui suit (\esp{tanto trabajo}, \esp{tanta paciencia}, \esp{tanto tiempo}, \esp{tantas horas}).
\item Devant un chiffre, \esp{más/menos de} : \esp{Más de cien candidatos.} (Plus de cent candidats.)
\end{itemize}
\end{propriete}

\begin{propriete}[Comparatifs irréguliers]
Ces formes remplacent \esp{más bueno}, \esp{más malo}, \esp{más grande}, \esp{más pequeño} quand l'adjectif est utilisé \textbf{seul} (sans complément de comparaison explicite \esp{que...}) ou par usage lexicalisé.
\begin{center}
\begin{tabular}{|l|l|l|l|}
\hline
\rowcolor{poporangeL} \textbf{Adjectif} & \textbf{Comparatif} & \textbf{Sens principal} & \textbf{Exemple} \\
\hline
\esp{bueno} & \esp{mejor} & meilleur (qualité) & \esp{Tu CV es mejor que el mío.} \\
\hline
\esp{malo} & \esp{peor} & pire & \esp{La situación es peor de lo que creía.} \\
\hline
\esp{grande} & \esp{mayor} & plus âgé / plus grand (taille, importance) & \esp{Mi hermana es mayor que yo.} (âge) \\
\hline
\esp{pequeño} & \esp{menor} & plus jeune / plus petit & \esp{Es un problema menor.} (importance) \\
\hline
\end{tabular}
\end{center}
\end{propriete}

\begin{attention}[Piège \esp{mayor/menor} vs \esp{más grande/pequeño}]
\esp{Mayor/Menor} s'emploient pour l'\textbf{âge} (\esp{mi hermano mayor}) et la \textbf{hiérarchie/importance} (\esp{un problema mayor}). Pour la taille physique d'un objet, on garde \esp{más grande/pequeño} : \esp{Esta caja es más grande que aquella.}
\end{attention}

\begin{attention}[Erreurs interdites — \cle{Faux amis francophones}]
\begin{itemize}
\item \textbf{Jamais \esp{más mejor}, \esp{más peor}, \esp{más mayor}, \esp{más menor}.} Ces formes contiennent déjà l'idée de « plus ».
\item \textbf{Accord de \esp{tanto} :} \esp{Tengo tanto trabajo como tú} (masc. sg.), \esp{Tanta prisa como yo} (fém. sg.), \esp{Tantos amigos como él} (masc. pl.), \esp{Tantas ganas como nosotras} (fém. pl.).
\item \textbf{\esp{Que} vs \esp{de} :} Après comparatif \rightarrow \esp{que}. Après superlatif relatif + groupe nominal \rightarrow \esp{de} : \esp{El mejor de la clase.}
\item \textbf{\esp{Ser} vs \esp{Estar} dans la comparaison :} \esp{Es más listo que yo} (qualité intrinsèque) vs \esp{Está más cansado que yo} (état passager).
\end{itemize}
\end{attention}

\courssub{Les superlatifs : relatif et absolu}

\begin{propriete}[Superlatif relatif — Le plus / le moins \emph{dans un groupe}]
\textbf{Forme} : \textbf{Article défini} (\esp{el/la/los/las}) + \esp{más / menos} + \emph{adjectif} + \esp{de} (groupe de référence).
\begin{itemize}
\item \esp{Es \textbf{la más preparada} \textbf{de} todos los candidatos.} (C'est la plus préparée de tous les candidats.)
\item \esp{Son \textbf{los menos experimentados} \textbf{de} la empresa.} (Ce sont les moins expérimentés de l'entreprise.)
\item Irréguliers : \esp{el mejor}, \esp{la peor}, \esp{los mayores}, \esp{las menores} (article obligatoire).
\end{itemize}
\end{propriete}

\begin{attention}[L'article est obligatoire] En français « le plus fort » = \esp{el más fuerte}. Dire \esp{más fuerte} seul signifie « plus fort » (comparatif), pas « le plus fort ».
\end{attention}


\begin{propriete}[Superlatif absolu — Très / Extrêmement (sans comparaison)]
Trois structures principales :
\begin{enumerate}
\item \esp{Muy + adjectif} (neutre, courant) : \esp{Estoy muy interesado.}
\item \esp{Adjectif + -ísimo / -ísima / -ísimos / -ísimas} (expressif, emphatique) : \esp{La tarea es facilísima.} (Accent sur la voyelle tonique de l'adjectif : \esp{fácil} $\rightarrow$ \esp{facilísimo} ; \esp{contento} $\rightarrow$ \esp{contentísimo}.)
\item \esp{Sumamente / extremadamente / altamente + adjectif} (formel, écrit) : \esp{Es sumamente importante.}
\end{enumerate}
\end{propriete}

\begin{exemplebox}[Variantes morphologiques -ísimo]
\esp{Grande} $\rightarrow$ \esp{grandísimo} (perte du \esp{d}) ; \esp{Joven} $\rightarrow$ \esp{jovenísimo} ; \esp{Feliz} $\rightarrow$ \esp{felicísimo} (\esp{z} $\rightarrow$ \esp{c}) ; \esp{Leal} $\rightarrow$ \esp{lealísimo}.
\end{exemplebox}


\courssub{Traduction guidée : thème et version}
Appliquez les règles ci-dessus. Traduisez en espagnol (thème) ou en français (version).

\begin{exercice}[Thème : Comparatifs et superlatifs]
Traduisez en espagnol.
\begin{enumerate}
\item Ce stage est plus intéressant que mon dernier emploi.
\item Je parle espagnol aussi bien que mon professeur.
\item Elle a moins d'expérience que moi, mais elle est plus motivée.
\item C'est le meilleur candidat de toute la promotion. (promoción)
\item Nous sommes très pressés (superlatif absolu emphatique).
\item Ce problème est moins grave qu'il n'y paraît. (que parece)
\item Tu as autant de diplômes que moi.
\item C'est la pire situation possible.
\item Mon frère aîné (mayor) travaille au Cameroun.
\item L'entretien a été extrêmement difficile. (sumamente)
\end{enumerate}
\end{exercice}

\begin{corrige}[Exercice Thème]
\begin{enumerate}
\item \esp{Este prácticas es más interesante que mi último empleo.} (\esp{prácticas} fém. sg. ; \esp{empleo} masc.)
\item \esp{Hablo español tan bien como mi profesor.}
\item \esp{Tiene menos experiencia que yo, pero está más motivada.} (\esp{estar} pour état passager \esp{motivada}).
\item \esp{Es el mejor candidato de toda la promoción.}
\item \esp{Estamos apuradísimos.} (ou \esp{muy apurados}).
\item \esp{Este problema es menos grave de lo que parece.} (\esp{de lo que} après comparatif + infinitif/proposition).
\item \esp{Tienes tantos títulos como yo.}
\item \esp{Es la peor situación posible.}
\item \esp{Mi hermano mayor trabaja en Camerún.}
\item \esp{La entrevista fue sumamente difícil.}
\end{enumerate}
\end{corrige}

\begin{exercice}[Version : Du espagnol au français]
Traduisez en français.
\begin{enumerate}
\item \esp{Mi sueldo es menor que el de mi jefe.}
\item \esp{Esta empresa es la más grande del sector.}
\item \esp{No eres tan responsable como tu hermana.}
\item \esp{Tengo tantas ganas de trabajar como tú.}
\item \esp{El candidato es malísimo para el puesto.}
\item \esp{Es mejor prevenir que curar.} (Proverbe)
\end{enumerate}
\end{exercice}

\begin{corrige}[Exercice Version]
\begin{enumerate}
\item Mon salaire est inférieur à celui de mon patron. (\esp{menor} = plus petit / inférieur).
\item Cette entreprise est la plus grande du secteur.
\item Tu n'es pas aussi responsable que ta sœur.
\item J'ai autant envie de travailler que toi.
\item Le candidat est très mauvais / nul pour le poste. (\esp{malísimo} = superlatif absolu de \esp{malo}).
\item Mieux vaut prévenir que guérir. (Litt. : Il est meilleur de prévenir que de guérir).
\end{enumerate}
\end{corrige}

\courssub{Le currículum vítae en espagnol : structure et lexique}

\begin{methode}[Rédiger un CV hispanophone (Pas à pas)]
\begin{enumerate}
\item \textbf{En-tête (\esp{Datos personales}) :} Nom complet (Prénom + 2 noms de famille si usage local), date de naissance (JJ/MM/AAAA), adresse, téléphone, email, nationalité. \textit{Pas de photo obligatoire en Espagne, courante en Amérique Latine / Guinée équatoriale.}
\item \textbf{Formation (\esp{Formación académica}) :} Ordre antéchronologique (récent $\rightarrow$ ancien). Diplôme, établissement, ville, année, mention. Ex: \esp{2023-2025: Bachillerato en Ciencias, Lycée de Yaoundé, Mención Bien.}
\item \textbf{Expérience (\esp{Experiencia profesional}) :} Ordre antéchronologique. Poste, entreprise, lieu, dates, 2-3 puces \textbf{actions/résultats} (verbes d'action : \esp{gestioné}, \esp{atendí}, \esp{traduje}, \esp{colaboré}).
\item \textbf{Langues (\esp{Idiomas}) :} Niveau CECRL ou descriptif : \esp{Lengua materna}, \esp{Nivel avanzado (C1/C2)}, \esp{Nivel intermedio (B1/B2)}, \esp{Nivel básico (A1/A2)}. Préciser certifications (DELE, SIELE).
\item \textbf{Compétences transverses (\esp{Otros datos de interés / Habilidades}) :} Informatique (\esp{Office, Python, Canva}), Permis (\esp{Carnet de conducir B}), Bénévolat, Sports, Centres d'intérêt.
\end{enumerate}
\end{methode}

\begin{definition}[Lexique spécialisé CV / Entretien]
\begin{center}
\begin{tabularx}{\linewidth}{|X|X|X|}
\hline
\rowcolor{popgreenL} \textbf{Français} & \textbf{Español} & \textbf{Contexte / Exemple} \\
\hline
Le poste / L'emploi & \esp{El puesto de trabajo / El empleo} & \esp{Solicito el puesto de traductor.} \\
Le salaire & \esp{El sueldo / El salario} & \esp{Un sueldo competitivo.} \\
L'entretien & \esp{La entrevista (de trabajo)} & \esp{Tengo una entrevista mañana.} \\
Embaucher / Recruter & \esp{Contratar / Reclutar} & \esp{Quieren contratar a dos becarios.} \\
Le stage & \esp{Las prácticas (profesionales)} & \esp{Hice prácticas en un banco.} \\
La lettre de recommandation & \esp{La carta de recomendación} & \esp{Adjunto carta de recomendación.} \\
Les connaissances & \esp{Los conocimientos} & \esp{Conocimientos de contabilidad.} \\
Les compétences (savoir-faire) & \esp{Las habilidades / Las competencias} & \esp{Habilidades comunicativas.} \\
Travailleur / -euse & \esp{Trabajador / Trabajadora} & \esp{Es una persona muy trabajadora.} \\
Responsable & \esp{Responsable} & \esp{Cargo de responsabilidad.} \\
Ponctuel / -le & \esp{Puntual} & \esp{Llegar puntual es vital.} \\
Disponibilité & \esp{Disponibilidad} & \esp{Disponibilidad inmediata.} \\
\hline
\end{tabularx}
\end{center}
\end{definition}

\begin{savaistu}[Le CV au Cameroun et dans l'espace hispanophone]
Au Cameroun, le CV « français » (photo, état civil détaillé) est la norme. En Espagne, la loi interdit de demander la photo, l'âge, l'état civil ou la nationalité pour lutter contre la discrimination ; le CV est anonyme (\esp{CV ciego}). Au Mexique, en Colombie ou en Guinée équatoriale, la photo et les données personnelles restent courantes. \textbf{Conseil :} Adaptez votre CV au pays de l'entreprise visée. Pour \emph{Cacao del Golfo S.A.} (siège Douala, capital hispano), la formule « complète » (avec photo, âge, nationalité) est attendue.
\end{savaistu}

\courssub{Production intégrée : Mon CV et mon entretien}

\begin{experience}[Tarea 2 — Rédaction de votre CV (25 min)]
Rédigez votre \esp{Currículum Vítae} réel ou fictif (élève Terminale A) sur une page A4. Respectez l'ordre des rubriques, la sobriété du style (noms, infinitifs ou verbes passé, pas de phrases complètes), et l'orthographe des accents. Utilisez le lexique de la leçon.
\end{experience}

\begin{experience}[Tarea 3 — Préparation du « Pitch » 1 minute (10 min)]
Préparez une présentation orale où vous vous « vendez » face au recruteur. \textbf{Contraintes obligatoires :}
\begin{itemize}
\item 2 comparatifs minimum (réguliers ou irréguliers).
\item 1 superlatif minimum (relatif ou absolu).
\item Vocabulaire du monde du travail.
\end{itemize}
\textit{Amorces :} \esp{« Soy más... que la media... » ; « Mi mayor fortaleza es... » ; « Tengo tanta ilusión como... » ; « Me considero el/la mejor candidato/a porque... »}
\end{experience}

\begin{experience}[Tarea 4 — Simulation d'entretien (20 min)]
Par binômes (A = RH, B = Candidat). Échangez les rôles.
\textbf{RH (script obligatoire) :}
\begin{enumerate}
\item \esp{¿Qué idiomas habla usted y a qué nivel?}
\item \esp{¿Tiene experiencia profesional o ha hecho prácticas?}
\item \esp{¿Por qué es usted el mejor candidato para Cacao del Golfo?}
\end{enumerate}
\textbf{Candidat :} Répondre en intégrant \textbf{au moins un comparatif ou superlatif par réponse}. Ex: \esp{Hablo francés como lengua materna y español mejor que la mayoría de mis compañeros.}
\end{experience}

\begin{exoresolu}[Modèle commenté : CV et Pitch d'Amina Mbarga (Terminale A, Yaoundé)]
Rédigez le CV et la présentation orale d'Amina Mbarga, candidate au stage de Cacao del Golfo S.A.
\tcblower
\textbf{Currículum Vítae de Amina Mbarga}

\esp{\textbf{Datos personales}\\
Nombre y apellidos: Amina Mbarga Ngo Bell\\
Fecha de nacimiento: 15 de marzo de 2007\\
Dirección: Barrio Nkolbisson, Yaoundé, Camerún\\
Teléfono: 690 00 00 00 \textbar\ Correo electrónico: amina.mbarga@email.cm\\
Nacionalidad: Camerunesa}

\esp{\textbf{Formación académica}\\
2023--2025: Bachillerato (Terminale A), Lycée Bilingue de Yaoundé. Mención: Bien.\\
2023: Probatoire (Brevet), Série A, Mención Bien.}

\esp{\textbf{Experiencia profesional}\\
Verano 2024: Vendedora, Mercado Mokolo (Yaoundé). Atención al cliente, gestión de caja, inventario.\\
2024: Intérprete voluntaria, Visita delegación estudiantil Guinea Ecuatorial. Traducción consecutiva español-francés.}

\esp{\textbf{Idiomas}\\
Francés: Lengua materna.\\
Español: Nivel avanzado (C1) -- DELE B2 obtenido 2024.\\
Inglés: Nivel intermedio (B1).\\
Ewondo: Lengua materna.}

\esp{\textbf{Habilidades informáticas}\\
Microsoft Office (Word, Excel, PowerPoint) -- Nivel usuario avanzado.\\
Google Workspace.}

\esp{\textbf{Otros datos de interés}\\
Deportes: Balonmano (capitana equipo lycée).\\
Disponibilidad: Inmediata. Carnet de conducir tipo B.}

\textbf{Pitch oral (1 minute)}

\esp{Buenos días. Me llamo Amina Mbarga y vengo a ofrecerles mi entusiasmo y mi rigor. Soy \textbf{más organizada que la mayoría} de mis compañeros: planifico mi semana cada domingo y nunca entrego tarde. Hablo español \textbf{tan fluidamente como} un estudiante de inmersión, y tengo \textbf{tanta ilusión por aprender como experiencia me falta}, ¡pero me adapto \textbf{rapidísimo}! Mi \textbf{mayor fortaleza} es el trato humano: en el mercado Mokolo fui \textbf{la vendedora mejor valorada} por los clientes. Por todo esto, me considero \textbf{la mejor candidata} para representar a Cacao del Golfo. Muchas gracias.}

\textbf{Commentaire linguistique détaillé.}
\begin{itemize}
\item \textbf{CV :} Respect des rubriques canoniques. Style télégraphique (noms, participes passés \esp{gestioné} implicites). \esp{Mención Bien} (maj. à \esp{Mención} seulement). Dates JJ/MM/AAAA. \esp{Guinea Ecuatorial} (pas \esp{Guinée équatoriale}).
\item \textbf{Pitch :} 1. \esp{más organizada que la mayoría} (comp. sup. régulier). 2. \esp{tan fluidamente como} (comp. égalité adv.). 3. \esp{tanta ilusión como experiencia me falta} (comp. égalité nom \esp{tanta} fém. sg. + structure elliptique rhétorique correcte : « autant d'envie que l'expérience qui me manque »). 4. \esp{rapidísimo} (superlatif absolu \esp{-ísimo}). 5. \esp{mayor fortaleza} (irrégulier \esp{mayor} = principal/majeure, pas âge). 6. \esp{la vendedora mejor valorada} (superlatif relatif irrégulier \esp{mejor} + article \esp{la} + participe). 7. \esp{la mejor candidata} (superlatif relatif irrégulier + article). \textbf{Zero faute.}
\end{itemize}
\end{exoresolu}

\begin{experience}[Tarea 5 — Co-évaluation et exposition (15 min)]
Chaque binôme note le CV et l'entretien de son partenaire sur 20 selon la grille :
\begin{center}
\begin{tabular}{|l|c|}
\hline
\rowcolor{poppurpleL} \textbf{Critère} & \textbf{Points} \\
\hline
Correction comparatifs / superlatifs (emploi, accord, article) & /4 \\
\hline
Structure CV complète + lexique professionnel & /4 \\
\hline
Richesse lexicale (variété, justesse) & /4 \\
\hline
Aisance orale (prononciation, fluidité, interaction) & /4 \\
\hline
Présentation / Ponctualité / Tenue professionnelle & /4 \\
\hline
\textbf{Total} & \textbf{/20} \\
\hline
\end{tabular}
\end{center}
Les 3 meilleurs CV et les 3 meilleures prestations orales sont affichés au « Mur de l'Emploi » de la classe.
\end{experience}

\courssub{Bilan de la leçon}
Cette \esp{Feria de Empleo} a ancré les acquis grammaticaux dans une tâche finale authentique (\emph{task-based learning}).
\begin{itemize}
\item \textbf{Grammaire :} Maîtrise des 4 comparatifs (\esp{más/menos/tan/tanto}), des 4 irréguliers (\esp{mejor/peor/mayor/menor}), du superlatif relatif (\esp{el más... de}, \esp{el mejor}) et absolu (\esp{muy/-ísimo/sumamente}). Évitement des pièges (\esp{más mejor}, accord \esp{tanto}, article relatif).
\item \textbf{Écriture professionnelle :} Production d'un CV normé, utilisable pour une vraie candidature au Cameroun ou à l'international.
\item \textbf{Oral :} Prise de parole en interaction (entretien), gestion du stress, argumentation comparative.
\item \textbf{Culture :} Conscience de la valeur marchande de l'espagnol au Cameroun (voisin de la Guinée équatoriale, partenaire Espagne/Amérique latine). L'espagnol n'est pas une « option », c'est un \esp{avantage competitivo}.
\end{itemize}

\begin{savaistu}[Prolongement : Le marché du travail hispanophone]
L'Espagne est le 2e investisseur étranger au Cameroun (après la France). La Guinée équatoriale (pays hispanophone frontalier) recrute massivement dans le pétrole, le bois, les infrastructures. Maîtriser le CV et l'entretien en espagnol, c'est s'ouvrir le marché de l'emploi de \textbf{21 pays} et de \textbf{500 millions} de locuteurs natifs. \esp{¡Tu futuro empieza hoy!}
\end{savaistu}

\newpage

\coursec{Méthodes et exercices résolus}

\courssub{Observation : La comparaison en contexte}

\begin{textebox}[Forum emploi : « ¿Qué valoran más? » -- Extraits authentiques]
\esp{-- Usuario1: « Busco trabajo en Málaga. ¿Alguien sabe si los salarios son \textbf{más altos} que en Madrid? »} \\
\esp{-- Usuario2: « En general, \textbf{menos altos}, pero el coste de vida es \textbf{menor}. »} \\
\esp{-- Usuario3: « Para mí, lo \textbf{mejor} es la calidad de vida. Mi \textbf{mayor} prioridad: tiempo libre. »} \\
\esp{-- Usuario4: « Tengo \textbf{tanta experiencia como} tú, pero me ofrecen \textbf{peores} condiciones. »} \\
\esp{-- Usuario5: « ¡Ojo! Un contrato \textbf{tan precario como} ese \textbf{no vale la pena}. »}
\end{textebox}

\begin{definition}[Les trois degrés de la comparaison]
En espagnol comme en français, la comparaison s'exprime selon trois degrés :
\begin{enumerate}
    \item \textbf{Comparatif de supériorité} : \emph{plus... que} $\rightarrow$ \esp{más... que}.
    \item \textbf{Comparatif d'infériorité} : \emph{moins... que} $\rightarrow$ \esp{menos... que}.
    \item \textbf{Comparatif d'égalité} :
    \begin{itemize}
        \item Qualité (adjectif/adverbe) : \emph{aussi... que} $\rightarrow$ \esp{tan... como}.
        \item Quantité (nom) : \emph{autant de... que} $\rightarrow$ \esp{tanto/tanta/tantos/tantas... como} (\textbf{accord}).
    \end{itemize}
\end{enumerate}
\cle{Repère} : Les mots-clés \esp{más, menos, tan, tanto, mejor, peor, mayor, menor, -ísimo} balisent la phrase.
\end{definition}

\begin{savaistu}[La Guinée équatoriale : unique pays hispanophone d'Afrique]
La Guinée équatoriale (\esp{Guinea Ecuatorial}) a l'espagnol comme langue officielle (avec le français et le portugais) depuis 1844. Sa capitale, \esp{Malabo} (sur l'île de Bioko), et \esp{Bata} (continent) sont des pôles économiques (pétrole, bois, cacao) où la maîtrise de l'espagnol professionnel est un atout majeur pour les jeunes Camerounais de la zone CEMAC. Le \esp{DELE} (Diploma de Español como Lengua Extranjera) y est organisé par l'Instituto Cervantes de Malabo et Bata.
\end{savaistu}

\courssub{Méthode 1 : Traduire les comparatifs de supériorité, d'infériorité et d'égalité}

\begin{methode}[Comparatifs réguliers : repérage, accord et traduction]
\textbf{Objectif :} Passer du français à l'espagnol (thème) et de l'espagnol au français (version) sans calque.
\begin{enumerate}
    \item \textbf{Repérer le type de comparatif} dans la phrase française : supériorité (\emph{plus... que}), infériorité (\emph{moins... que}), égalité (\emph{aussi... que}, \emph{autant... que}).
    \item \textbf{Choisir la structure espagnole correspondante} :
    \begin{itemize}
        \item Supériorité : \esp{más + adjectif/adverbe + que}
        \item Infériorité : \esp{menos + adjectif/adverbe + que}
        \item Égalité (qualité) : \esp{tan + adjectif/adverbe + como}
        \item Égalité (quantité) : \esp{tanto/tanta/tantos/tantas + nom + como} (\textbf{accord obligatoire} en genre et nombre avec le nom).
    \end{itemize}
    \item \textbf{Vérifier l'accord de l'adjectif} (si comparatif d'adjectif) avec le sujet de la phrase espagnole.
    \item \textbf{Traduire \emph{que} par \esp{que} (conjonction) ou \esp{de} (chiffre/quantité) :} \esp{Más de cien personas} (Plus de cent personnes).
    \item \textbf{Relire} en remplaçant mentalement \esp{más/menos/tan} par \emph{plus/moins/aussi} pour valider le sens.
\end{enumerate}
\end{methode}

\begin{exoresolu}[Thème et version : comparatifs réguliers]
\textbf{Énoncé :}
\begin{enumerate}
    \item \textbf{Thème :} Traduisez en espagnol.
    \begin{enumerate}
        \item Yaoundé est plus peuplée que Douala.
        \item Ce travail est moins difficile que je ne le pensais.
        \item Elle parle aussi couramment l'espagnol que son frère.
        \item Nous avons autant d'opportunités qu'eux.
        \item Le salaire est supérieur à 200 000 francs CFA. (Indice : \esp{superior a} / \esp{más de})
    \end{enumerate}
    \item \textbf{Version :} Traduisez en français.
    \begin{enumerate}
        \item \esp{Los jóvenes cameruneses tienen menos oportunidades que los europeos.}
        \item \esp{Este contrato es tan ventajoso como el anterior.}
        \item \esp{En Guinea Ecuatorial hay tantos hispanohablantes como en España.}
        \item \esp{Mi jefe es más exigente que el tuyo.}
    \end{enumerate}
\end{enumerate}

\tcblower

\textbf{Solution détaillée :}

\textbf{1. Thème}
\begin{enumerate}
    \item \textbf{Repérage :} Supériorité (\emph{plus... que}) sur adjectif \emph{peuplée}. Sujet : \emph{Yaoundé} (féminin singulier).
    \textbf{Structure :} \esp{más + adj (fém. sg.) + que}.
    \textbf{Réponse :} \esp{Yaoundé está más poblada que Douala.} (Note : \esp{estar} pour état/résultat démographique, ou \esp{ser} pour caractéristique intrinsèque ; \esp{estar más poblada} est fréquent pour la densité actuelle).
    \item \textbf{Repérage :} Infériorité (\emph{moins... que}) sur adjectif \emph{difficile}. Sujet : \emph{este trabajo} (masc. sg.). \emph{Que je ne le pensais} $\rightarrow$ \esp{de lo que yo pensaba} (néophronème \emph{lo}).
    \textbf{Réponse :} \esp{Este trabajo es menos difícil de lo que yo pensaba.}
    \item \textbf{Repérage :} Égalité qualité (\emph{aussi... que}) sur adverbe \emph{couramment} (\esp{con fluidez} / \esp{bien}). Sujet : \emph{ella}.
    \textbf{Réponse :} \esp{Ella habla tan fluidamente el español como su hermano.} (ou \esp{tan bien como}).
    \item \textbf{Repérage :} Égalité quantité (\emph{autant... que}) sur nom \emph{opportunités} (\esp{oportunidades}, fém. pl.).
    \textbf{Accord :} \esp{tantas oportunidades}.
    \textbf{Réponse :} \esp{Tenemos tantas oportunidades como ellos.}
    \item \textbf{Repérage :} Supériorité sur chiffre (\emph{supérieur à / plus de}).
    \textbf{Règle :} Devant chiffre $\rightarrow$ \esp{más de} (pas \esp{que}).
    \textbf{Réponse :} \esp{El salario es superior a 200.000 francos CFA.} (ou \esp{más de 200.000 francos CFA}).
\end{enumerate}

\textbf{2. Version}
\begin{enumerate}
    \item \esp{Los jóvenes cameruneses tienen menos oportunidades que los europeos.}
    $\rightarrow$ \emph{Les jeunes Camerounais ont moins d'opportunités que les Européens.} (Infériorité quantité : \esp{menos oportunidades}).
    \item \esp{Este contrato es tan ventajoso como el anterior.}
    $\rightarrow$ \emph{Ce contrat est aussi avantageux que le précédent.} (Égalité qualité : \esp{tan + adj + como}). Accord \esp{ventajoso} (masc. sg.) avec \esp{contrato}.
    \item \esp{En Guinea Ecuatorial hay tantos hispanohablantes como en España.}
    $\rightarrow$ \emph{En Guinée équatoriale, il y a autant d'hispanophones qu'en Espagne.} (Égalité quantité : \esp{tantos hispanohablantes}, masc. pl.). \esp{Hay} = \emph{il y a} (impersonnel).
    \item \esp{Mi jefe es más exigente que el tuyo.}
    $\rightarrow$ \emph{Mon patron est plus exigeant que le tien.} (Supériorité : \esp{más exigente}). \esp{El tuyo} = \emph{le tien} (pronom possessif, masc. sg. pour \esp{jefe}).
\end{enumerate}

\textbf{Conclusion :} La clé est l'identification du mot pivot (\emph{plus, moins, aussi, autant}) et l'accord rigoureux de \esp{tanto/tanta/tantos/tantas} et des adjectifs.
\end{exoresolu}

\courssub{Méthode 2 : Traduire les comparatifs irréguliers (mejor, peor, mayor, menor)}

\begin{methode}[Comparatifs irréguliers : choix lexical et pièges à éviter]
\textbf{Objectif :} Maîtriser les quatre formes irrégulières et leurs nuances d'emploi (âge, taille, qualité, quantité).
\begin{enumerate}
    \item \textbf{Identifier l'adjectif français source} : \emph{meilleur, pire, plus grand/âgé, plus petit/jeune}.
    \item \textbf{Appliquer la correspondance irrégulière} (jamais \esp{más bueno}, \esp{más malo}, \esp{más grande} pour l'âge/hiérarchie) :
    \begin{center}
    \begin{tabularx}{\linewidth}{|X|X|X|}\hline
    \rowcolor{popblueL} Français & Espagnol irrégulier & Nuance / Contexte \\ \hline
    \emph{meilleur} & \esp{mejor} & Qualité, skill, santé (\esp{estar mejor}) \\ \hline
    \emph{pire} & \esp{peor} & Qualité, gravité, santé (\esp{estar peor}) \\ \hline
    \emph{plus grand / plus âgé / supérieur (hiérarchie)} & \esp{mayor} & Âge, rang, importance, taille (abstrait) \\ \hline
    \emph{plus petit / plus jeune / inférieur} & \esp{menor} & Âge, rang, importance, taille (abstrait) \\ \hline
    \end{tabularx}
    \end{center}
    \item \textbf{Distinction cruciale \esp{mayor/menor} vs \esp{más grande/más pequeño}} :
    \begin{itemize}
        \item \esp{Mayor/Menor} : Âge (\esp{mi hermano mayor}), hiérarchie (\esp{un cargo mayor}), importance (\esp{un problema mayor}). \textbf{Jamais} pour la taille physique d'un objet.
        \item \esp{Más grande/Más pequeño} : Taille physique, dimensions concrètes (\esp{una casa más grande}).
    \end{itemize}
    \item \textbf{Construire la phrase} : Sujet + \esp{ser/estar} + \esp{mejor/peor/mayor/menor} + \esp{que} + terme de comparaison.
    \item \textbf{Attention au faux ami :} \esp{Superior/Inferior} existent mais sont formels (hiérarchie, qualité) ; \esp{Mayor/Menor} sont plus courants pour l'âge et l'importance.
\end{enumerate}
\end{methode}

\begin{exoresolu}[Transformation et thème : comparatifs irréguliers]
\textbf{Énoncé :}
\begin{enumerate}
    \item \textbf{Transformez} les phrases suivantes en remplaçant l'adjectif entre parenthèses par le comparatif irrégulier correct.
    \begin{enumerate}
        \item Este café es \textbf{(bueno)} que el de ayer.
        \item Mi situación económica es \textbf{(mala)} que la tuya.
        \item Mi hermana es \textbf{(grande)} que yo. (âge)
        \item Este problema es \textbf{(pequeño)} que el anterior. (importance)
        \item Esta casa es \textbf{(grande)} que la vuestra. (taille physique)
    \end{enumerate}
    \item \textbf{Thème :} Traduisez en espagnol.
    \begin{enumerate}
        \item Mon frère aîné est médecin. (Indice : \esp{médico})
        \item C'est une erreur plus grave que la précédente.
        \item Il se sent mieux aujourd'hui.
        \item Le directeur est supérieur en grade à l'adjoint. (Indice : \esp{rango} / \esp{cargo})
    \end{enumerate}
\end{enumerate}

\tcblower

\textbf{Solution détaillée :}

\textbf{1. Transformation}
\begin{enumerate}
    \item \esp{Este café es \textbf{mejor} que el de ayer.} (\emph{bueno} $\rightarrow$ \esp{mejor}). \esp{El de ayer} = \emph{celui d'hier} (pronom neutre \esp{lo} + \esp{de} évité ici par \esp{el de} pour masc. sg.).
    \item \esp{Mi situación económica es \textbf{peor} que la tuya.} (\emph{mala} $\rightarrow$ \esp{peor}). \esp{La tuya} (fém. sg. pour \esp{situación}).
    \item \esp{Mi hermana es \textbf{mayor} que yo.} (\emph{grande} = plus âgée $\rightarrow$ \esp{mayor}). \textbf{Piège :} \esp{más grande} signifierait "plus grande physiquement".
    \item \esp{Este problema es \textbf{menor} que el anterior.} (\emph{petit} = moins important $\rightarrow$ \esp{menor}). \esp{El anterior} (masc. sg.).
    \item \esp{Esta casa es \textbf{más grande} que la vuestra.} (\emph{grande} = taille physique $\rightarrow$ \esp{más grande}). \textbf{Règle :} \esp{mayor} interdit pour dimensions physiques.
\end{enumerate}

\textbf{2. Thème}
\begin{enumerate}
    \item \textbf{Frère aîné} $\rightarrow$ \esp{hermano mayor}. \textbf{Réponse :} \esp{Mi hermano mayor es médico.}
    \item \textbf{Erreur plus grave} $\rightarrow$ \emph{pire} / \emph{plus grave} $\rightarrow$ \esp{peor} (ou \esp{más grave}, régulier admis pour \emph{grave}, mais \esp{peor} plus fort). \textbf{Réponse :} \esp{Es un error peor que el anterior.}
    \item \textbf{Se sent mieux} (santé/état) $\rightarrow$ \esp{estar mejor}. \textbf{Réponse :} \esp{Hoy se siente mejor.} (\esp{Sentirse} pronominal).
    \item \textbf{Supérieur en grade} (hiérarchie) $\rightarrow$ \esp{mayor rango} / \esp{cargo mayor}. \textbf{Réponse :} \esp{El director tiene un cargo mayor que el subdirector.} (ou \esp{es de mayor rango}).
\end{enumerate}

\textbf{Conclusion :} L'erreur fatale au Bac est d'écrire \esp{más bueno}, \esp{más malo}, \esp{más grande} (pour l'âge) ou \esp{más pequeño} (pour l'importance). Mémorisez le quatuor \cle{mejor, peor, mayor, menor} comme des blocs indivisibles.
\end{exoresolu}

\courssub{Méthode 3 : Traduire les superlatifs (relatif et absolu)}

\begin{methode}[Superlatifs : structure, accord et intensité]
\textbf{Objectif :} Traduire \emph{le plus / le moins} (relatif) et \emph{très / extrêmement / -issime} (absolu).
\begin{enumerate}
    \item \textbf{Superlatif Relatif} (\emph{le plus / le moins} + adjectif + \emph{de/dans}) :
    \textbf{Formule :} \esp{El/La/Los/Las + más/menos + adjectif (accordé) + de/en + lieu/groupe}.
    \begin{itemize}
        \item \esp{De} pour l'appartenance à un ensemble (\esp{el mejor de la clase}).
        \item \esp{En} pour un lieu physique (\esp{el más alto en el edificio}).
    \end{itemize}
    \item \textbf{Superlatif Absolu synthétique (suffixe \textbf{-ísimo})} :
    \textbf{Formule :} Adjectif (perte voyelle finale \emph{-o/-a}) + \esp{-ísimo/-ísima/-ísimos/-ísimas}.
    \begin{itemize}
        \item Accent obligatoire sur le \emph{í} : \esp{guapísimo, riquísima, facilísimos}.
        \item Adjectifs en \emph{-ble} : \esp{amable $\rightarrow$ amabilísimo} (perte \emph{e}). En \emph{-co/-ca/-go/-ga} : \esp{qu/gu} (\esp{rico $\rightarrow$ riquísimo, largo $\rightarrow$ larguísimo}).
    \end{itemize}
    \item \textbf{Superlatif Absolu analytique (adverbes d'intensité)} :
    \esp{Muy + adj} (standard), \esp{Sumamente / Extremadamente / Inmensamente + adj} (soutenu, écrit, formel).
    \item \textbf{Traduction version $\rightarrow$ français} :
    \begin{itemize}
        \item \esp{El más... de} $\rightarrow$ \emph{Le plus... de}.
        \item \esp{-ísimo} $\rightarrow$ \emph{Très...} ou \emph{...issime} (selon registre).
        \item \esp{Muy / Sumamente} $\rightarrow$ \emph{Très / Extrêmement}.
    \end{itemize}
\end{enumerate}
\end{methode}

\begin{exoresolu}[Thème/Version mixte : superlatifs dans un contexte professionnel]
\textbf{Énoncé :}
\begin{enumerate}
    \item \textbf{Thème :} Traduisez en espagnol.
    \begin{enumerate}
        \item C'est le candidat le plus compétent de tous.
        \item C'est la tâche la moins agréable du poste.
        \item Le directeur est extrêmement occupé aujourd'hui.
        \item Nous avons une responsabilité énorme. (Indice : \esp{enorme} $\rightarrow$ \esp{enormísimo}).
        \item C'est un problème très facile à résoudre.
    \end{enumerate}
    \item \textbf{Version :} Traduisez en français.
    \begin{enumerate}
        \item \esp{Es el empleado más puntual de la empresa.}
        \item \esp{La entrevista fue dificilísima.}
        \item \esp{Los requisitos son sumamente estrictos.}
        \item \esp{Esta es la oferta menos interesante del mercado.}
    \end{enumerate}
\end{enumerate}

\tcblower

\textbf{Solution détaillée :}

\textbf{1. Thème}
\begin{enumerate}
    \item \textbf{Relatif supériorité :} \emph{Le plus compétent de tous}.
    \textbf{Accord :} \esp{candidato} (masc. sg.) $\rightarrow$ \esp{competente} (invariable genre, marque nombre). \textbf{Réponse :} \esp{Es el candidato más competente de todos.}
    \item \textbf{Relatif infériorité :} \emph{La moins agréable}. \esp{Tarea} (fém. sg.) $\rightarrow$ \esp{agradable}. \textbf{Réponse :} \esp{Es la tarea menos agradable del puesto.} (\esp{del} = \esp{de el}).
    \item \textbf{Absolu analytique (soutenu) :} \emph{Extrêmement occupé} $\rightarrow$ \esp{sumamente / extremadamente ocupado}. \textbf{Réponse :} \esp{El director está sumamente ocupado hoy.}
    \item \textbf{Absolu synthétique :} \emph{Enorme} $\rightarrow$ \esp{enorme} (termine par \emph{-e}) $\rightarrow$ \esp{enormísimo}. \textbf{Accord :} \esp{responsabilidad} (fém. sg.) $\rightarrow$ \esp{enormísima}. \textbf{Réponse :} \esp{Tenemos una responsabilidad enormísima.} (Accent sur \emph{í} obligatoire !).
    \item \textbf{Absolu analytique (standard) :} \emph{Très facile} $\rightarrow$ \esp{muy fácil}. \textbf{Réponse :} \esp{Es un problema muy fácil de resolver.}
\end{enumerate}

\textbf{2. Version}
\begin{enumerate}
    \item \esp{Es el empleado más puntual de la empresa.}
    $\rightarrow$ \emph{C'est l'employé le plus ponctuel de l'entreprise.} (Relatif \esp{más puntual de}).
    \item \esp{La entrevista fue dificilísima.}
    $\rightarrow$ \emph{L'entretien a été extrêmement difficile / d'une difficulté extrême.} (Absolu synthétique \esp{dificilísima} : \esp{difícil} $\rightarrow$ perte accent, +\esp{ísimo} $\rightarrow$ \esp{dificilísimo}). Accord fém. sg. avec \esp{entrevista}.
    \item \esp{Los requisitos son sumamente estrictos.}
    $\rightarrow$ \emph{Les exigences sont extrêmement strictes.} (Absolu analytique \esp{sumamente}). Accord masc. pl. \esp{estrictos}.
    \item \esp{Esta es la oferta menos interesante del mercado.}
    $\rightarrow$ \emph{C'est l'offre la moins intéressante du marché.} (Relatif infériorité \esp{menos interesante del}).
\end{enumerate}

\textbf{Conclusion :} Au Bac, la distinction \textbf{Relatif (article défini + más/menos + de/en)} vs \textbf{Absolu (-ísimo / muy / sumamente)} est systématiquement testée. Vérifiez toujours l'accord (genre/nombre) de l'adjectif et l'accent tonique sur le \emph{í} de \esp{-ísimo}.
\end{exoresolu}

\courssub{Méthode 4 : Analyser et traduire un texte mêlant comparatifs et superlatifs (Version)}

\begin{methode}[Version littéraire/journalistique : stratégie globale]
\textbf{Objectif :} Traduire un paragraphe cohérent (extrait de presse ou rapport) sans perdre les nuances de comparaison.
\begin{enumerate}
    \item \textbf{Lecture globale :} Identifier le thème (ici : marché de l'emploi jeunes Cameroun/Espagne), le ton (analytique, comparatif) et la structure (constats $\rightarrow$ comparaison $\rightarrow$ conclusion).
    \item \textbf{Surligner tous les comparatifs/superlatifs} (repérage visuel : \esp{más, menos, tan, tanto, mejor, peor, mayor, menor, -ísimo, muy, sumamente}).
    \item \textbf{Segmenter} : Découper en unités de sens (propositions principales / subordonnées).
    \item \textbf{Traduire phrase par phrase} en respectant :
    \begin{itemize}
        \item La tense (souvent présent/imparfait/passé composé en narration).
        \item Les connecteurs logiques (\esp{sin embargo, por tanto, en cambio}).
        \item Le vocabulaire spécifique (faux amis : \esp{carrera} $\neq$ carrière, \esp{éxito} $\neq$ exit).
    \end{itemize}
    \item \textbf{Relire la version française} : Doit être un français naturel, pas un "franglais" calqué sur la syntaxe espagnole.
\end{enumerate}
\end{methode}

\begin{exoresolu}[Version : Rapport sur l'employabilité des jeunes]
\textbf{Énoncé :} Traduisez le texte suivant en français soigné.

\begin{textebox}[Rapport fictif : Observatorio Laboral Hispano-Camerunés, 2024]
En los últimos años, la situación laboral de los jóvenes graduados en Camerún se ha vuelto \textbf{más compleja} que en la década anterior. Aunque hay \textbf{tantas titulaciones universitarias como antes}, la calidad del primer empleo es \textbf{menor}. Los salarios iniciales son \textbf{menos competitivos} que en el sector público y \textbf{mucho más bajos} que los de sus homólogos en España o Guinea Ecuatorial.

Sin embargo, los jóvenes cameruneses muestran una \textbf{adaptabilidad riquísima} y una resiliencia \textbf{mejor} que la media africana. El \textbf{mayor} obstáculo no es la formación, sino la falta de redes profesionales. \textbf{Sumamente} preparados, esperan oportunidades \textbf{tan dignas como} sus cualificaciones.
\end{textebox}

\tcblower

\textbf{Solution détaillée :}

\textbf{Repérage des structures (surlignées en gras dans le texte) :}
\begin{itemize}
    \item \esp{más compleja} (Comparatif supériorité régulier)
    \item \esp{tantas titulaciones... como} (Comparatif égalité quantité, accord fém. pl.)
    \item \esp{menor} (Comparatif irrégulier infériorité : \emph{plus petite / inférieure})
    \item \esp{menos competitivos} (Comparatif infériorité régulier, accord masc. pl.)
    \item \esp{mucho más bajos} (Comparatif supériorité intensifié par \esp{mucho}, accord masc. pl.)
    \item \esp{adaptabilidad riquísima} (Superlatif absolu synthétique \esp{rica} $\rightarrow$ \esp{riquísima}, accent + \emph{qu})
    \item \esp{mejor} (Comparatif irrégulier supériorité : \emph{meilleure})
    \item \esp{El mayor obstáculo} (Superlatif relatif / Adjectif irrégulier \emph{principal/majeur} \textbf{devant le nom})
    \item \esp{Sumamente preparados} (Superlatif absolu analytique, accord masc. pl. participe)
    \item \esp{tan dignas como} (Comparatif égalité qualité, accord fém. pl. avec \esp{oportunidades})
\end{itemize}

\textbf{Traduction proposée :}
\begin{quote}
{\em « Ces dernières années, la situation professionnelle des jeunes diplômés au Cameroun est devenue \textbf{plus complexe} qu'il y a dix ans. Bien qu'il y ait \textbf{autant de diplômes universitaires qu'auparavant}, la qualité du premier emploi est \textbf{inférieure}. Les salaires initiaux sont \textbf{moins compétitifs} que dans le secteur public et \textbf{beaucoup plus bas} que ceux de leurs homologues en Espagne ou en Guinée équatoriale.

Cependant, les jeunes Camerounais font preuve d'une \textbf{adaptabilité extraordinaire} et d'une résilience \textbf{meilleure} que la moyenne africaine. L'\textbf{obstacle majeur} n'est pas la formation, mais l'absence de réseaux professionnels. \textbf{Extrêmement bien} préparés, ils attendent des opportunités \textbf{aussi dignes que} leurs qualifications. »}
\end{quote}

\textbf{Justifications clés :}
\begin{itemize}
    \item \esp{Menor} $\rightarrow$ \emph{inférieure} (qualité abstraite), pas \emph{plus petite}.
    \item \esp{Riquísima} $\rightarrow$ \emph{extraordinaire / d'une grande richesse} (éviter \emph{très riche} qui sonne mal pour \emph{adaptabilité}).
    \item \esp{El mayor obstáculo} : \esp{Mayor} placé \textbf{avant} le nom = \emph{principal, majeur} (valeur superlative relative ou adjectif classificateur). Placé après (\esp{un obstáculo mayor}) = \emph{un obstacle plus grand}.
    \item \esp{Sumamente preparados} : Participe passé adjectivé, accord masc. pl. (\esp{jóvenes} sous-entendu).
    \item \esp{Tan dignas como} : Accord fém. pl. avec \esp{oportunidades}.
\end{itemize}
\end{exoresolu}

\courssub{Méthode 5 : Rédiger un Currículum Vítae en español}

\begin{methode}[Rédaction du CV : structure normalisée, vocabulaire précis, verbes d'action]
\textbf{Objectif :} Produire un CV standard hispanophone (format européen / \esp{Europass} adapté), clair et sans fautes.
\begin{enumerate}
    \item \textbf{En-tête (\esp{Datos Personales}) :}
    \esp{Nombre y apellidos} (Prénom NOM), \esp{Domicilio} (Adresse : Quartier, Ville, Pays), \esp{Teléfono} (avec indicatif \esp{+237}), \esp{Correo electrónico} (pro), \esp{LinkedIn} (optionnel), \esp{Fecha y lugar de nacimiento}, \esp{Nacionalidad}. \textbf{Pas de photo, pas d'âge, pas d'état civil} sur le modèle \esp{Europass} moderne (mais encore courant au Cameroun/Guinée : \esp{Foto}, \esp{Edad}, \esp{Estado civil}).
    \item \textbf{Profil (\esp{Perfil Profesional} / \esp{Objetivo}) :} 3-4 lignes. \emph{Qui suis-je ? Que cherche-je ? Quelle valeur ajoutée ?} Utiliser l'infinitif ou noms : \esp{Graduado en... con experiencia en... Busco...}
    \item \textbf{Formation (\esp{Formación Académica}) :} \textbf{Antéchronologique} (la plus récente en haut).
    \textbf{Modèle :} \esp{2021--2024 \textbf{Grado en Administración de Empresas} (Licence), Universidad de Yaoundé II, Soa. \textit{Mención: Bien. TFG: "Estrategias de exportación de cacao".}}
    \item \textbf{Expérience (\esp{Experiencia Profesional / Prácticas}) :} Antéchronologique.
    \textbf{Modèle :} \esp{06/2023--08/2023 \textbf{Becario en Marketing Digital}, Empresa XYZ, Douala. \textit{- Gestión de redes sociales. - Análisis de métricas (Google Analytics). - Redacción de contenido SEO.}} \textbf{Verbes d'action (infinitif ou nominalisés) :} \esp{Gestionar, Coordinar, Desarrollar, Implementar, Resolver, Supervisar, Redactar, Analizar, Negociar, Atender.}
    \item \textbf{Langues (\esp{Idiomas}) :} Niveau \esp{MCER} (A1--C2). \esp{Francés: Lengua materna / C2. Español: C1 (DELE). Inglés: B2.}
    \item \textbf{Informatique (\esp{Competencias Digitales / Informática}) :} \esp{Paquete Office (Avanzado), Python (Básico), Canva, ERP SAP.}
    \item \textbf{Autres (\esp{Otros Datos / Información Adicional}) :} Permis (\esp{Carnet de conducir B}), Bénévolat (\esp{Voluntariado}), Associations, Centres d'intérêt (\esp{Intereses}).
\end{enumerate}
\end{methode}

\begin{exoresolu}[Rédaction guidée : CV d'un étudiant Terminale A]
\textbf{Énoncé :} Vous êtes \textbf{Marie Ngo Bibaa}, élève de Terminale A au Lycée Leclerc (Yaoundé). Vous postulez pour un \emph{stage d'été (prácticas de verano)} dans une agence de tourisme réceptive à \textbf{Malabo (Guinée équatoriale)}. Rédigez votre CV en espagnol (format simplifié, une page). Inventez des détails cohérents (notes, langues, bénévolat).

\tcblower

\textbf{Solution détaillée (Modèle de CV rédigé) :}

\begin{textebox}[Currículum Vítae -- Marie Ngo Bibaa]
\textbf{DATOS PERSONALES} \\
\textbf{Nombre:} Marie Ngo Bibaa \\
\textbf{Domicilio:} Quartier Mvog-Ada, Yaoundé, Camerún (+237 6XX XX XX XX) \\
\textbf{Email:} marie.ngo.bibaa@email.com \quad \textbf{LinkedIn:} linkedin.com/in/mariego \\
\textbf{Fecha de nacimiento:} 12/03/2006 \quad \textbf{Nacionalidad:} Camerunesa \\
\textbf{Permiso de conducir:} Tipo B (en trámite) \\

\textbf{PERFIL PROFESIONAL} \\
Estudiante de Terminale A (Bachillerato Literario) con nivel C1 en español (DELE B2 obtenido) y certificado en inglés (B1). Apasionada por el turismo cultural y la hospitalidad. Busco prácticas de verano (julio--agosto) para aplicar mis competencias lingüísticas y organizativas en una agencia receptiva en Malabo. Disponibilidad inmediata, gran adaptabilidad y sentido del servicio. \\

\textbf{FORMACIÓN ACADÉMICA} \\
\textbf{2021--2024} \quad \textbf{Bachillerato General (Serie A -- Literatura/Lenguas)}, Lycée Leclerc, Yaoundé. \\
\quad \textit{Promedio general: 14,5/20. Mención: Bien. Especialidad: Español (18/20), Historia (16/20).} \\
\textbf{2022} \quad \textbf{DELE Nivel B2 (Diploma de Español como Lengua Extranjera)}, Instituto Cervantes, Yaoundé. \\

\textbf{EXPERIENCIA PROFESIONAL Y PRÁCTICAS} \\
\textbf{07/2023--08/2023} \quad \textbf{Voluntaria / Animadora sociocultural}, Asociación "Jeunesse Active", Yaoundé. \\
\quad \textit{- Organización de talleres de lectura en español para niños (10--14 años).} \\
\quad \textit{- Coordinación de un equipo de 5 voluntarios.} \\
\quad \textit{- Gestión de agenda y logística de salidas culturales.} \\
\textbf{01/2023--06/2023} \quad \textbf{Delegada de curso}, Lycée Leclerc. \\
\quad \textit{- Representación de los alumnos ante la administración.} \\
\quad \textit{- Mediación de conflictos y organización de la fiesta de fin de curso.} \\

\textbf{IDIOMAS (MCER)} \\
\textbf{Francés:} Lengua materna (C2). \quad \textbf{Español:} C1 (DELE B2 certificado, práctica oral fluida). \\
\textbf{Inglés:} B1 (PET Cambridge). \quad \textbf{Bassa:} Nativo (oral). \\

\textbf{COMPETENCIAS DIGITALES} \\
Paquete Office (Word, Excel, PowerPoint: Nivel Avanzado). Redes Sociales (Instagram, TikTok: Gestión de contenido). Canva (Diseño gráfico básico). Navegación web y correo electrónico profesional. \\

\textbf{INFORMACIÓN ADICIONAL} \\
\textbf{Voluntariado:} Cruz Roja Juvenil Camerún (2022--presente) -- Campañas de donación de sangre. \\
\textbf{Intereses:} Literatura africana e hispanoamericana, Danza tradicional (Bikutsi), Fútbol, Cocina camerunesa. \\
\textbf{Referencias:} Disponibles a petición (Prof. Principal, Lycée Leclerc).
\end{textebox}

\textbf{Commentaires pédagogiques :}
\begin{itemize}
    \item \textbf{Vocabulaire clé :} \esp{Bachillerato} (Lycée), \esp{Promedio} (Moyenne), \esp{Mención} (Mention), \esp{Delegada de curso} (Déléguée de classe), \esp{Prácticas} (Stage), \esp{Voluntariado} (Bénévolat), \esp{Carnet de conducir} (Permis).
    \item \textbf{Verbes d'action (noms en -\emph{ción}/\emph{-sión} comme titres de puces) :} \esp{Organización, Coordinación, Gestión, Representación, Mediación}. C'est la norme espagnole (nominalisation). Éviter la 1ère personne (\esp{Organicé...}) dans un CV moderne.
    \item \textbf{Mise en page :} Alignement à gauche, police sobre (Arial, Calibri, Times), gras pour les titres de poste/diplôme, italique pour lieu/détails.
    \item \textbf{Adaptation contexte :} Mentionner \esp{Guinea Ecuatorial} comme cible montre la motivation. Niveau \esp{C1} justifié par \esp{DELE B2} + pratique.
\end{itemize}
\end{exoresolu}

\courssub{Méthode 6 : Se présenter à l'oral à partir de son CV (Expression orale guidée)}

\begin{methode}[Présentation orale type "Entretien / Elevator Pitch" (2-3 min)]
\textbf{Objectif :} Structurer une prise de parole fluide, corrigeant les hésitations, pour un entretien de stage ou d'embauche.
\begin{enumerate}
    \item \textbf{Accroche (30 s) :} Salutation formelle (\esp{Buenos días, estimados señores / Buenas tardes}). Prénom, Nom, Âge, Ville d'origine, Situation actuelle. \esp{"Me llamo Marie Ngo Bibaa, tengo 18 años, vengo de Yaundé y soy estudiante de Terminale A..."}
    \item \textbf{Parcours \& Motivation (1 min) :} Diplôme en cours + Spécialité + Niveau langues (certificats). \textbf{Pourquoi ce poste ?} Lien études/poste. \esp{"Mi formación en serie A me ha dado una base cultural sólida... Mi nivel C1 en español me permite atender a clientes hispanohablantes..."}
    \item \textbf{Expériences \& Compétences transverses (1 min) :} 1 expérience significative (bénévolat, job d'été, délégué). Compétences \emph{soft skills} prouvées. \esp{"Como delegada de curso, desarrollé liderazgo y resolución de conflictos... En la asociación, organicé talleres..."}
    \item \textbf{Atouts techniques / Savoir-faire (30 s) :} Informatique, Permis, Connaissances spécifiques (Tourisme, Cacao, Culture). \esp{"Domino Office y Canva... Conozco la realidad turística de la zona CEMAC..."}
    \item \textbf{Conclusion \& Disponibilité (15 s) :} Réitérer intérêt, disponibilité, question ouverte. \esp{"Estoy muy motivada para integrar su equipo. Tengo disponibilidad total en julio y agosto. Quedo a su disposición para cualquier pregunta. Muchas gracias."}
\end{enumerate}
\textbf{Astuce :} Préparer des "mots-clés" sur une fiche (pas des phrases apprises par cœur). S'entraîner à chronométrer.
\end{methode}

\begin{exoresolu}[Simulation orale : Entretien pour le stage à Malabo]
\textbf{Énoncé :} À partir du CV de Marie (Exercice résolu 5), rédigez le \textbf{script complet} de sa présentation orale de 2 minutes (environ 250 mots), en espagnol naturel, prêt à être appris/révisé. Incluez les connecteurs logiques et les formules de politesse.

\tcblower

\textbf{Script de présentation (Modèle corrigé) :}

\begin{textebox}[Simulación de Entrevista -- Marie Ngo Bibaa]
\textbf{Marie :} \esp{Buenos días, señor Director, señora, señores. Me llamo Marie Ngo Bibaa, tengo dieciocho años y vengo de Yaundé, la capital de Camerún.}

\esp{Actualmente, curso el Bachillerato General, Serie A (Literatura y Lenguas), en el Lycée Leclerc. Mi expediente académico es sólido: tengo un promedio de 14,5 sobre 20, con mención "Bien", y destaco especialmente en Español --saqué 18 sobre 20-- e Historia. Además, obtuve el diploma DELE B2 del Instituto Cervantes el año pasado, y mi nivel oral roza ya el C1.}

\esp{Mi motivación para solicitar estas prácticas en su agencia de Malabo es doble. Primero, mi formación literaria me ha dado una gran sensibilidad cultural y una capacidad de redacción impecable, esenciales para la comunicación con el cliente. Segundo, domino el español a un nivel profesional y conozco bien el contexto de la zona CEMAC, lo que facilita la adaptación.}

\esp{En cuanto a mi experiencia, aunque soy joven, he asumido responsabilidades reales. Fui delegada de mi curso, lo que me enseñó a mediar conflictos y a organizar eventos para más de 50 personas. El verano pasado, fui voluntaria en la asociación "Jeunesse Active": organicé talleres de lectura en español para niños y coordiné un equipo de cinco personas. Estas experiencias me han dado \textbf{liderazgo, paciencia y sentido de la organización}.}

\esp{A nivel técnico, manejo perfectamente el paquete Office, diseño en Canva y gestiono redes sociales. Tengo el carnet de conducir tipo B en trámite y total disponibilidad para julio y agosto.}

\esp{Estoy convencida de que mi adaptabilidad --una \textbf{riqueza} que cultivo cada día-- y mis ganas de aprender serán un activo para su equipo. Quedo a su entera disposición para ampliar cualquier aspecto. Muchas gracias por su tiempo.}
\end{textebox}

\textbf{Analyse linguistique (points de vigilance pour l'oral) :}
\begin{itemize}
    \item \textbf{Connecteurs :} \esp{Actualmente, Además, En cuanto a, Por tanto / Por eso (implícito), Muchas gracias.}
    \item \textbf{Temps :} Présent (\esp{curso, tengo, domino}), Passé composé (\esp{he asumido, fui, obtuve, organicé}) pour expériences terminées.
    \item \textbf{Subjonctif (naturel) :} \esp{Quedo a su disposición} (présent indicatif mais valeur de politesse), \esp{Es esencial que...} (si développé).
    \item \textbf{Vocabulaire "pro" :} \esp{Expediente académico, Mención, Sensibilidad cultural, Mediar conflictos, Activo, Disponibilidad.}
    \item \textbf{Prononciation :} Travailler \esp{ll} (yeísmo standard), \esp{j/g} (guttural), accents toniques (\esp{expediénte, académico, motivación, organización, disponibilidad}).
\end{itemize}
\end{exoresolu}

\courssub{Attention : Les 5 erreurs qui coûtent des points au Bac}

\begin{attention}[Les 5 erreurs fatales sur les comparatifs, superlatifs et le CV]
\begin{enumerate}
    \item \textbf{L'accent sur \esp{más} (comparatif) vs \esp{mas} (conjonction = \emph{mais}).}

    \cle{Règle :} \esp{Más} (monosyllabe tonique) \textbf{porte toujours l'accent} quand c'est un quantificateur (\emph{plus}). \esp{Mas} (conjonction adversative) \textbf{jamais d'accent}.
    \item \textbf{L'accord de \esp{tanto/tanta/tantos/tantas} (comparatif d'égalité quantité).}
    \cle{Piège :} Oublier l'accord avec le nom. \esp{Tanto trabajo} (masc. sg.), \esp{Tanta gente} (fém. sg.), \esp{Tantos libros} (masc. pl.), \esp{Tantas oportunidades} (fém. pl.). \textbf{Jamais} \esp{*tanto oportunidades}.
    \item \textbf{Confusion \esp{mayor/menor} (âge, rang, importance) vs \esp{más grande/más pequeño} (taille physique).}
    \cle{Erreur type :} \esp{*Mi hermano es más grande que yo} (pour l'âge). \textbf{Correct :} \esp{Mi hermano es mayor que yo.} \esp{Mayor/Menor} s'emploient aussi \textbf{devant le nom} pour signifier \emph{principal/majeur} (\esp{El mayor problema}).
    \item \textbf{L'accent sur le suffixe \esp{-ísimo} (superlatif absolu synthétique).}
    \cle{Règle d'or :} L'accent tombe \textbf{toujours} sur le \emph{í} : \esp{-ísim-o, -ísim-a, -ísim-os, -ísim-as}. \textbf{Jamais} \esp{*facilisimo, *rapidisima}. Attention aux changements d'orthographe : \esp{rico $\rightarrow$ riquísimo, largo $\rightarrow$ larguísimo, amable $\rightarrow$ amabilísimo}.
    \item \textbf{Faux amis et calques du français dans le CV (\emph{Rédac}).}
    \begin{center}
    \begin{tabularx}{\linewidth}{|X|X|X|}\hline
    \rowcolor{poporangeL} \textbf{Français (à éviter)} & \textbf{Espagnol correct} & \textbf{Note} \\ \hline
    \emph{Carrière} / \emph{Parcours} & \esp{Trayectoria profesional} / \esp{Experiencia} & \esp{Carrera} = Carrière universitaire (course) ou filière. \\ \hline
    \emph{Licence} & \esp{Grado} (Bologne) / \esp{Licenciatura} (ancien) & \esp{Licenciatura} = Master (4-5 ans) avant Bologne. \\ \hline
    \emph{Stage} & \esp{Prácticas} / \esp{Prácticas profesionales} & \esp{Stage} n'existe pas. \esp{Estancia} = séjour. \\ \hline
    \emph{Responsable de...} & \esp{Encargado de...} / \esp{Responsable de...} & \esp{Responsable} existe (adj/nom) mais \esp{Encargado} plus précis pour "chargé de". \\ \hline
    \emph{Compétences} & \esp{Competencias} / \esp{Habilidades} & \esp{Competencias} (savoir-faire), \esp{Habilidades} (soft skills). \\ \hline
    \end{tabularx}
    \end{center}
\end{enumerate}
\textbf{Conseil final :} Relisez votre copie \textbf{à l'envers} (dernière phrase $\rightarrow$ première) pour chasser les fautes d'accords et d'accents (aveuglement de la relecture linéaire).

\end{attention}


\newpage

\coursec{Exercices}

\courssub{Je vérifie mes connaissances}

\begin{exercice}[QCM : comparatifs et superlatifs]
Pour chaque phrase, choisis la bonne réponse (a, b ou c).
\begin{enumerate}
\item Mi hermano es \ldots{} alto que yo. \quad a) tan \quad b) más \quad c) muy
\item Este hotel es \ldots{} caro \ldots{} el otro. \quad a) tan / que \quad b) más / como \quad c) tan / como
\item Es la \ldots{} empleada de la empresa. \quad a) más buena \quad b) mejor \quad c) buenísima
\item Tengo \ldots{} experiencia que tú. \quad a) tanta \quad b) tan \quad c) tanto
\item El examen fue \ldots{} difícil. (superlatif absolu) \quad a) dificilísimo \quad b) más difícil \quad c) muy difícil que
\item Hablo \ldots{} de tres idiomas. (plus de) \quad a) más que \quad b) más de \quad c) tan
\item Mi padre es \ldots{} que mi madre. (plus âgé) \quad a) más viejo \quad b) mayor \quad c) más mayor
\item Este mercado es \ldots{} de Douala. (le moins cher) \quad a) el menos barato \quad b) el más barato \quad c) el menos caro
\end{enumerate}
\end{exercice}

\begin{exercice}[Vrai ou faux ? Justifie ta réponse]
Dis si les affirmations suivantes sont vraies (V) ou fausses (F), puis justifie en une phrase en français.
\begin{enumerate}
\item On peut dire \esp{más bueno} pour « meilleur ».
\item \esp{Mayor} peut signifier « plus âgé » ou « plus grand » selon le contexte.
\item Devant un chiffre, « plus de » se traduit par \esp{más de} et non \esp{más que}.
\item Le superlatif absolu en \esp{-ísimo} s'accorde en genre et en nombre avec le nom.
\item \esp{Tan\ldots{} como} s'emploie devant un nom, \esp{tanto\ldots{} como} devant un adjectif.
\item Dans un CV espagnol, la rubrique \esp{Formación académica} correspond à « Expérience professionnelle ».
\end{enumerate}
\end{exercice}

\begin{exercice}[Vocabulaire du monde du travail]
Complète chaque phrase avec le mot espagnol qui convient, choisi dans la liste : \esp{empresa}, \esp{puesto}, \esp{salario}, \esp{entrevista}, \esp{candidato}, \esp{experiencia laboral}, \esp{idiomas}, \esp{referencias}.
\begin{enumerate}
\item Quiero solicitar el \ldots{} de secretario en vuestra \ldots{}.
\item El \ldots{} mensual es de ciento cincuenta mil francos.
\item Tengo cinco años de \ldots{} en el sector del cacao.
\item Mañana tengo una \ldots{} de trabajo en Yaoundé.
\item Cada \ldots{} debe enviar su currículum vítae antes del lunes.
\item En la rubrica \esp{Idiomas} del CV, indico que hablo español, francés e inglés : son tres \ldots{}.
\item Puedo dar dos \ldots{} : mis antiguos jefes.
\end{enumerate}
\end{exercice}

\begin{exercice}[Formation du superlatif absolu]
Forme le superlatif absolu en \esp{-ísimo} des adjectifs suivants, en respectant les changements orthographiques, puis traduis le résultat en français.
\begin{enumerate}
\item caro (una casa \ldots)
\item fácil (un examen \ldots)
\item feliz (unos niños \ldots)
\item grande (una empresa \ldots)
\item bueno (unas noticias \ldots)
\item rápido (un tren \ldots)
\item pobre (un barrio \ldots)
\item interesante (unas ofertas \ldots)
\end{enumerate}
\end{exercice}

\courssub{Je m'entraîne}

\begin{exercice}[Texte à trous : deux entreprises camerounaises]
Complète le texte suivant avec : \esp{más}, \esp{menos}, \esp{tan}, \esp{tanta}, \esp{que}, \esp{como}, \esp{de}, \esp{mejor}, \esp{peor}, \esp{mayor}, \esp{menor}. Chaque mot ne peut servir qu'une seule fois.
\begin{textebox}[Dos empresas de Douala]
La empresa Njoka es \ldots{} grande \ldots{} la empresa Wouri : tiene \ldots{} de doscientos empleados, mientras que Wouri tiene \ldots{} de cien. El director de Njoka es \ldots{} que el de Wouri : tiene sesenta años. Los empleados de Wouri ganan \ldots{} dinero \ldots{} los de Njoka, pero trabajan \ldots{} horas. Según los clientes, el servicio de Njoka es el \ldots{} de la ciudad, y su precio es el \ldots{} : ¡es carísimo! Sin embargo, Wouri no tiene \ldots{} clientes \ldots{} Njoka, porque sus productos son de \ldots{} calidad. La competencia entre las dos empresas es cada vez \ldots{} fuerte.
\end{textebox}
\end{exercice}

\begin{exercice}[Reformulation : muy et -ísimo]
\begin{enumerate}
\item Transforme les phrases suivantes en remplaçant \esp{muy} + adjectif par le superlatif absolu en \esp{-ísimo} :

a) El salario es muy alto. \quad b) La entrevista fue muy larga. \quad c) Este candidato es muy joven. \quad d) La oficina está muy limpia.
\item Transforme les phrases suivantes en remplaçant le superlatif absolu par \esp{muy} + adjectif :

a) El jefe es riquísimo. \quad b) Las preguntas eran dificilísimas. \quad c) Mi colega es simpatiquísimo. \quad d) El viaje fue larguísimo.
\end{enumerate}
\end{exercice}

\begin{exercice}[Appariements : les rubriques du CV]
Associe chaque rubrique espagnole (1 à 6) à son contenu (a à f).
\begin{enumerate}
\item \esp{Datos personales}
\item \esp{Formación académica}
\item \esp{Experiencia laboral}
\item \esp{Idiomas}
\item \esp{Informática}
\item \esp{Otros datos de interés}
\end{enumerate}
\begin{itemize}
\item[a)] Word, Excel, Internet
\item[b)] español : nivel intermedio ; francés : lengua materna
\item[c)] nombre, dirección, teléfono, fecha de nacimiento
\item[d)] permiso de conducir, aficiones, disponibilidad
\item[e)] bachillerato, serie A4, 2023, Lycée de Yaoundé
\item[f)] vendedor en un supermercado, 2022-2023
\end{itemize}
\end{exercice}

\begin{exercice}[Correction d'erreurs typiques de francophones]
Chaque phrase contient une faute. Souligne-la et corrige-la.
\begin{enumerate}
\item Mi hermana es más buena que yo en matemáticas.
\item Tengo tanto experiencia como mi colega.
\item Es el más alto de la clase que todos.
\item Este puesto es más peor que el anterior.
\item Gano más que cien mil francos al mes.
\item Ella es tan inteligente como su hermano es.
\item La reunión fue muy larguísima.
\item Mi padre es más viejo que mi tío : tiene setenta años. (pour « plus âgé »)
\end{enumerate}
\end{exercice}

\courssub{Je me prépare au Bac}

\begin{exercice}[Texte, compréhension et langue (type Bac)]
Lis le texte suivant, puis réponds aux questions.
\begin{textebox}[El mercado de trabajo en España]
En España, encontrar un primer empleo es cada vez más difícil para los jóvenes. Los licenciados son hoy más numerosos que hace veinte años, pero las empresas ofrecen menos puestos estables que antes. Muchos jóvenes aceptan trabajos peor pagados que los de sus padres, y el salario de un principiante es a menudo tan bajo como el de un obrero sin estudios. Sin embargo, los candidatos que hablan tantos idiomas como piden las empresas internacionales tienen mejores oportunidades. Según los expertos, la situación es dificilísima en las regiones del sur, donde el paro juvenil es altísimo. Preparar un buen currículum vítae es hoy más importante que nunca : es la primera imagen que el candidato da de sí mismo.
\end{textebox}
\textbf{Questions de compréhension} (réponds en espagnol, avec des phrases complètes) :
\begin{enumerate}
\item ¿Por qué es difícil encontrar un primer empleo en España, según el texto ?
\item ¿Qué diferencia hay entre los salarios de los jóvenes y los de sus padres ?
\item ¿Qué candidatos tienen mejores oportunidades ? ¿Por qué ?
\item ¿Dónde es más grave el paro juvenil ?
\end{enumerate}
\textbf{Questions de langue} :
\begin{enumerate}
\item Relève dans le texte trois comparatifs et deux superlatifs absolus, puis traduis-les en français.
\item Traduis en français la dernière phrase du texte.
\end{enumerate}
\end{exercice}

\begin{exercice}[Thème et version (type Bac)]
\textbf{Thème.} Traduis en espagnol les cinq phrases suivantes :
\begin{enumerate}
\item Mon frère est plus âgé que moi.
\item Il gagne moins d'argent que sa femme.
\item C'est le meilleur employé de l'entreprise.
\item Elle parle autant de langues que toi.
\item Plus de cinquante candidats ont postulé.
\end{enumerate}
\textbf{Version.} Traduis en français le texte suivant :
\begin{textebox}[Trabajar en Yaoundé]
En Yaoundé, la vida es más cara que en las ciudades pequeñas del país. Los alquileres son altísimos en el centro, y los transportes son menos cómodos que en Douala. Sin embargo, hay más ofertas de empleo que en otras regiones, y los jóvenes llegan cada día más numerosos. Un empleado de banca gana tanto como dos profesores juntos. Para los empresarios, el candidato ideal es puntualísimo, serio y habla por lo menos dos idiomas. Buscar trabajo es hoy más difícil que antes, pero no imposible : con un buen currículum y mucha paciencia, todo es posible.
\end{textebox}
\end{exercice}

\begin{exercice}[Rédaction guidée : mon CV et ma lettre de motivation (type Bac)]
Lis l'annonce d'emploi suivante, puis rédige ta production.
\begin{textebox}[Se busca empleado]
Empresa camerunesa de importación-exportación, Douala. Se busca empleado/a de oficina. Requisitos : bachillerato, buen nivel de español y de francés, conocimientos de informática, experiencia de un año mínimo. Se ofrece : contrato de un año, salario interesante. Enviar currículum vítae y carta de motivación antes del 30 de junio.
\end{textebox}
\textbf{Consigne.} Tu es candidat à ce poste. Rédige en espagnol :
\begin{enumerate}
\item ton \textbf{currículum vítae} complet, avec les 6 rubriques suivantes : \esp{Datos personales}, \esp{Formación académica}, \esp{Experiencia laboral}, \esp{Idiomas}, \esp{Informática}, \esp{Otros datos de interés} ;
\item le début de ta \textbf{carta de motivación} : 5 phrases au minimum, dans lesquelles tu utiliseras au moins deux comparatifs et un superlatif.
\end{enumerate}
Nombre de mots attendu : 120 à 150 mots au total. Tu soigneras la présentation, l'orthographe et les accents.
\end{exercice}

\newpage

\coursec{Corrigés des exercices}

\coursec{Corrigés des exercices — Leçon E17}

\begin{corrige}[Exercice 1 — QCM : comparatifs et superlatifs]
\begin{enumerate}
\item \textbf{b) más} : \esp{Mi hermano es más alto que yo.} — Comparatif de supériorité : \esp{más} + adjectif + \esp{que}.
\item \textbf{c) tan / como} : \esp{Este hotel es tan caro como el otro.} — Comparatif d'égalité devant un adjectif.
\item \textbf{b) mejor} : \esp{Es la mejor empleada de la empresa.} — « Meilleur » se dit \esp{mejor} ; \esp{más buena} est fautif ici (comparatif irrégulier, \esp{más buena} ne s'emploie qu'au sens moral « plus vertueuse »).
\item \textbf{a) tanta} : \esp{Tengo tanta experiencia como tú.} — Devant un nom féminin, \esp{tanto} s'accorde : \esp{tanta experiencia} ; le second terme de l'égalité est introduit par \esp{como}.
\item \textbf{a) dificilísimo} : \esp{El examen fue dificilísimo.} — Superlatif absolu en \esp{-ísimo} (changement d'accent : \esp{difícil} $\rightarrow$ \esp{dificilísimo}) ; \esp{más difícil} est un comparatif, \esp{muy difícil que} est incorrect.
\item \textbf{b) más de} : \esp{Hablo más de tres idiomas.} — Devant un chiffre, « plus de » = \esp{más de}, jamais \esp{más que}.
\item \textbf{b) mayor} : \esp{Mi padre es mayor que mi madre.} — Pour l'âge, on utilise le comparatif irrégulier \esp{mayor}.
\item \textbf{a) el menos barato} : \esp{Este mercado es el menos barato de Douala.} — « Le moins cher » = superlatif relatif d'infériorité. (La réponse c, \esp{el menos caro}, signifie « le moins coûteux » et serait acceptable en espagnol courant, mais la formulation attendue avec \esp{barato} est la a.)
\end{enumerate}
\end{corrige}

\begin{corrige}[Exercice 2 — Vrai ou faux ?]
\begin{enumerate}
\item \textbf{F.} « Meilleur » se traduit par le comparatif irrégulier \esp{mejor} ; \esp{más bueno} n'est admis qu'au sens moral (« plus bon, plus charitable »).
\item \textbf{V.} \esp{Mayor} signifie « plus âgé » pour les personnes (\esp{mi hermano mayor}) et « plus grand / plus important » dans d'autres contextes (\esp{un problema mayor}).
\item \textbf{V.} Devant un chiffre ou une quantité déterminée, « plus de » = \esp{más de} : \esp{más de cien candidatos}.
\item \textbf{V.} L'adjectif en \esp{-ísimo} s'accorde en genre et en nombre : \esp{carísimo, carísima, carísimos, carísimas}.
\item \textbf{F.} C'est l'inverse : \esp{tan\ldots{} como} s'emploie devant un adjectif ou un adverbe (\esp{tan alto como}), \esp{tanto\ldots{} como} devant un nom (\esp{tanto dinero como}).
\item \textbf{F.} \esp{Formación académica} correspond à « Formation / Études » ; « Expérience professionnelle » se dit \esp{Experiencia laboral}.
\end{enumerate}
\end{corrige}

\begin{corrige}[Exercice 3 — Vocabulaire du monde du travail]
\begin{enumerate}
\item \esp{Quiero solicitar el \textbf{puesto} de secretario en vuestra \textbf{empresa}.} — « Je voudrais solliciter le poste de secrétaire dans votre entreprise. »
\item \esp{El \textbf{salario} mensual es de ciento cincuenta mil francos.} — « Le salaire mensuel est de cent cinquante mille francs. »
\item \esp{Tengo cinco años de \textbf{experiencia laboral} en el sector del cacao.}
\item \esp{Mañana tengo una \textbf{entrevista} de trabajo en Yaoundé.} — Attention au faux ami : \esp{entrevista} = « entretien d'embauche ».
\item \esp{Cada \textbf{candidato} debe enviar su \textbf{currículum vítae} antes del lunes.}
\item \esp{En la \textbf{rúbrica} \esp{Idiomas} del CV, indico que hablo español, francés e inglés : son tres \textbf{idiomas}.} — \cle{rúbrica} (accent sur le \emph{u}) = rubrique.
\item \esp{Puedo dar dos \textbf{referencias} : mis antiguos jefes.}
\end{enumerate}
\end{corrige}

\begin{corrige}[Exercice 4 — Formation du superlatif absolu en \esp{-ísimo}]
\begin{enumerate}
\item \esp{una casa \textbf{carísima}} — « une maison très chère » (voyelle finale \emph{-o} qui tombe).
\item \esp{un examen \textbf{facilísimo}} — « un examen très facile » (changement orthographique \esp{c} $\rightarrow$ \esp{qu} devant \emph{i} pour garder le son [k]).
\item \esp{unos niños \textbf{felicísimos}} — « des enfants très heureux » (\esp{z} $\rightarrow$ \esp{c} devant \emph{i}).
\item \esp{una empresa \textbf{grandísima}} — « une très grande entreprise » (forme pleine \esp{grande} + \esp{ísimo}).
\item \esp{unas noticias \textbf{buenísimas}} — « de très bonnes nouvelles » (accent sur la diphtongue rompue : \esp{buenísimo}, \esp{buenísima}).
\item \esp{un tren \textbf{rapidísimo}} — « un train très rapide » (l'accent de \esp{rápido} se déplace sur le \emph{í} du suffixe : \esp{rapidísimo}).
\item \esp{un barrio \textbf{pobrísimo}} — « un quartier très pauvre » (forme irrégulière, à côté de \esp{paupérrimo}, plus soutenu).
\item \esp{unas ofertas \textbf{interesantísimas}} — « des offres très intéressantes ».
\end{enumerate}
\textbf{Rappel des règles :} la voyelle finale tombe (\esp{caro} $\rightarrow$ \esp{carísimo}) ; \esp{c} $\rightarrow$ \esp{qu}, \esp{z} $\rightarrow$ \esp{c}, \esp{g} $\rightarrow$ \esp{gu} devant \emph{i} ; l'accent tonique porte toujours sur le premier \emph{í} du suffixe.
\end{corrige}

\begin{corrige}[Exercice 5 — Texte à trous : deux entreprises de Douala]
\textbf{Grille de correction (chaque mot de la liste utilisé une seule fois) :}

\begin{center}
\begin{tabularx}{\linewidth}{|X|X|}\hline
\rowcolor{popblueL} Trou & Réponse \\ \hline
es \ldots{} grande \ldots{} & tan / como \\ \hline
\ldots{} de doscientos & más \\ \hline
\ldots{} de cien & menos \\ \hline
el director es \ldots{} que & mayor \\ \hline
ganan \ldots{} dinero \ldots{} los de Njoka & menos / que \\ \hline
trabajan \ldots{} horas & más \\ \hline
el servicio es el \ldots{} & mejor \\ \hline
su precio es el \ldots{} & peor \\ \hline
no tiene \ldots{} clientela \ldots{} Njoka & tanta / como \\ \hline
de \ldots{} calidad & menor \\ \hline
cada vez \ldots{} fuerte & más \\ \hline
\end{tabularx}
\end{center}

\textbf{Texte corrigé (version harmonisée) :}

\begin{textebox}[Comparatif d'entreprises à Douala]
\esp{La empresa Njoka es \textbf{tan} grande \textbf{como} la empresa Wouri : tiene \textbf{más} de doscientos empleados, mientras que Wouri tiene \textbf{menos} de cien. El director de Njoka es \textbf{mayor} que el de Wouri : tiene sesenta años. Los empleados de Wouri ganan \textbf{menos} dinero \textbf{que} los de Njoka, pero trabajan \textbf{más} horas. Según los clientes, el servicio de Njoka es el \textbf{mejor} de la ciudad, y su precio es el \textbf{peor} : \textbf{¡es carísimo!} Sin embargo, Wouri no tiene \textbf{tanta} clientela \textbf{como} Njoka, porque sus productos son de \textbf{menor} calidad. La competencia entre las dos empresas es cada vez \textbf{más} fuerte.}
\end{textebox}

\textbf{Traduction :} « L'entreprise Njoka est aussi grande que l'entreprise Wouri : elle compte plus de deux cents employés, tandis que Wouri en compte moins de cent. Le directeur de Njoka est plus âgé que celui de Wouri : il a soixante ans. Les employés de Wouri gagnent moins d'argent que ceux de Njoka, mais travaillent plus d'heures. Selon les clients, le service de Njoka est le meilleur de la ville, et son prix le pire : il est très cher ! Cependant, Wouri n'a pas autant de clientèle que Njoka, car ses produits sont de moindre qualité. La concurrence entre les deux entreprises est de plus en plus forte. »
\end{corrige}

\begin{corrige}[Exercice 6 — Reformulation : \esp{muy} et \esp{-ísimo}]
\textbf{1. De \esp{muy} + adjectif au superlatif absolu (exprimé) :}
\begin{enumerate}
\item[a)] \esp{El salario es \textbf{altísimo}.} (\esp{alto} $\rightarrow$ \esp{altísimo})
\item[b)] \esp{La entrevista fue \textbf{larguísima}.} (\esp{larga}, \esp{g} $\rightarrow$ \esp{gu} devant \emph{i})
\item[c)] \esp{Este candidato es \textbf{jovencísimo}.} (forme pleine : \esp{joven} $\rightarrow$ \esp{jovencísimo})
\item[d)] \esp{La oficina está \textbf{limpísima}.} (\esp{limpia} $\rightarrow$ \esp{limpísima})
\end{enumerate}

\textbf{2. Du superlatif absolu à \esp{muy} + adjectif (neutre) :}
\begin{enumerate}
\item[a)] \esp{El jefe es \textbf{muy rico}.} (\esp{riquísimo})
\item[b)] \esp{Las preguntas eran \textbf{muy difíciles}.} (\esp{dificilísimas})
\item[c)] \esp{Mi colega es \textbf{muy simpático}.} (\esp{simpatiquísimo})
\item[d)] \esp{El viaje fue \textbf{muy largo}.} (\esp{larguísimo})
\end{enumerate}

\textbf{Point de vigilance :} le sens est équivalent, mais le registre diffère : \esp{-ísimo} est plus expressif, plus courant à l'oral en Espagne ; \esp{muy} + adjectif est neutre. On ne cumule jamais les deux (\esp{*muy larguísima}).
\end{corrige}

\begin{corrige}[Exercice 7 — Appariements : les rubriques du CV]
\begin{center}
\begin{tabularx}{\linewidth}{|X|X|}\hline
\rowcolor{popblueL} Rubrique & Contenu \\ \hline
1. \esp{Datos personales} & c) nombre, dirección, teléfono, fecha de nacimiento \\ \hline
2. \esp{Formación académica} & e) bachillerato, serie A4, 2023, Lycée de Yaoundé \\ \hline
3. \esp{Experiencia laboral} & f) vendedor en un supermercado, 2022-2023 \\ \hline
4. \esp{Idiomas} & b) español : nivel intermedio ; francés : lengua materna \\ \hline
5. \esp{Informática} & a) Word, Excel, Internet \\ \hline
6. \esp{Otros datos de interés} & d) permiso de conducir, aficiones, disponibilidad \\ \hline
\end{tabularx}
\end{center}

\textbf{Astuce mémo :} \esp{formación} = études et diplômes ; \esp{experiencia laboral} = emplois occupés ; \esp{otros datos de interés} = tout ce qui peut intéresser l'employeur (permis, loisirs, disponibilité immédiate).
\end{corrige}

\begin{corrige}[Exercice 8 — Correction d'erreurs typiques de francophones]
\begin{enumerate}
\item \esp{\emph{Más buena} $\rightarrow$ \textbf{mejor}} : \esp{Mi hermana es \textbf{mejor} que yo en matemáticas.} — « Meilleur » = comparatif irrégulier \esp{mejor}.
\item \esp{\emph{Tanto} $\rightarrow$ \textbf{tanta}} : \esp{Tengo \textbf{tanta} experiencia como mi colega.} — \esp{tanto} s'accorde avec le nom féminin \esp{experiencia}.
\item \esp{\emph{Es el más alto de la clase que todos} $\rightarrow$ \textbf{Es el más alto de la clase}.} — Le superlatif relatif se construit avec \esp{de} (\esp{de la clase}), pas avec \esp{que} (calque du français « que tous »).
\item \esp{\emph{Más peor} $\rightarrow$ \textbf{peor}} : \esp{Este puesto es \textbf{peor} que el anterior.} — \esp{peor} est déjà un comparatif : jamais \esp{más peor}.
\item \esp{\emph{Más que} $\rightarrow$ \textbf{más de}} : \esp{Gano \textbf{más de} cien mil francos al mes.} — Devant un chiffre : \esp{más de}.
\item \esp{\emph{Tan inteligente como su hermano es} $\rightarrow$ \textbf{tan inteligente como su hermano}} : \esp{Ella es \textbf{tan inteligente como su hermano}.} — Pas de verbe après \esp{como} (calque de « comme son frère l'est »).
\item \esp{\emph{Muy larguísima} $\rightarrow$ \textbf{larguísima}} : \esp{La reunión fue \textbf{larguísima}.} — On ne cumule jamais \esp{muy} et \esp{-ísimo}.
\item \esp{\emph{Más viejo} $\rightarrow$ \textbf{mayor}} : \esp{Mi padre es \textbf{mayor} que mi tío : tiene setenta años.} — Pour l'âge : \esp{mayor} ; \esp{viejo} signifie « vieux » et ne prend pas de comparatif analytique dans ce sens.
\end{enumerate}
\end{corrige}

\begin{corrige}[Exercice 9 — Texte, compréhension et langue (type Bac)]
\textbf{Barème indicatif :} compréhension 8 points (2 par question) ; questions de langue 6 points ; qualité de la langue 6 points. Total : 20.

\textbf{Questions de compréhension — réponses modèles :}
\begin{enumerate}
\item \esp{Es difícil porque los licenciados son hoy más numerosos que antes, pero las empresas ofrecen menos puestos estables.}
\item \esp{Muchos jóvenes aceptan trabajos peor pagados que los de sus padres ; el salario de un principiante es a menudo tan bajo como el de un obrero sin estudios.}
\item \esp{Tienen mejores oportunidades los candidatos que hablan tantos idiomas como piden las empresas internacionales, porque las lenguas son muy valoradas.}
\item \esp{El paro juvenil es más grave en las regiones del sur, donde es altísimo.}
\end{enumerate}

\textbf{Questions de langue :}
\begin{enumerate}
\item \textbf{Comparatifs relevés (trois au choix) :} \esp{más difícil} (supériorité) ; \esp{más numerosos que} (supériorité) ; \esp{menos puestos estables que} (infériorité) ; \esp{peor pagados que} (infériorité irrégulière) ; \esp{tan bajo como} (égalité adj.) ; \esp{tantos idiomas como} (égalité nom) ; \esp{más importante que nunca} (supériorité).
\item \textbf{Superlatifs absolus relevés :} \esp{dificilísima} (très difficile) ; \esp{altísimo} (très élevé).
\item \textbf{Traduction de la dernière phrase :} « Préparer un bon curriculum vitae est aujourd'hui plus important que jamais : c'est la première image que le candidat donne de lui-même. »
\end{enumerate}

\textbf{Points de vigilance :} répondre par des phrases complètes en reprenant les mots de la question ; ne pas confondre \esp{peor pagados} (« moins bien payés ») avec « plus mal payés » ; \esp{paro juvenil} = « chômage des jeunes » (faux ami : \esp{paro} $\neq$ « parade »).
\end{corrige}

\begin{corrige}[Exercice 10 — Thème et version (type Bac)]
\textbf{Barème indicatif :} thème 10 points (2 par phrase) ; version 10 points. Total : 20.

\textbf{Thème — corrigé :}
\begin{enumerate}
\item \esp{Mi hermano es \textbf{mayor} que yo.} — « Plus âgé » : \esp{mayor}, jamais \esp{más viejo}.
\item \esp{Gana \textbf{menos} dinero que su mujer.} (ou \esp{su esposa}) — Comparatif d'infériorité devant un nom.
\item \esp{Es \textbf{el mejor empleado} de la empresa.} — Superlatif relatif : article + \esp{mejor} + nom + \esp{de}.
\item \esp{Ella habla \textbf{tantos idiomas como} tú.} — \esp{tanto} s'accorde avec \esp{idiomas} (masculin pluriel).
\item \esp{\textbf{Más de} cincuenta candidatos han \textbf{solicitado} el puesto.} — Devant un chiffre : \esp{más de} ; « postuler » = \esp{solicitar (un puesto)}.
\end{enumerate}

\textbf{Version — traduction modèle :}

« À Yaoundé, la vie est plus chère que dans les petites villes du pays. Les loyers sont très élevés au centre, et les transports sont moins confortables qu'à Douala. Cependant, il y a plus d'offres d'emploi que dans les autres régions, et les jeunes arrivent chaque jour plus nombreux. Un employé de banque gagne autant que deux professeurs réunis. Pour les chefs d'entreprise, le candidat idéal est très ponctuel, sérieux et parle au moins deux langues. Chercher du travail est aujourd'hui plus difficile qu'avant, mais pas impossible : avec un bon CV et beaucoup de patience, tout est possible. »

\textbf{Points de vigilance :} \esp{altísimos} et \esp{puntualísimo} doivent être rendus par « très » ou « extrêmement » ; \esp{tanto como dos profesores juntos} = « autant que deux professeurs réunis » ; \esp{empresarios} = « chefs d'entreprise » (pas « entrepreneurs » au sens d'entreprises) ; \esp{por lo menos} = « au moins ».
\end{corrige}

\begin{corrige}[Exercice 11 — Rédaction guidée : mon CV et ma lettre de motivation (type Bac)]
\textbf{Barème indicatif (sur 20) :} respect des consignes et présentation du CV (6 rubriques obligatoires) : 6 points ; richesse du vocabulaire du travail : 4 points ; correction grammaticale (accords, accents, comparatifs et superlatif exigés) : 6 points ; cohérence et nombre de mots (120-150) : 4 points.

\textbf{Proposition de rédaction modèle :}

\begin{textebox}[Curriculum Vítae — Modèle élève]
\esp{\textbf{CURRÍCULUM VÍTAE}}

\esp{\textbf{Datos personales} : Ngo Bell Sandrine ; Calle de los Mangos, 45, Douala ; Tel. : 690 00 00 00 ; nacida el 12 de marzo de 2005 en Yaoundé ; nacionalidad camerunesa.}

\esp{\textbf{Formación académica} : 2023 : bachillerato, serie A4, Lycée de Douala ; 2020 : BEPC.}

\esp{\textbf{Experiencia laboral} : 2023-2024 : vendedora en un supermercado de Douala ; 2022 : prácticas de secretariado en una empresa de importación-exportación.}

\esp{\textbf{Idiomas} : francés : lengua materna ; español : nivel intermedio-alto ; inglés : nivel básico.}

\esp{\textbf{Informática} : Word, Excel, Internet.}

\esp{\textbf{Otros datos de interés} : permiso de conducir ; disponibilidad inmediata ; aficiones : fútbol y lectura.}
\end{textebox}

\begin{textebox}[Carta de motivación — Extrait]
\esp{Estimado \textbf{Señor Director} :}

\esp{He leído su anuncio y me interesa \textbf{muchísimo} el puesto de empleada de oficina. Creo que soy \textbf{la mejor candidata} porque hablo español \textbf{mejor que la mayoría} de los jóvenes de mi edad y tengo \textbf{más experiencia que} muchos candidatos. Mi nivel de informática es \textbf{tan bueno como} mi nivel de idiomas. Soy seria, \textbf{puntualísima} y muy motivada. Quedo a su disposición para una entrevista.}

\esp{Atentamente,}

\esp{Ngo Bell Sandrine}
\end{textebox}

\textbf{Points de vigilance :}
\begin{itemize}
\item Vérifier la présence des \textbf{6 rubriques} et de l'en-tête \esp{Currículum vítae} (accent sur le \emph{ú}).
\item La lettre doit contenir \textbf{au moins deux comparatifs} (ici : \esp{mejor que la mayoría}, \esp{más experiencia que}, \esp{tan bueno como}) et \textbf{un superlatif} (ici : \esp{puntualísima}, \esp{muchísimo}).
\item Accents obligatoires : \esp{Currículum}, \esp{informática}, \esp{académica}, \esp{motivación}, \esp{atención} (dans \esp{Atentamente}), \esp{prácticas}, \esp{importación-exportación}.
\item Formules de politesse : \esp{Estimado Señor Director\ldots} / \esp{Atentamente} ; ne pas traduire mot à mot « Cher Monsieur » par \esp{Querido señor} (trop familier).
\item Faux amis à éviter : « postuler » = \esp{solicitar}, non \esp{postular} (rare) ; « une expérience » professionnelle = \esp{experiencia laboral}.
\end{itemize}
\end{corrige}

\newpage

\coursec{Fiche bilan}

\begin{popfigure}
\begin{tikzpicture}[mindmap, grow cyclic, every node/.style={concept, text=white, font=\small, align=center}, root concept/.append style={concept color=popink, font=\bfseries}, level 1/.append style={level distance=3.4cm, sibling angle=60}, level 2/.append style={level distance=2.1cm, sibling angle=45, font=\scriptsize}]
\node[root concept, align=center]{E17\\Comparación\\y CV}
  child[concept color=popblue]{ node[align=center]{Comparativo\\de superioridad}
    child { node[concept color=popblue!70, align=center]{más \ldots\\que} }
    child { node[concept color=popblue!70, align=center]{mejor\\mayor} }
  }
  child[concept color=poporange]{ node[align=center]{Comparativo\\de inferioridad}
    child { node[concept color=poporange!70, align=center]{menos \ldots\\que} }
    child { node[concept color=poporange!70, align=center]{peor\\menor} }
  }
  child[concept color=popgreen]{ node[align=center]{Comparativo\\de igualdad}
    child { node[concept color=popgreen!70, align=center]{tan \ldots\\como} }
    child { node[concept color=popgreen!70, align=center]{tanto \ldots\\como} }
  }
  child[concept color=poppurple]{ node[align=center]{Superlativo\\relativo}
    child { node[concept color=poppurple!70, align=center]{el más \ldots\\de} }
  }
  child[concept color=poppink]{ node[align=center]{Superlativo\\absoluto}
    child { node[concept color=poppink!70] {-ísimo} }
    child { node[concept color=poppink!70, align=center]{muy,\\sumamente} }
  }
  child[concept color=popgold]{ node[align=center]{El currículum\\vítae}
    child { node[concept color=popgold!70, align=center]{datos,\\formación,\\experiencia} }
  };
\end{tikzpicture}
\legende{Carte mentale de la leçon E17 : comparer en espagnol et rédiger un CV.}
\end{popfigure}

\begin{bilan}
\textbf{1. Lexique clé (espagnol -- français)}
\begin{itemize}
  \item \esp{la experiencia laboral} : l'expérience professionnelle ; \esp{la formación académica} : la formation scolaire.
  \item \esp{los datos personales} : les informations personnelles ; \esp{el puesto de trabajo} : le poste de travail.
  \item \esp{las competencias / habilidades} : les compétences ; \esp{los idiomas} : les langues.
  \item \esp{solicitar un empleo} : postuler à un emploi ; \esp{la carta de presentación} : la lettre de motivation.
  \item \esp{el candidato / la candidata} : le/la candidat(e) ; \esp{el entrevistador} : le recruteur.
\end{itemize}

\textbf{2. Grammaire à connaître par cœur}
\begin{center}
\begin{tabularx}{\linewidth}{|l|X|X|}
\hline
\rowcolor{popblueL} \textbf{Structure} & \textbf{Espagnol} & \textbf{Français} \\
\hline
Supériorité & \esp{más alto que} & plus grand que \\
\hline
Infériorité & \esp{menos caro que} & moins cher que \\
\hline
Égalité (adjectif/adverbe) & \esp{tan rápido como} & aussi rapide que \\
\hline
Égalité (nom) & \esp{tanto dinero como} & autant d'argent que \\
\hline
Irréguliers & \esp{mejor, peor, mayor, menor} & meilleur, pire, plus âgé/grand, plus jeune/petit \\
\hline
Superlatif relatif & \esp{el más trabajador de la clase} & le plus travailleur de la classe \\
\hline
Superlatif absolu & \esp{guapísimo, muy bueno, sumamente útil} & très beau, extrêmement utile \\
\hline
\end{tabularx}
\end{center}

\textbf{3. Points de vigilance}
\begin{itemize}
  \item Jamais \esp{más mejor} : on dit \esp{mejor} seul. Jamais \esp{más bueno} en comparatif standard.
  \item Devant un nombre, « plus de » se traduit \esp{más de} : \esp{más de veinte alumnos} (plus de vingt élèves).
  \item \esp{mayor / menor} parlent surtout de l'âge ou de l'importance, pas de la taille physique (\esp{más grande}).
  \item Formation de l'absolu : \esp{grande} $\rightarrow$ \esp{grandísimo} ; \esp{fácil} $\rightarrow$ \esp{facilísimo} (retour du c latin).
  \item Le superlatif relatif exige l'article : \esp{la chica más lista} (la fille la plus intelligente).
\end{itemize}

\textbf{4. Méthode express : rédiger un CV en espagnol}
\begin{enumerate}
  \item Titre : \esp{Currículum Vítae} + \esp{Datos personales} (nombre, fecha de nacimiento, dirección, teléfono, correo).
  \item \esp{Formación académica} : diplômes du plus récent au plus ancien.
  \item \esp{Experiencia laboral} : postes, dates, entreprises.
  \item \esp{Idiomas} et \esp{Informática} : niveaux précis (\esp{nivel intermedio}, \esp{lengua materna}).
  \item \esp{Otros datos de interés} : loisirs, permis, centres d'intérêt.
  \item Style : phrases nominales, verbes à l'infinitif ou participe, une page maximum.
\end{enumerate}
\end{bilan}

\begin{aretenir}[Checklist avant le Bac]
\begin{itemize}
  \item $\square$ Je connais les trois comparatifs : \esp{más\ldots que}, \esp{menos\ldots que}, \esp{tan\ldots como}.
  \item $\square$ J'accorde \esp{tanto/tanta/tantos/tantas} avec le nom dans \esp{tanto\ldots como}.
  \item $\square$ Je traduis « meilleur » par \esp{mejor} et « pire » par \esp{peor}, sans \esp{más}.
  \item $\square$ Je distingue \esp{mayor/menor} (âge, importance) et \esp{más grande/más pequeño} (taille).
  \item $\square$ Je traduis « plus de + nombre » par \esp{más de} et non \esp{más que}.
  \item $\square$ Je forme le superlatif relatif avec article + \esp{más/menos} + adjectif + \esp{de}.
  \item $\square$ Je forme le superlatif absolu en \esp{-ísimo} avec les changements orthographiques (\esp{rico} $\rightarrow$ \esp{riquísimo}).
  \item $\square$ Je sais utiliser \esp{muy} et \esp{sumamente} pour renforcer un adjectif.
  \item $\square$ Je connais les rubriques d'un CV : \esp{datos personales, formación, experiencia, idiomas}.
  \item $\square$ Je sais décrire mon parcours au passé (\esp{estudié, trabajé, obtuve}).
  \item $\square$ Je vérifie les accents : \esp{currículum, vítae, académica, informática}.
  \item $\square$ Je relis ma traduction en comparant le nombre de comparatifs en français et en espagnol.
\end{itemize}
\end{aretenir}

\findecours

\end{document}
